% !TEX program = lualatex
% ---------------------------------------------------------------------------------------------
%  Readback note 47: Crystalline classes with few low weights, sparse large weights and tame types: a length bound and rank-free exponents
%
%  Build:  python3 tools/rb.py note <slug>
%
%  One unit of a Lean artifact, rewritten as ordinary mathematics. The prose is hand-written. Lean names,
%  line ranges and counts come from data.tex; quoted Lean statements are copied from the artifact by the build.
%
%  Rules of the text: standard terminology; every symbol is explained where it first appears and listed again
%  in the notation at the end; no letter has two meanings; a proof is complete or says what it takes for granted;
%  Lean names appear only in the LEAN lines and in the concordance.
%
%  The signature rule (the build enforces it): a function, relation or predicate that the artifact defines is
%  written only through its notation macro, which the build makes from the spec (`notation`) with one slot for
%  each explicit argument of the Lean definition. No argument is fixed and dropped, also not "for brevity" and
%  not for one section. Every argument and every hypothesis of a stated declaration stands inside its
%  statement, or in a Setting that the statement cites by number; only implicit and instance arguments (the
%  ambient objects) may stand in the conventions.
% ---------------------------------------------------------------------------------------------
\documentclass[11pt,a4paper]{article}
\ifdefined\pdfvariable\pdfvariable suppressoptionalinfo \numexpr 1+2+4+8\relax\fi
\usepackage[a4paper,hmargin=35mm,top=22mm,bottom=27mm,footskip=12mm,nohead]{geometry}

\usepackage{fontspec}
\directlua{
  luaotfload.add_fallback("gfallback", {"STIXTwoMath-Regular.otf:mode=node;", "DejaVuSans.ttf:mode=node;"})
  luaotfload.add_fallback("gmonofallback", {"DejaVuSansMono.ttf:mode=node;", "STIXTwoMath-Regular.otf:mode=node;"})
}
\setmainfont{NewCM10}[Extension=.otf,UprightFont=*-Book,ItalicFont=*-BookItalic,
  BoldFont=*-Bold,BoldItalicFont=*-BoldItalic,SmallCapsFont=*-Book,SmallCapsFeatures={Letters=SmallCaps},
  RawFeature={fallback=gfallback}]
\setsansfont{NewCMSans10}[Extension=.otf,UprightFont=*-Book,ItalicFont=*-BookOblique,
  BoldFont=*-Bold,BoldItalicFont=*-BoldOblique,RawFeature={fallback=gfallback}]
\setmonofont{JetBrainsMono}[Extension=.otf,UprightFont=*-Regular,ItalicFont=*-Italic,
  BoldFont=*-SemiBold,BoldItalicFont=*-SemiBoldItalic,Scale=MatchLowercase,
  RawFeature={-calt;-liga;fallback=gmonofallback}]
\usepackage{amsmath}
\usepackage{amsthm}
\usepackage[math-style=TeX]{unicode-math}
\setmathfont{NewCMMath-Book.otf}
\emergencystretch=2.5em\relax

\usepackage{xcolor}
\usepackage{tikz}
\usetikzlibrary{calc,arrows.meta,positioning,backgrounds}
\usepackage[protrusion=true,expansion=true]{microtype}
\usepackage{enumitem}
\usepackage{needspace}
\usepackage{booktabs}
\usepackage{array}
\usepackage{longtable}
\usepackage{fancyhdr}
\usepackage{lastpage}
\usepackage{multicol}
\usepackage[hidelinks,pdfusetitle]{hyperref}

\definecolor{ink}{HTML}{0F172A}
\definecolor{inksoft}{HTML}{475569}
\definecolor{inkfaint}{HTML}{94A3B8}
\definecolor{rulegrey}{HTML}{CBD5E1}
\definecolor{arith}{HTML}{C2410C}
\definecolor{est}{HTML}{7C3AED}
\definecolor{ana}{HTML}{0891B2}
\definecolor{ok}{HTML}{047857}
\definecolor{leancol}{HTML}{3730A3}

\title{Crystalline classes with few low weights, sparse large weights and tame types: a length bound and rank-free exponents}
\author{grh-compact}
\date{}

% ------------------------------------------------------------------------------------- macros
\newcommand{\leanfont}{\ttfamily\fontsize{8.6}{10.5}\selectfont\color{leancol}}
\newcommand{\lean}[1]{\csname LN@#1\endcsname}
\newcommand{\fillin}[1]{{\color{arith}\itshape[#1]}}  % a place of the template that is not written yet; the build reports every one that is left
\newcommand{\leanquote}[1]{\csname LQ@#1\endcsname}
% generated by rb.py note; do not edit
\newcommand{\ConcDecls}{19}
\newcommand{\Pin}{72fead10}
\newcommand{\PinText}{that of commit {\leanfont 72fead10}}
\newcommand{\ArtifactName}{grh-compact}
\expandafter\newcommand\csname LN@IsTameType\endcsname{\href{https://github.com/lsernau/grh-lean/blob/72fead102f663d5c7a407964a13dc0cf9d977207/grh-compact/src/New/Lift/ExponentTheorems.lean\#L39-L40}{{\leanfont Lift.\allowbreak{}IsTameType}}}
\expandafter\newcommand\csname LN@letterJumps\endcsname{\href{https://github.com/lsernau/grh-lean/blob/72fead102f663d5c7a407964a13dc0cf9d977207/grh-compact/src/New/Lift/ExponentTheorems.lean\#L44-L45}{{\leanfont Lift.\allowbreak{}KummerChart.\allowbreak{}letterJumps}}}
\expandafter\newcommand\csname LN@universalSide\endcsname{\href{https://github.com/lsernau/grh-lean/blob/72fead102f663d5c7a407964a13dc0cf9d977207/grh-compact/src/New/Lift/ExponentTheorems.lean\#L50-L52}{{\leanfont Lift.\allowbreak{}KummerChart.\allowbreak{}universalSide}}}
\expandafter\newcommand\csname LN@memcupOrthogonalofscaledCupBounds\endcsname{\href{https://github.com/lsernau/grh-lean/blob/72fead102f663d5c7a407964a13dc0cf9d977207/grh-compact/src/New/Lift/ExponentTheorems.lean\#L56-L59}{{\leanfont Lift.\allowbreak{}KummerChart.\allowbreak{}mem\_cupOrthogonal\_of\_scaledCupBounds}}}
\expandafter\newcommand\csname LN@fewlowjumps\endcsname{\href{https://github.com/lsernau/grh-lean/blob/72fead102f663d5c7a407964a13dc0cf9d977207/grh-compact/src/New/Lift/ExponentTheorems.lean\#L69-L99}{{\leanfont Lift.\allowbreak{}KummerChart.\allowbreak{}few\_low\_jumps}}}
\expandafter\newcommand\csname LQ@fewlowjumps\endcsname{\leanquoteline{0}{lemma KummerChart.few\_low\_jumps (hk : IsUnramified) (hp : 5 ≤ p) \{a η W\} \{b : \_ → ℕ\} \{ε : \_ →}\leanquoteline{2}{\_\} (ha : ∀ n, b n + W + η ≤ (p - 3) * a) (hη : 1 ≤ η) (hb : ∀ n, 6 * (a + η + W) < b n) (hs : ∀ n,}\leanquoteline{2}{NoMultipleNear p (6 * (a + η + W)) (b n)) (hL : InClass ⊤ ⊥ U' L) (hV : ∀ i, IsInJumpWindow η}\leanquoteline{2}{(-ε i) (b i) W (k.letterJumps ζ U i) ∧ IsInJumpWindow η (ε i) (b i) W (k.letterJumps ζ U' i)) :}\leanquoteline{2}{((∀ g u, u ∈ U g ∨ u ∈ U' g → u ∉ Ico (-↑a) 0) ∨ (∀ g, U' g ⊆ Icc (-η) 0) → crysCocycles L.ρ ≤}\leanquoteline{2}{cupOrthogonal (InClass ⊤ ⊥ U) L) ∧ ∀ m, (∀ g, ∑ u ∈ Icc (-↑η) ↑η, jump L.ρ g u ≤ m) →}\leanquoteline{2}{relLength (crysCocycles L.ρ) (cupOrthogonal (InClass ⊤ ⊥ U) L) ≤}\leanquoteline{2}{↑(11 * (k.f * (m * (a + 1) + finrank O L * (W + 1)))) * length O (O ⧸ Ideal.span \{(p : O)\})}}
\expandafter\newcommand\csname LN@tametypesexponent\endcsname{\href{https://github.com/lsernau/grh-lean/blob/72fead102f663d5c7a407964a13dc0cf9d977207/grh-compact/src/New/Lift/ExponentTheorems.lean\#L107-L139}{{\leanfont Lift.\allowbreak{}tame\_types\_exponent}}}
\expandafter\newcommand\csname LQ@tametypesexponent\endcsname{\leanquoteline{0}{lemma tame\_types\_exponent (hj : Surjective Γ.overGp.toGp) \{F K\} (hL : InClass F (K ·⁻¹) U' L) (hK :}\leanquoteline{2}{∀ s, IsWeilPair (toGp s) 0 → IsTameType p r N (K s)) \{η W b ε\} (hU : ∀ g,}\leanquoteline{2}{IsInJumpWindow η (-ε) b W (U g) ∧ IsInJumpWindow η ε b W (U' g)) (hp : b + 6 * (W + η) +}\leanquoteline{2}{r \textasciicircum{} 2 * (N + 1) \textasciicircum{} 2 + 2 < p) : ∀ a ∈ crysCocycles L.ρ,}\leanquoteline{2}{(p : O) \textasciicircum{} (3 * (r \textasciicircum{} 2 * (N + 1) \textasciicircum{} 2) + 150 * (η + W + 1)) • a ∈ cupOrthogonal (InClass F K U) L}}
\expandafter\newcommand\csname LN@deepConstants\endcsname{\href{https://github.com/lsernau/grh-lean/blob/72fead102f663d5c7a407964a13dc0cf9d977207/grh-compact/src/New/Lift/FewLowJumps.lean\#L20-L21}{{\leanfont Lift.\allowbreak{}KummerChart.\allowbreak{}PairingSide.\allowbreak{}deepConstants}}}
\expandafter\newcommand\csname LN@relLengthdeepConstantsle\endcsname{\href{https://github.com/lsernau/grh-lean/blob/72fead102f663d5c7a407964a13dc0cf9d977207/grh-compact/src/New/Lift/FewLowJumps.lean\#L35-L94}{{\leanfont Lift.\allowbreak{}KummerChart.\allowbreak{}PairingSide.\allowbreak{}relLength\_deepConstants\_le}}}
\expandafter\newcommand\csname LN@existsrelLengthdeepConstantsle\endcsname{\href{https://github.com/lsernau/grh-lean/blob/72fead102f663d5c7a407964a13dc0cf9d977207/grh-compact/src/New/Lift/FewLowJumps.lean\#L101-L143}{{\leanfont Lift.\allowbreak{}KummerChart.\allowbreak{}PairingSide.\allowbreak{}exists\_relLength\_deepConstants\_le}}}
\expandafter\newcommand\csname LN@existsscaledCupBoundsfewlowjumps\endcsname{\href{https://github.com/lsernau/grh-lean/blob/72fead102f663d5c7a407964a13dc0cf9d977207/grh-compact/src/New/Lift/FewLowJumps.lean\#L148-L163}{{\leanfont Lift.\allowbreak{}KummerChart.\allowbreak{}PairingSide.\allowbreak{}exists\_scaledCupBounds\_few\_low\_jumps}}}
\expandafter\newcommand\csname LN@TameChart\endcsname{\href{https://github.com/lsernau/grh-lean/blob/72fead102f663d5c7a407964a13dc0cf9d977207/grh-compact/src/New/Lift/SparseThickLetters.lean\#L27-L33}{{\leanfont Lift.\allowbreak{}TameChart}}}
\expandafter\newcommand\csname LN@existszsmuleqsumtameGen\endcsname{\href{https://github.com/lsernau/grh-lean/blob/72fead102f663d5c7a407964a13dc0cf9d977207/grh-compact/src/New/Lift/SparseThickLetters.lean\#L43-L61}{{\leanfont Lift.\allowbreak{}TameChart.\allowbreak{}exists\_z\_smul\_eq\_sum\_tameGen}}}
\expandafter\newcommand\csname LN@aevallinearFrobextLatticeeqzero\endcsname{\href{https://github.com/lsernau/grh-lean/blob/72fead102f663d5c7a407964a13dc0cf9d977207/grh-compact/src/New/Lift/SparseThickLetters.lean\#L67-L76}{{\leanfont Lift.\allowbreak{}TameChart.\allowbreak{}aeval\_linearFrob\_extLattice\_eq\_zero}}}
\expandafter\newcommand\csname LN@killsCokernelofisKilledBy\endcsname{\href{https://github.com/lsernau/grh-lean/blob/72fead102f663d5c7a407964a13dc0cf9d977207/grh-compact/src/New/Lift/SparseThickLetters.lean\#L81-L116}{{\leanfont Lift.\allowbreak{}TameChart.\allowbreak{}killsCokernel\_of\_isKilledBy}}}
\expandafter\newcommand\csname LN@existstameChart\endcsname{\href{https://github.com/lsernau/grh-lean/blob/72fead102f663d5c7a407964a13dc0cf9d977207/grh-compact/src/New/Lift/SparseThickLetters.lean\#L124-L159}{{\leanfont Lift.\allowbreak{}exists\_tameChart}}}
\expandafter\newcommand\csname LN@killsCokernelatcut\endcsname{\href{https://github.com/lsernau/grh-lean/blob/72fead102f663d5c7a407964a13dc0cf9d977207/grh-compact/src/New/Lift/SparseThickLetters.lean\#L178-L194}{{\leanfont Lift.\allowbreak{}KummerChart.\allowbreak{}PairingSide.\allowbreak{}killsCokernel\_at\_cut}}}
\expandafter\newcommand\csname LN@scaledCupBoundsofcutwindows\endcsname{\href{https://github.com/lsernau/grh-lean/blob/72fead102f663d5c7a407964a13dc0cf9d977207/grh-compact/src/New/Lift/SparseThickLetters.lean\#L199-L243}{{\leanfont Lift.\allowbreak{}KummerChart.\allowbreak{}PairingSide.\allowbreak{}SidePairing.\allowbreak{}scaledCupBounds\_of\_cut\_windows}}}
\expandafter\newcommand\csname LN@scaledCupBoundssparsethick\endcsname{\href{https://github.com/lsernau/grh-lean/blob/72fead102f663d5c7a407964a13dc0cf9d977207/grh-compact/src/New/Lift/SparseThickLetters.lean\#L253-L270}{{\leanfont Lift.\allowbreak{}KummerChart.\allowbreak{}PairingSide.\allowbreak{}SidePairing.\allowbreak{}scaledCupBounds\_sparse\_thick}}}
\expandafter\newcommand\csname LN@scaledCupBoundstwosizes\endcsname{\href{https://github.com/lsernau/grh-lean/blob/72fead102f663d5c7a407964a13dc0cf9d977207/grh-compact/src/New/Lift/SparseThickLetters.lean\#L273-L286}{{\leanfont Lift.\allowbreak{}KummerChart.\allowbreak{}PairingSide.\allowbreak{}SidePairing.\allowbreak{}scaledCupBounds\_two\_sizes}}}
\newcommand{\ConcordanceRows}{\crow{IsTameType}{\href{https://github.com/lsernau/grh-lean/blob/72fead102f663d5c7a407964a13dc0cf9d977207/grh-compact/src/New/Lift/ExponentTheorems.lean\#L39-L40}{{\leanfont Lift.\allowbreak{}IsTameType}}}{ExponentTheorems.lean}{39--40}
\crow{letterJumps}{\href{https://github.com/lsernau/grh-lean/blob/72fead102f663d5c7a407964a13dc0cf9d977207/grh-compact/src/New/Lift/ExponentTheorems.lean\#L44-L45}{{\leanfont Lift.\allowbreak{}KummerChart.\allowbreak{}letterJumps}}}{ExponentTheorems.lean}{44--45}
\crow{universalSide}{\href{https://github.com/lsernau/grh-lean/blob/72fead102f663d5c7a407964a13dc0cf9d977207/grh-compact/src/New/Lift/ExponentTheorems.lean\#L50-L52}{{\leanfont Lift.\allowbreak{}KummerChart.\allowbreak{}universalSide}}}{ExponentTheorems.lean}{50--52}
\crow{memcupOrthogonalofscaledCupBounds}{\href{https://github.com/lsernau/grh-lean/blob/72fead102f663d5c7a407964a13dc0cf9d977207/grh-compact/src/New/Lift/ExponentTheorems.lean\#L56-L59}{{\leanfont Lift.\allowbreak{}KummerChart.\allowbreak{}mem\_cupOrthogonal\_of\_scaledCupBounds}}}{ExponentTheorems.lean}{56--59}
\crow{fewlowjumps}{\href{https://github.com/lsernau/grh-lean/blob/72fead102f663d5c7a407964a13dc0cf9d977207/grh-compact/src/New/Lift/ExponentTheorems.lean\#L69-L99}{{\leanfont Lift.\allowbreak{}KummerChart.\allowbreak{}few\_low\_jumps}}}{ExponentTheorems.lean}{69--99}
\crow{tametypesexponent}{\href{https://github.com/lsernau/grh-lean/blob/72fead102f663d5c7a407964a13dc0cf9d977207/grh-compact/src/New/Lift/ExponentTheorems.lean\#L107-L139}{{\leanfont Lift.\allowbreak{}tame\_types\_exponent}}}{ExponentTheorems.lean}{107--139}
\crow{deepConstants}{\href{https://github.com/lsernau/grh-lean/blob/72fead102f663d5c7a407964a13dc0cf9d977207/grh-compact/src/New/Lift/FewLowJumps.lean\#L20-L21}{{\leanfont Lift.\allowbreak{}KummerChart.\allowbreak{}PairingSide.\allowbreak{}deepConstants}}}{FewLowJumps.lean}{20--21}
\crow{relLengthdeepConstantsle}{\href{https://github.com/lsernau/grh-lean/blob/72fead102f663d5c7a407964a13dc0cf9d977207/grh-compact/src/New/Lift/FewLowJumps.lean\#L35-L94}{{\leanfont Lift.\allowbreak{}KummerChart.\allowbreak{}PairingSide.\allowbreak{}relLength\_deepConstants\_le}}}{FewLowJumps.lean}{35--94}
\crow{existsrelLengthdeepConstantsle}{\href{https://github.com/lsernau/grh-lean/blob/72fead102f663d5c7a407964a13dc0cf9d977207/grh-compact/src/New/Lift/FewLowJumps.lean\#L101-L143}{{\leanfont Lift.\allowbreak{}KummerChart.\allowbreak{}PairingSide.\allowbreak{}exists\_relLength\_deepConstants\_le}}}{FewLowJumps.lean}{101--143}
\crow{existsscaledCupBoundsfewlowjumps}{\href{https://github.com/lsernau/grh-lean/blob/72fead102f663d5c7a407964a13dc0cf9d977207/grh-compact/src/New/Lift/FewLowJumps.lean\#L148-L163}{{\leanfont Lift.\allowbreak{}KummerChart.\allowbreak{}PairingSide.\allowbreak{}exists\_scaledCupBounds\_few\_low\_jumps}}}{FewLowJumps.lean}{148--163}
\crow{TameChart}{\href{https://github.com/lsernau/grh-lean/blob/72fead102f663d5c7a407964a13dc0cf9d977207/grh-compact/src/New/Lift/SparseThickLetters.lean\#L27-L33}{{\leanfont Lift.\allowbreak{}TameChart}}}{SparseThickLetters.lean}{27--33}
\crow{existszsmuleqsumtameGen}{\href{https://github.com/lsernau/grh-lean/blob/72fead102f663d5c7a407964a13dc0cf9d977207/grh-compact/src/New/Lift/SparseThickLetters.lean\#L43-L61}{{\leanfont Lift.\allowbreak{}TameChart.\allowbreak{}exists\_z\_smul\_eq\_sum\_tameGen}}}{SparseThickLetters.lean}{43--61}
\crow{aevallinearFrobextLatticeeqzero}{\href{https://github.com/lsernau/grh-lean/blob/72fead102f663d5c7a407964a13dc0cf9d977207/grh-compact/src/New/Lift/SparseThickLetters.lean\#L67-L76}{{\leanfont Lift.\allowbreak{}TameChart.\allowbreak{}aeval\_linearFrob\_extLattice\_eq\_zero}}}{SparseThickLetters.lean}{67--76}
\crow{killsCokernelofisKilledBy}{\href{https://github.com/lsernau/grh-lean/blob/72fead102f663d5c7a407964a13dc0cf9d977207/grh-compact/src/New/Lift/SparseThickLetters.lean\#L81-L116}{{\leanfont Lift.\allowbreak{}TameChart.\allowbreak{}killsCokernel\_of\_isKilledBy}}}{SparseThickLetters.lean}{81--116}
\crow{existstameChart}{\href{https://github.com/lsernau/grh-lean/blob/72fead102f663d5c7a407964a13dc0cf9d977207/grh-compact/src/New/Lift/SparseThickLetters.lean\#L124-L159}{{\leanfont Lift.\allowbreak{}exists\_tameChart}}}{SparseThickLetters.lean}{124--159}
\crow{killsCokernelatcut}{\href{https://github.com/lsernau/grh-lean/blob/72fead102f663d5c7a407964a13dc0cf9d977207/grh-compact/src/New/Lift/SparseThickLetters.lean\#L178-L194}{{\leanfont Lift.\allowbreak{}KummerChart.\allowbreak{}PairingSide.\allowbreak{}killsCokernel\_at\_cut}}}{SparseThickLetters.lean}{178--194}
\crow{scaledCupBoundsofcutwindows}{\href{https://github.com/lsernau/grh-lean/blob/72fead102f663d5c7a407964a13dc0cf9d977207/grh-compact/src/New/Lift/SparseThickLetters.lean\#L199-L243}{{\leanfont Lift.\allowbreak{}KummerChart.\allowbreak{}PairingSide.\allowbreak{}SidePairing.\allowbreak{}scaledCupBounds\_of\_cut\_windows}}}{SparseThickLetters.lean}{199--243}
\crow{scaledCupBoundssparsethick}{\href{https://github.com/lsernau/grh-lean/blob/72fead102f663d5c7a407964a13dc0cf9d977207/grh-compact/src/New/Lift/SparseThickLetters.lean\#L253-L270}{{\leanfont Lift.\allowbreak{}KummerChart.\allowbreak{}PairingSide.\allowbreak{}SidePairing.\allowbreak{}scaledCupBounds\_sparse\_thick}}}{SparseThickLetters.lean}{253--270}
\crow{scaledCupBoundstwosizes}{\href{https://github.com/lsernau/grh-lean/blob/72fead102f663d5c7a407964a13dc0cf9d977207/grh-compact/src/New/Lift/SparseThickLetters.lean\#L273-L286}{{\leanfont Lift.\allowbreak{}KummerChart.\allowbreak{}PairingSide.\allowbreak{}SidePairing.\allowbreak{}scaledCupBounds\_two\_sizes}}}{SparseThickLetters.lean}{273--286}}
% the notations of the defined functions: one slot for each explicit argument of the Lean declaration
\newcommand{\TameType}[4]{\operatorname{Tame}_{#1}(#2,#3;#4)}
\newcommand{\letterJ}[4]{\mathcal{U}_{#1,#2}(#3)(#4)}
\newcommand{\univS}[3]{\mathsf{S}^{\mathrm{u}}_{#1,#2}(#3)}
\newcommand{\deepC}[4]{\operatorname{DC}_{#1}(#2;#3,#4)}
\newcommand{\UnitMembers}{141}
\newcommand{\UnitInNote}{19}
\newcommand{\UnitProofs}{69}
\newcommand{\UnitSteps}{375}
\newcommand{\UnitAllSteps}{430}


\linespread{1.08}                       % 11 pt on 14.7 pt
\setlength{\parindent}{0pt}
\setlength{\parskip}{5.5pt plus 1.5pt minus 0.5pt}
\setlength{\abovedisplayskip}{7.5pt plus 2pt minus 3pt}
\setlength{\belowdisplayskip}{7.5pt plus 2pt minus 3pt}
\setlength{\abovedisplayshortskip}{3pt plus 2pt}
\setlength{\belowdisplayshortskip}{6pt plus 2pt minus 2pt}
\emergencystretch=1.5em
\clubpenalty=10000 \widowpenalty=10000 \displaywidowpenalty=10000
\relpenalty=7000 \binoppenalty=9000
\raggedbottom
\pagestyle{fancy}\fancyhf{}\renewcommand{\headrulewidth}{0pt}
\fancyfoot[C]{\sffamily\fontsize{7.6}{9}\selectfont\color{inkfaint}\thepage\,/\,\pageref*{LastPage}}
\setlist{topsep=2pt,partopsep=0pt,parsep=0pt,itemsep=4pt,leftmargin=2.1em,labelsep=0.55em}

\newcommand{\Q}{\mathbb{Q}}
\newcommand{\Z}{\mathbb{Z}}
\newcommand{\F}{\mathbb{F}}
\newcommand{\len}{\operatorname{len}}
\newcommand{\cok}{\operatorname{cok}}
\newcommand{\End}{\operatorname{End}}
\newcommand{\Hom}{\operatorname{Hom}}
\newcommand{\Gal}{\operatorname{Gal}}
\newcommand{\diag}{\operatorname{diag}}
\newcommand{\kw}[1]{\emph{#1}}
\newcommand{\m}[1]{\mbox{$#1$}}        % a short formula that must not be broken across lines
\newcommand{\head}[1]{\par\addvspace{17pt plus 4pt minus 2pt}\Needspace{6\baselineskip}%
  \noindent{\fontsize{13.2}{16}\selectfont\bfseries\color{ink}#1}\par\nopagebreak\vspace{-1pt}}
\newcommand{\sect}[2]{\refstepcounter{section}\label{#2}\par\addvspace{19pt plus 4pt minus 2pt}\Needspace{7\baselineskip}%
  \noindent\llap{{\fontsize{13.2}{16}\selectfont\bfseries\color{inkfaint}\thesection}\hspace{9pt}}%
  {\fontsize{13.2}{16}\selectfont\bfseries\color{ink}#1}\par\nopagebreak\vspace{-1pt}}
\newcommand{\subhead}[1]{\par\addvspace{9pt plus 2pt}\Needspace{4\baselineskip}\noindent\textbf{#1}\enspace\ignorespaces}
% where a statement lives in the Lean archive
\newcommand{\leantag}{{\sffamily\bfseries\fontsize{6.3}{8}\selectfont\color{leancol}\addfontfeatures{LetterSpace=7}LEAN}}
\newcommand{\sep}{\unskip\nobreak\hspace{5pt}{\color{inkfaint}\textperiodcentered}\hspace{5pt plus 2pt}\ignorespaces}
\newcommand{\leanline}[1]{\par\nopagebreak\vspace{-2.5pt}{\fontsize{8.9}{12.4}\selectfont\color{inksoft}\raggedright
  \noindent\llap{\leantag\hspace{9pt}}#1\par}\vspace{1pt}}
% a displayed statement in a tinted box: #1 colour, #2 text
\newcommand{\statement}[2]{\par\addvspace{7pt}\begingroup\setlength{\fboxsep}{8.5pt}\setlength{\fboxrule}{0.6pt}%
  \noindent\fcolorbox{#1!30!white}{#1!3!white}{\parbox{\dimexpr\linewidth-2\fboxsep-2\fboxrule\relax}{#2}}%
  \endgroup\par\addvspace{7pt}}
\newcommand{\nt}[2]{\par\noindent\hangindent=1.1em\hangafter=1 #1\enspace{\color{inksoft}#2}\par\vspace{1pt}}
\newcommand{\legdot}[1]{\tikz[baseline=-0.62ex]\fill[#1] (0,0) circle (2.5pt);\hspace{3pt}}
\newcommand{\crow}[4]{\csname CS@#1\endcsname & #2 & \textcolor{inksoft}{#3} & \textcolor{inksoft}{#4}\\}
\newcommand{\leanquoteline}[2]{\par\noindent\hangindent=\dimexpr#1\fontdimen2\font+2.2em\relax\hangafter=1
  \hspace*{#1\fontdimen2\font}#2}
\newcommand{\quoted}[1]{\par\vspace{3pt}{\ttfamily\fontsize{7.5}{10.4}\selectfont\color{leancol}\raggedright
  \setlength{\parskip}{0pt}\leanquote{#1}\par}\vspace{4pt}}
\newcommand{\bx}[3]{{\sffamily\bfseries\fontsize{7.3}{9.2}\selectfont\color{#1!82!black}#2}\\[1.5pt]#3}
\newcommand{\cs}[2]{\expandafter\newcommand\csname CS@#1\endcsname{#2}}

% theorems
\newtheoremstyle{plainx}{7pt plus 2pt}{7pt plus 2pt}{\itshape}{}{\bfseries}{.}{0.5em}{}
\newtheoremstyle{uprightx}{7pt plus 2pt}{7pt plus 2pt}{}{}{\bfseries}{.}{0.5em}{}
\theoremstyle{plainx}
\newtheorem{theorem}{Theorem}[section]
\newtheorem{lemma}[theorem]{Lemma}
\newtheorem{proposition}[theorem]{Proposition}
\newtheorem{corollary}[theorem]{Corollary}
\theoremstyle{uprightx}
\newtheorem{definition}[theorem]{Definition}
\newtheorem{setting}[theorem]{Setting}          % hypotheses shared by several statements; each of them cites it by number
\newtheorem{remark}[theorem]{Remark}
\newtheorem{example}[theorem]{Example}
\renewcommand{\qedsymbol}{{\setlength{\fboxsep}{0pt}\setlength{\fboxrule}{0.5pt}\fbox{\rule{0pt}{0.46em}\hspace{0.46em}}}}
% what holds throughout the note: the ambient objects only (implicit and instance arguments), never a hypothesis
\newenvironment{conventions}{\par\addvspace{5pt}\noindent\textbf{Conventions.}\enspace\ignorespaces}{\par\addvspace{3pt}}
\makeatletter
\renewenvironment{proof}[1][Proof]{\par\pushQED{\qed}\normalfont\topsep0pt\partopsep0pt
  \trivlist\item[\hskip\labelsep\itshape #1.]\ignorespaces}{\popQED\endtrivlist\@endpefalse}
\makeatother

% where each declaration of the concordance is treated: \cs{key}{place}
\cs{fewlowjumps}{Theorem A}
\cs{tametypesexponent}{Theorem B}
\cs{scaledCupBoundssparsethick}{Theorem C}
\cs{scaledCupBoundstwosizes}{Theorem D}
\cs{deepConstants}{Definition~\ref{d:dc}}
\cs{TameChart}{Definition~\ref{d:tame}}
\cs{IsTameType}{Definition~\ref{d:univ}}
\cs{letterJumps}{Definition~\ref{d:univ}}
\cs{universalSide}{Definition~\ref{d:univ}}
\cs{relLengthdeepConstantsle}{Lemma~\ref{l:rel}}
\cs{existsrelLengthdeepConstantsle}{Lemma~\ref{l:col}}
\cs{existsscaledCupBoundsfewlowjumps}{Proposition~\ref{p:few}}
\cs{existszsmuleqsumtameGen}{Lemma~\ref{l:zsum}}
\cs{aevallinearFrobextLatticeeqzero}{Lemma~\ref{l:extkill}}
\cs{killsCokernelofisKilledBy}{Proposition~\ref{p:tamekill}}
\cs{existstameChart}{Proposition~\ref{p:tameexist}}
\cs{killsCokernelatcut}{Lemma~\ref{l:atcut}}
\cs{scaledCupBoundsofcutwindows}{Proposition~\ref{p:cuts}}
\cs{memcupOrthogonalofscaledCupBounds}{Lemma~\ref{l:orth}}
% ---------- common symbols of the notes on unit 136
\providecommand{\Zp}{\mathbb{Z}_p}\providecommand{\Qp}{\mathbb{Q}_p}\providecommand{\GQp}{G_{\mathbb{Q}_p}}
\providecommand{\Ainf}{A_{\inf}}\providecommand{\WC}{W(C^{\flat})}\providecommand{\BdR}{B_{\mathrm{dR}}}
\providecommand{\Fil}{\operatorname{Fil}}\providecommand{\fS}{\mathfrak{S}}\providecommand{\cE}{\mathcal{E}}
\providecommand{\cO}{\mathcal{O}}\providecommand{\fc}{\mathfrak{c}}\providecommand{\fd}{\mathfrak{d}}
\providecommand{\WW}{\mathbb{W}}\providecommand{\Per}{\operatorname{Per}}\providecommand{\pr}{\operatorname{pr}}
\providecommand{\Minf}{\mathfrak{M}^{\inf}}\providecommand{\Mloc}{\mathfrak{M}_{\mathrm{loc}}}
\providecommand{\Mlocinf}{\mathfrak{M}^{\inf}_{\mathrm{loc}}}\providecommand{\std}{\operatorname{std}}
\providecommand{\ret}{\operatorname{ret}}
\providecommand{\inputbox}[1]{\statement{inksoft}{#1}}
% ---------- notations of definitions that Part 1 owns (same forms as there)
\providecommand{\cost}[2]{\operatorname{cost}(#1,#2)}
\providecommand{\wExp}[3]{\mathfrak{X}_{#1}(#2;#3)}
\providecommand{\locLP}[2]{\mathfrak{N}_{#1,#2}}
\providecommand{\locKP}[1]{(#1)_{\fS}}
\providecommand{\invMod}[1]{\operatorname{Inv}(#1)}
\providecommand{\impE}[4]{\operatorname{imp}(#1,#2)(#3,#4)}
\providecommand{\wPair}[3]{{#1}_{W}(#2,#3)}
\providecommand{\frobX}[5]{\operatorname{fr}_{#1,#2}(#3)(#4,#5)}
\providecommand{\extL}[1]{\tilde{Y}_{#1}}
\providecommand{\trivL}[1]{C^{\mathrm{tr}}_{#1}}
\providecommand{\projC}[1]{\pr^{C}_{#1}}
\providecommand{\JumpWin}[5]{#5\subseteq[-#1,#1]\cup#2\cdot[#3,#3+#4]}
\providecommand{\raisedE}[3]{\operatorname{rs}_{#1}(#2,#3)}
\providecommand{\twoWin}[2]{\mathcal{W}_{#1}(#2)}
\providecommand{\NoMult}[3]{\operatorname{NM}(#1;#2,#3)}
\providecommand{\constAt}[3]{\operatorname{cst}_{#1,#2}(#3)}
\providecommand{\liftId}[4]{\mathfrak{a}_{#1}(#2;#3,#4)}
\providecommand{\Kills}[4]{\operatorname{Kill}_{#1}(#2;#3,#4)}
% the cost of a map (the spec says why it has no generated macro): side, cost, map
\providecommand{\Reaches}[3]{\operatorname{Cost}_{#1}(#3)\le #2}

% ---------- notations of definitions that Part 2 owns (same forms as there)
\providecommand{\projY}[1]{\pr^{Y}_{#1}}
\providecommand{\inclY}[1]{\operatorname{in}^{Y}_{#1}}
\providecommand{\projCloc}[1]{\pr^{C,\mathrm{loc}}_{#1}}
\providecommand{\tensL}[2]{\mathsf{T}_{#1,#2}}
\providecommand{\tateL}[1]{#1(1)}
\providecommand{\third}[1]{S^{\vee}_{#1}}
\providecommand{\pairW}[1]{\mathsf{b}_{#1}}
\providecommand{\Safe}[4]{\operatorname{Safe}_{#1}(#2;#3,#4)}
\providecommand{\CupB}[4]{\operatorname{Cup}_{#1}(#2;#3,#4)}
\providecommand{\kerX}[3]{\mathfrak{X}^{0}_{#1}(#2;#3)}
\providecommand{\deepE}[3]{\operatorname{Deep}_{#1}(#2;#3)}
\providecommand{\ptsOver}[6]{\operatorname{Pt}_{#1}(#2,#3;#4;#5,#6)}
\providecommand{\raiseX}[5]{\operatorname{up}_{#2,#3}(#1)(#4,#5)}
\providecommand{\shStd}[4]{\operatorname{sh}_{#1}(#2)(#3,#4)}
\providecommand{\constE}[3]{\mathfrak{K}_{#1,#2,#3}}
% ---------- notions of the background library used in Part 3
\providecommand{\InCl}[3]{\operatorname{Cl}(#1,#2;#3)}
\providecommand{\cupOrth}[2]{{#2}^{\perp#1}}
\providecommand{\relLen}[3]{\operatorname{len}_{#1}\bigl(#2/(#2\cap#3)\bigr)}
\providecommand{\linF}[3]{\beta^{#1}_{#2,#3}}

\begin{document}
\color{ink}

{\raggedright
{\sffamily\bfseries\fontsize{7.2}{9}\selectfont\color{leancol}\addfontfeatures{LetterSpace=8}READBACK NOTE · CHAPTER VIII · UNIT 136 · PART 3 OF 3\par}
\vspace{5pt}
{\fontsize{21.5}{25.5}\selectfont\bfseries Crystalline classes with few low weights, sparse large weights and tame types: a length bound and rank-free exponents\par}
\vspace{5pt}
{\fontsize{12.4}{16}\selectfont\itshape\color{inksoft} One unit of the Lean artifact \ArtifactName, read back into
mathematics\par}}
\vspace{8pt}
{\color{rulegrey}\hrule height 0.6pt}
\vspace{7pt}
\thispagestyle{empty}

\textbf{Summary.} This note draws the consequences of Parts 1 and~2 that the bound for Selmer groups at~$\ell$
uses. For a crystalline lattice~$L$ over an unramified base whose weights lie in a low window and a far window,
the crystalline classes of~$L$ whose cup product with every crystalline class of every dual-type lattice
bounds have, inside all crystalline classes, a colength that is linear in the number of low weights and in the
rank, with an explicit constant; and all crystalline classes have this property when no weight is small and
negative. For lattices of tame inertial type over~$\Qp$, a power of~$p$ whose exponent does not depend on the
ramification index makes every crystalline class orthogonal in this sense: the cokernel of Kisin modules is
killed by a power of~$p$ read off a polynomial that kills the tame generator. Two further rank-free exponents
are proved, for weights of two sizes and for sparse large weights, from a general statement with cuts on the
Frobenius orbits.

\statement{leancol}{\fontsize{9.6}{13.2}\selectfont\textbf{What this note is.} The artifact \ArtifactName{} is a
body of Lean code: the compact edition of a formal proof, due to L.~Alpöge with the AI system Claude, of the Riemann
hypothesis and of its generalisation to Dirichlet $L$-functions. That headline is a claim. The authors report that
every file compiles with Lean v4.35.0 against a pinned version of Mathlib and that Lean's kernel accepted the
proofs; an independent build is recorded in the handout's record of checks when it reports. No mathematician has
read the argument through, and nothing here says that the Riemann hypothesis is proved. This note is the third of
three on one unit of the artifact, the folder on lifting; it treats \UnitInNote{} of the unit's \UnitMembers{}
declarations, those of the last three files. Theorems A and~B, Proposition~\ref{p:few} and the results on tame
charts are proved in full from Parts 1 and~2 and from Inputs of the background library, which are stated and not
proved. Proposition~\ref{p:cuts} is proved in full as well, including the final comparisons of exponents, which
the formal text settles by a decision procedure for linear arithmetic and which are carried out by hand here;
Theorems C and~D are deduced from it. No proof of this note is partial. Two Inputs of the background library, the
crystalline criterion~\ref{k:crit} and the finiteness of the module of crystalline cocycles~\ref{k:bk}, are
statements of the formal library that are close to, but not the same as, published theorems; the boxes say
exactly how. Every statement was compared with its Lean text, argument by argument. The Lean text read is \PinText; Appendix~A is the concordance of \ConcDecls{} declarations.}

\sect{The statements}{s:intro}
The objects are those of Setting~\ref{set:base}: the Kummer chart with its letters, lattices over a coefficient
ring~$\cO$, classes $\InCl{\mathcal{F}}{\mathsf{K}}{U}$ of lattices, the module $\cupOrth{\mathcal{Q}}{L}$ of the
maps orthogonal to a class, and the notation of Parts 1 and~2.

\statement{arith}{\phantomsection\label{t:A}\textbf{Theorem A} (few low weights).\enspace
\emph{In Setting~\ref{set:base}, let $\Gamma$ carry a topology for which $j_\Gamma$ is an open embedding. Let
$\fc$ be a Kummer chart for~$\Gamma$ that is unramified, with residue degree~$h$, let \m{\omega\in\cO} with
\m{\omega^{p^h-1}=1}, let $L$ be a lattice over~$\cO$, and let $U=(U(\gamma))_\gamma$ and
$U'=(U'(\gamma))_\gamma$ be families of finite sets of integers indexed by the $\Zp$-algebra maps
$\gamma\colon\cO\to C$. Suppose \m{p\ge5}. Let \m{\rho,\eta,W\ge0} be integers, and for every letter~$i$ let
\m{\beta(i)\ge0} be an integer and \m{\mathsf{s}(i)} a sign, such that for every letter~$i$
\[\begin{gathered}\beta(i)+W+\eta\le(p-3)\rho,\qquad\eta\ge1,\\
6(\rho+\eta+W)<\beta(i),\qquad\NoMult{p}{6(\rho+\eta+W)}{\beta(i)},\end{gathered}\]
$\JumpWin{\eta}{-\mathsf{s}(i)}{\beta(i)}{W}{\letterJ{\fc}{\omega}{U}{i}}$ and
$\JumpWin{\eta}{\mathsf{s}(i)}{\beta(i)}{W}{\letterJ{\fc}{\omega}{U'}{i}}$. Suppose that
$L\in\InCl{\top}{\bot}{U'}$. Then:
\begin{enumerate}[label=\textup{(\arabic*)},leftmargin=2.3em,itemsep=2pt]
\item if no element of $U(\gamma)\cup U'(\gamma)$ lies in \m{[-\rho,0)}, for any~$\gamma$, or if
$U'(\gamma)\subseteq[-\eta,0]$ for every~$\gamma$, then
$Z^{1}_{\mathrm{f}}(L)\subseteq\cupOrth{\InCl{\top}{\bot}{U}}{L}$;
\item for every integer \m{m_0\ge0} such that the multiplicities of the jumps \m{u\in[-\eta,\eta]} of~$L$
at~$\gamma$ have sum at most~$m_0$, for every~$\gamma$,
\[\begin{gathered}\relLen{\cO}{Z^{1}_{\mathrm{f}}(L)}{\cupOrth{\InCl{\top}{\bot}{U}}{L}}\\
\le\ 11\,h\,\bigl(m_0(\rho+1)+\operatorname{rk}_{\cO}(L)\,(W+1)\bigr)\cdot\operatorname{len}_{\cO}(\cO/p\cO).\end{gathered}\]
\end{enumerate}}}
\leanline{\lean{fewlowjumps}}

\statement{arith}{\phantomsection\label{t:B}\textbf{Theorem B} (tame types: an exponent independent of the
ramification index).\enspace
\emph{In Setting~\ref{set:base}, let $\Gamma$ carry a topology for which $j_\Gamma$ is an open embedding, and
suppose that $j_\Gamma$ is surjective. Let $L$ be a lattice over~$\cO$, $U$ and~$U'$ families of finite sets of
integers indexed by the maps $\gamma\colon\cO\to C$, $\mathcal{F}$ a filter on~$\Gamma$, $\mathsf{K}$ a relation
between elements of~$\Gamma$ and polynomials over~$\cO$, and $\mathsf{r},\mathsf{n}\ge0$ integers. Write
$\mathsf{K}^{-}$ for the relation $(s,Q)\mapsto\mathsf{K}(s^{-1},Q)$. Suppose that
$L\in\InCl{\mathcal{F}}{\mathsf{K}^{-}}{U'}$, and that $\TameType{p}{\mathsf{r}}{\mathsf{n}}{\mathsf{K}(s,\cdot)}$
for every \m{s\in\Gamma} with \m{j_\Gamma(s)} in the inertia group. Let \m{\eta,W,\beta\ge0} be integers and
$\mathsf{s}$ a sign with $\JumpWin{\eta}{-\mathsf{s}}{\beta}{W}{U(\gamma)}$ and
$\JumpWin{\eta}{\mathsf{s}}{\beta}{W}{U'(\gamma)}$ for every~$\gamma$, and suppose
\[\beta+6(W+\eta)+\mathsf{r}^{2}(\mathsf{n}+1)^{2}+2<p .\]
Then for every $a\in Z^{1}_{\mathrm{f}}(L)$,
\[p^{\,3\mathsf{r}^{2}(\mathsf{n}+1)^{2}+150(\eta+W+1)}\cdot a\ \in\
\cupOrth{\InCl{\mathcal{F}}{\mathsf{K}}{U}}{L}.\]}}
\leanline{\lean{tametypesexponent}}

\statement{arith}{\phantomsection\label{t:C}\textbf{Theorem C} (sparse thick letters).\enspace
\emph{In Setting~\ref{set:cut}, of whose objects only $S$, $S'$, $\Pi$, $\mathsf{thick}$, $\mathsf{s}$,
$\beta$, $\eta$, $V$, $W$, $X$, $B$, $J$ and~$G$ are used, assume of its conditions only \textup{(W1)},
\textup{(W2)}, \textup{(W5)}, \textup{(W6)}, \textup{(W7)} and~\textup{(W12)}. Suppose moreover that the Kummer chart is unramified and
\m{p\ge7}, or that $\mathsf{s}(i)=\mathsf{s}(i')$ for all letters $i$,~$i'$; that
\[G+V\le\eta,\qquad11\,(p^{J}+1)\,W\le B,\qquad4\,\bigl(J+1+\fd+\lfloor\log_p(2B)\rfloor\bigr)\,W\le X;\]
and that for every letter~$i$ and all integers \m{a,a'\in[1,J]}: if $i$ and \m{i^{p^{a}}} are thick, then
\m{i^{p^{a+a'}}} is not thick. Then $\CupB{\Pi}{10^{4}X(J+1)}{c}{c'}$ for all \m{c\in C_S} and
\m{c'\in C_{S'}}.}}
\leanline{\lean{scaledCupBoundssparsethick}}

\statement{arith}{\phantomsection\label{t:D}\textbf{Theorem D} (two sizes).\enspace
\emph{In Setting~\ref{set:cut}, of whose objects only $S$, $S'$, $\Pi$, $\mathsf{thick}$, $\mathsf{s}$,
$\beta$, $\eta$, $V$, $W$, $B$ and~$G$ are used, assume of its conditions only \textup{(W1)}, \textup{(W2)},
\textup{(W7)} and~\textup{(W12)}. Suppose moreover that the Kummer chart is unramified and \m{p\ge7}, and let \m{K_0\ge0} be
an integer with
\[5W\le G,\qquad G\le B,\qquad B<K_0\,p^{2},\qquad2G<K_0\,p,\qquad W\le K_0 .\]
Then $\CupB{\Pi}{10^{4}K_0(\lfloor\log_p(2B)\rfloor+1)}{c}{c'}$ for all \m{c\in C_S} and \m{c'\in C_{S'}}.}}
\leanline{\lean{scaledCupBoundstwosizes}}

The unit on Selmer groups uses these four statements and nothing else of the three parts. In all four the
exponent or the length does not involve the rank of a partner lattice, and in Theorem~B not the ramification
index of the field that makes the lattices crystalline.

\sect{The setting}{s:setting}
\begin{conventions}\phantomsection\label{c:conventions}
Throughout, $p$ is a prime number. The declarations of this unit take the following ambient objects as implicit or
instance arguments; they are fixed by the Settings and named again in each statement that needs them: the
prime~$p$; the group~$\Gamma$ over~$\GQp$, in the last file with a topology; the Kummer chart~$\fc$
for~$\Gamma$; the coefficient ring~$\cO$; a lattice~$L$ or a lattice datum~$L$; sides $S$, $S'$ and a
pairing~$\Pi$; a letter~$i$; exponent functions $o$, $o'$; the integers and functions of
Settings \ref{set:few} and~\ref{set:cut}. Tensor products without a subscript are over~$\Zp$;
\m{v_p(n)} is the exponent of~$p$ in an integer~$n$, with \m{v_p(0)=0}; $\lfloor\log_pn\rfloor$ is the
largest \m{a\ge0} with \m{p^{a}\le n}, and $0$ for \m{n=0}. A \kw{sign} is $1$ or~$-1$. Exponent functions
are functions $\mathbb{N}\times\mathbb{Z}\to\mathbb{Z}$. The Lean text calls the artifact's period ring
\emph{Bcris}; it is Colmez's ring~$B_{\max}$.
\end{conventions}

\begin{setting}[the base; from the background library and from Parts 1 and~2]\label{set:base}
The notation is that of Parts 1 and~2, which this Setting recalls, with some further notions of the background
library.

\emph{The chart and lattice data.} $C$ is the completed algebraic closure of~$\Qp$; \m{\Ainf\subset\WC},
$\varphi$, $\theta$ and $\BdR$ with \m{\Fil^{n}\BdR} are the period rings of Part~1. As there (Part~1,
Setting~2.1), $A_{\max}$ is the $p$-adic completion of \m{\Ainf[X]/(pX-\xi)}, where
\m{\xi=\sum_{0\le a<p}[\varepsilon^{1/p}]^{a}} generates the kernel of~$\theta$; \m{B^{+}=A_{\max}[1/p]};
and \m{B_{\max}=B^{+}[1/([\varepsilon]-1)]}. Frobenius~$\varphi$ and~$\GQp$ act on these rings; the maps
\m{A_{\max}\to B^{+}\to B_{\max}} are injective, and \m{B_{\max}} has an injective homomorphism
to~$\BdR$ that commutes with~$\GQp$ and maps \m{B^{+}} into \m{\Fil^{0}\BdR}. A lattice~$T$ is
\kw{potentially crystalline} if the elements of \m{B_{\max}\otimes T} fixed near~$1$ span it
over~$B_{\max}$, and \kw{crystalline} on a group if those fixed by the group do. $\Gamma$ is a
group with a homomorphism $j_\Gamma\colon\Gamma\to\GQp$ of open image, $\chi$ its cyclotomic character; an
element \m{s\in\Gamma} is \kw{inertial} if \m{j_\Gamma(s)} lies in the inertia group. A Kummer chart~$\fc$
for~$\Gamma$ has its ring \m{\fS\subset\Ainf}, a quotient of \m{\WW[[X]]} with $\WW=W(\mathbb{F}_{p^h})$, its
residue degree~$h$, its coordinate~$z$, its Eisenstein element~$\cE$ with $\cE_\nu=\varphi^{\nu}(\cE)$, its
index~$e_{\fc}$, the number~$\fd$ (equal to~$0$ if \m{e_{\fc}=1} and to \m{v_p(e_{\fc})+1} otherwise), and its
subgroup~$\Gamma'$ of index prime to~$p$ containing the elements near~$1$; it is \kw{unramified} if
\m{\cE=z+p}, \m{e_{\fc}=1}, \m{\fd=0}. The \kw{letters} are the elements of~$\mathbb{F}_{p^h}^\times$.
$\cO$ is the ring of integers of a finite extension of~$\Qp$; with an element \m{\omega\in\cO},
\m{\omega^{p^h-1}=1}, it is a coefficient ring for the chart. A \kw{lattice} over~$\cO$ is a finite free
$\cO$-module with an $\cO$-linear action of~$\Gamma$; $Z^{1}_{\mathrm{f}}(L)$ is the $\cO$-module of its
crystalline cocycles, and for $\gamma\colon\cO\to C$ and \m{u\in\mathbb{Z}} the jump~$u$ of~$L$ at~$\gamma$
has a multiplicity. A \kw{lattice datum} is a lattice crystalline on~$\Gamma'$ with finite sets
$U_L(i)\subset\mathbb{Z}$, one for each letter, containing its jumps at the~$\gamma$ with
$\gamma(\omega)=\theta([i])$; \m{m_L(i,u)} is the sum of the multiplicities of the jump~$u$ over
these~$\gamma$. It has the modules \m{\Minf(L)}, \m{\Mlocinf(L)}, \m{\Mloc(L)} and the Kisin module \m{\mathfrak{M}(L)}, the
periods \m{\Per(L)}, the jump bound~$h_L$, the projectors~$\pr_i$ onto the parts \m{(\WC\otimes L)_i}, and for
an exponent function~$o$ the module \m{\wExp{L}{i}{o}} of the elements of $\Mlocinf(L)\cap(\WC\otimes L)_i$
whose periods at \m{i^{p^{-\nu}}} and~$u$, after~$\varphi^{-\nu}$, lie in $\Fil^{o(\nu,u)}\BdR$; a larger~$o$
gives a smaller module. Further: \m{\std_l(\nu,u)=0} for \m{\nu<l} and \m{-u} for \m{\nu\ge l};
$\locLP{l}{i}(L)=\Mloc(L)\cap\wExp{L}{i}{\std_l}$; \m{\locKP{Q}} is the set of the elements of \m{\Mloc(L)} in
a submodule~$Q$; \m{\invMod{Q}} is the set of the~$x$ with \m{gx-x\in Q} for all \m{g\in\Gamma'};
$\cost{\nu}{n}=v_p(n)+\fd+3\nu$; $\impE{o}{o'}{\nu}{u}$ is \m{o'(\nu,u)} if \m{o(\nu,u)>0} and
\m{\min(o'(\nu,u),0)} otherwise; $\frobX{l}{n}{o}{\nu}{u}$ is~$n$ for \m{\nu<l} and \m{o(\nu-l,u)} for
\m{\nu\ge l}; $\raiseX{o}{\nu}{u}{\nu'}{v}$ is \m{o(\nu',v)+1} if \m{(\nu',v)=(\nu,u)} and \m{o(\nu',v)}
otherwise; and $\shStd{l}{o_1}{\nu}{u}$, for a function~$o_1$ of~$u$, is $\std_l(\nu,u)+o_1(u)$ for \m{\nu=0}
and \m{\std_l(\nu,u)} otherwise.

\emph{Sides and points.} A side $S=(Y_S,C_S,\alpha_S)$ over~$\cO$ consists of a lattice datum~$Y_S$ with
\m{0\in U_{Y_S}(i)} for all~$i$, a lattice~$C_S$ and an $\cO$-linear map $\alpha_S\colon C_S\to
Z^{1}_{\mathrm{f}}(Y_S)$. Its extension \m{\extL{S}} is \m{Y_S\times C_S} with $g(y,c)=(gy+\alpha_S(c)(g),c)$
and the sets \m{U_{Y_S}(i)}; \m{\trivL{S}} is $C_S$ with the trivial action and the sets~$\{0\}$;
$\inclY{S}\colon Y_S\to\extL{S}$ and $\projC{S}\colon\extL{S}\to\trivL{S}$ are the inclusion and the
projection, and \m{\projCloc{S}} is the map \m{1\otimes\projC{S}} on \m{\Mloc(\extL{S})}. For \m{c\in C_S},
$\constAt{S}{i}{c}=\pr_i(1\otimes c)$; $\liftId{S}{r}{l}{i}$ is the lifting ideal of~$\Ainf$ (Part~1,
Definition~2.8; its characterization is Input~\ref{r:ideal}(a)); $\Kills{S}{i}{\kappa}{\rho}$ means
$p^{\kappa}\cE^{\rho}\in\liftId{S}{0}{0}{i}$. \m{\kerX{S}{i}{o}} is the module of the
$x\in\wExp{\extL{S}}{i}{o}$ with $(1\otimes\projC{S})(x)=0$; \m{\deepE{S}{i}{o}} is the $\Zp$-module of
the $x\in\Mloc(\extL{S})$ such that $(1\otimes\projC{S})(x)=\constAt{S}{i}{c}$ for some \m{c\in C_S} and
$x\in\invMod{\wExp{\extL{S}}{i}{o}}$; and $\ptsOver{S}{c}{i}{e}{o_1}{o'}$ is the set of the
$x\in\Mloc(\extL{S})$ with $(1\otimes\projC{S})(x)=p^{e}\constAt{S}{i}{c}$, $x\in\wExp{\extL{S}}{i}{o_1}$
and $x\in\invMod{\wExp{\extL{S}}{i}{o'}}$, the points over~$p^{e}c$ at~$i$ with exponents~$o_1$ and
depth~$o'$. An additive map~$\mathsf{A}$ of \m{\Mloc(\extL{S})} \kw{has cost~$e$}, written
$\Reaches{S}{e}{\mathsf{A}}$, if $\mathsf{A}(x)-p^{e}x\in\locKP{\kerX{S}{j}{o}}$ for all letters~$j$,
exponent functions~$o$ and $x\in\deepE{S}{j}{o}$. A pairing~$\Pi$ of sides $S$ and~$S'$ is a pairing of
the lattices $Y_S$ and~$Y_{S'}$ (Part~2, Setting~2.2); it has its third side \m{\third{\Pi}}, whose lattice datum
is the Tate twist of the dual of $Y_S\otimes_{\cO}Y_{S'}$ with the jump sets
$\{0\}\cup\{-u-v-1\}$, \m{u\in U_{Y_S}(i)}, \m{v\in U_{Y_{S'}}(i)}; and $\CupB{\Pi}{e}{c}{c'}$ means that,
for every topology on~$\Gamma$ for which $j_\Gamma$ is an open embedding, the cup product of \m{\alpha_S(c)}
and $p^{e}\alpha_{S'}(c')$ under~$\Pi$ bounds. The windows: $\raisedE{n}{\nu}{u}$ is \m{n-u} for \m{\nu=0}
and \m{-u} for \m{\nu>0}; $\twoWin{n}{c}=[-n,n]\cup[c-n,c+n]$; $\NoMult{p}{n}{\beta}$ says that $p$ divides
no non-zero element of $\twoWin{n}{\beta}\cup\twoWin{n}{-\beta}$; and
$\JumpWin{\eta}{\mathsf{s}}{\beta}{W}{U}$ says that every \m{u\in U} has \m{|u|\le\eta} or
$\beta\le\mathsf{s} u\le\beta+W$.

\emph{The linear Frobenius.} Let $T$ be a potentially crystalline lattice, \m{s\in\Gamma} and \m{m\ge0} such
that \m{j_\Gamma(s)^{-1}} acts on the residue field of~$\bar{\mathbb{Q}}_p$ as the $p^{m}$-th power (for
\m{m=0}: $s$ is inertial). Then \m{\linF{T}{s}{m}} is the $B_{\max}$-linear endomorphism of
\m{B_{\max}\otimes T} that agrees with $\varphi^{m}j_\Gamma(s)\otimes s$ on the elements fixed near~$1$. A
polynomial~$Q$ over~$\cO$ \kw{kills \m{(s,m)} on~$T$} if $T$ is potentially crystalline and
$Q(\linF{T}{s}{m})=0$, the coefficients acting through~$T$.

\emph{Classes.} Let $\mathcal{F}$ be a filter on~$\Gamma$, $\mathsf{K}$ a relation between elements of~$\Gamma$
and polynomials over~$\cO$, and $U=(U(\gamma))_\gamma$ a family of finite sets of integers. A lattice~$Y$
over~$\cO$ lies in the \kw{class} $\InCl{\mathcal{F}}{\mathsf{K}}{U}$ if the elements of
\m{B_{\max}\otimes Y} that are fixed by $j_\Gamma(g)\otimes g$ for all~$g$ near~$1$ and for all~$g$ in some
member of~$\mathcal{F}$ span \m{B_{\max}\otimes Y}; every jump of~$Y$ at~$\gamma$ with non-zero multiplicity
lies in \m{U(\gamma)}; and $Q$ kills \m{(s,0)} on~$Y$ whenever \m{\mathsf{K}(s,Q)} and $s$ is inertial. Here
$\top$ is the filter whose only member is~$\Gamma$, so that the first condition for \m{\mathcal{F}=\top} says that
$Y$ is crystalline on~$\Gamma$, and $\bot$ is the empty relation.

\emph{Orthogonality.} A pairing of lattices $Y$ and~$L$ and the statement that the cup product of two maps
\m{a\colon\Gamma\to Y}, \m{c\colon\Gamma\to L} under it bounds, for a topology on~$\Gamma$, are as in Part~2,
Setting~2.2. For a topology on~$\Gamma$, a class~$\mathcal{Q}$ of lattices over~$\cO$ and a lattice~$L$,
$\cupOrth{\mathcal{Q}}{L}$ is the $\cO$-module of the maps \m{c\colon\Gamma\to L} such that for every
\m{Y\in\mathcal{Q}}, every pairing of $Y$ and~$L$ and every $a\in Z^{1}_{\mathrm{f}}(Y)$ the cup product of
$a$ and~$c$ bounds.

\emph{Lengths.} For a ring~$R$ and submodules $N$, $N'$ of an $R$-module, \m{\relLen{R}{N}{N'}} is the length
of the $R$-module \m{N/(N\cap N')}, a natural number or~$\infty$.
\end{setting}

\begin{definition}[constants that lift to deep points]\label{d:dc}
In Setting~\ref{set:base}, for a side~$S$ over~$\cO$, a letter~$i$, an integer~$n$ and an exponent
function~$o$, \m{\deepC{S}{i}{n}{o}} is the image under \m{\projCloc{S}} of the $\Zp$-module of the
$x\in\Mloc(\extL{S})$ with $x\in\wExp{\extL{S}}{i}{\raisedE{n}{\cdot}{\cdot}}$ and
$x\in\deepE{S}{i}{o}$; it is a $\Zp$-submodule of \m{\WC\otimes C_S}.
\end{definition}
\leanline{\lean{deepConstants}}

\begin{definition}[tame charts]\label{d:tame}
In Setting~\ref{set:base}, a \kw{tame chart} for~$\Gamma$ is a Kummer chart~$\fc$ for~$\Gamma$ together with: an
element \m{\varrho\in\WW} whose reduction modulo~$p$ is a primitive $e_{\fc}$-th root of unity; an inertial
element \m{\sigma\in\Gamma}, the \kw{tame generator}, such that for every \m{k\ge0}, \m{\sigma^{k}} fixes~$\WW$
and $\sigma^{k}z=\varrho^{k}z$; and an element $\varpi\in\bar{\mathbb{Q}}_p$ with \m{\varpi^{e_{\fc}}=p}
that is fixed by \m{j_\Gamma(g)} for all \m{g\in\Gamma'}; such that $p$ does not divide~$e_{\fc}$ and
$\Gamma'$ is normal in~$\Gamma$.
\end{definition}
\leanline{\lean{TameChart}}

\begin{definition}[tame types; the universal side]\label{d:univ}
In Setting~\ref{set:base}:
\begin{enumerate}[label=\textup{(\roman*)},leftmargin=2.3em,itemsep=2pt]
\item For a natural number~$p$, integers $\mathsf{r},\mathsf{n}\ge0$ and a property~$\mathsf{K}_0$ of
polynomials over~$\cO$, $\TameType{p}{\mathsf{r}}{\mathsf{n}}{\mathsf{K}_0}$ means that $\mathsf{K}_0$ holds
for \m{X^{M_0}-1}, where $M_0$ is the least common multiple of the \m{p^{a}-1} for $1\le a\le\mathsf{r}$,
and for the image in \m{\cO[X]} of some monic polynomial over~$\Zp$ of degree at most~$\mathsf{n}$.
\item For a Kummer chart~$\fc$, \m{\omega\in\cO} with \m{\omega^{p^h-1}=1}, a family
$U=(U(\gamma))_\gamma$ of finite sets of integers and a letter~$i$,
\[\letterJ{\fc}{\omega}{U}{i}=\{0\}\cup\bigcup_{\gamma(\omega)=\theta([i])}U(\gamma).\]
\item Let $\fc$ and~$\omega$ be as in~(ii), $\mathcal{F}$ a filter on~$\Gamma$ every member of which
contains~$\Gamma'$, and $L$ a lattice with $L\in\InCl{\mathcal{F}}{\bot}{U}$. The \kw{universal side}
$\univS{\fc}{\omega}{\mathcal{F}}{L;U}$ is the side over \m{(\cO,\omega)} whose lattice datum is~$L$ with the sets
$\letterJ{\fc}{\omega}{U}{i}$, whose lattice is $Z^{1}_{\mathrm{f}}(L)$ with the trivial action, and whose
cocycle map is the identity. (It is a lattice datum because $L$ is crystalline on~$\Gamma'$; and
$Z^{1}_{\mathrm{f}}(L)$ is finitely generated by Input~\ref{k:bk} and has no torsion, so it is free.)
\end{enumerate}
\end{definition}
\leanline{\lean{IsTameType}\sep \lean{letterJumps}\sep \lean{universalSide}}

\begin{setting}[few low weights]\label{set:few}
In Setting~\ref{set:base}, the Kummer chart is unramified, $S$ is a side over~$\cO$, \m{\rho,\eta,W\ge0} are
integers, and for every letter~$i$ there are integers $\beta(i),n_{\mathrm{lo}}(i),n_{\mathrm{all}}(i)\ge0$
and a sign~\m{\mathsf{s}(i)}. The following conditions are referred to by their names.
\begin{enumerate}[label=\textup{(F\arabic*)},leftmargin=2.9em,itemsep=2pt]
\item \m{p\ge5}.
\item $\beta(i)+W+\eta\le(p-3)\rho$ for every letter~$i$.
\item $6(\rho+\eta+W)<\beta(i)$ for every letter~$i$.
\item $\NoMult{p}{6(\rho+\eta+W)}{\beta(i)}$ for every letter~$i$.
\item $\JumpWin{\eta}{\mathsf{s}(i)}{\beta(i)}{W}{U_{Y_S}(i)}$ for every letter~$i$.
\item $\sum_{u\in[-\eta,\eta]}h\,m_{Y_S}(i,u)\le n_{\mathrm{lo}}(i)$ and
$\sum_{u\in U_{Y_S}(i)}h\,m_{Y_S}(i,u)\le n_{\mathrm{all}}(i)$ for every letter~$i$.
\end{enumerate}
\end{setting}

\begin{setting}[cuts and windows]\label{set:cut}
In Setting~\ref{set:base}, $S$ and~$S'$ are sides over~$\cO$ and $\Pi$ is a pairing of $S$ and~$S'$;
$\mathsf{cut}$ and $\mathsf{thick}$ are properties of letters; for every letter~$i$ there are a
sign~\m{\mathsf{s}(i)} and integers $\rho(i),\beta(i)\ge0$; and
$\eta,V,W,X,\kappa,B,n,J,\rho_0,G\ge0$ are integers. The following conditions are referred to by their names;
each statement says which of them it assumes.
\begin{enumerate}[label=\textup{(W\arabic*)},leftmargin=3.2em,itemsep=2pt]
\item \m{V+4\eta\le W}.
\item $\JumpWin{\eta}{-\mathsf{s}(i)}{\beta(i)}{V}{U_{Y_S}(i)}$ and
$\JumpWin{\eta}{\mathsf{s}(i)}{\beta(i)}{V}{U_{Y_{S'}}(i)}$ for every letter~$i$.
\item $\bigl(J+1+\fd+\lfloor\log_p(2B)\rfloor\bigr)\,n+\fd\,G\le X$.
\item $12\,(J+1)\,X<\kappa$.
\item $e_{\fc}\cdot2B\le p^{J+1}X$.
\item $8\,e_{\fc}\le p^{J+1}$.
\item \m{\eta\ge1}.
\item \m{W\le n}, \m{W\le G}, and $\beta(i)+\rho(i)+W\le2B$ for every letter~$i$.
\item \m{G\le(p-1)\,n}.
\item For every letter~$i$ with the property~$\mathsf{cut}$: either \m{\rho(i)=0} and
$\beta(i^{p^{a}})+W\le G$ for all \m{a\in[1,J]}; or \m{\rho(i)<\beta(i)}, the chart is unramified,
\m{p\ge5} and $\rho_0\le\kappa+\rho(i)$.
\item If \m{p\ge5}, then $\beta(i)+W\le(p-3)\rho_0$ for every letter~$i$.
\item \m{\beta(i)=B} if $i$ is thick, and \m{\beta(i)=G} otherwise.
\end{enumerate}
\end{setting}

\inputbox{\textbf{Inputs from Parts 1 and~2.} From Part~1 (\emph{Exponents of elements of Kisin modules, the cost
of deepening them, and cokernel killers}) and Part~2 (\emph{Cup products of crystalline classes vanish up to a
power of~$p$: the walk along Frobenius orbits}); $L$ is a lattice datum, $S$ a side.
\begin{enumerate}[label=\textup{Q\arabic*.},ref=Q\arabic*,leftmargin=2.8em,itemsep=2.5pt]
\item\label{r:exp} \emph{Exponent modules} (Part~1, Lemmas 4.1, 4.2, Theorem~B, Lemma~7.3). Monotonicity in
the exponents; \m{\cE_{\nu_0}^{\,n}}, for an integer~$n$ of either sign, raises the exponents at~$\nu_0$ by~$n$ and preserves \m{\Mloc(L)};
if $rx\in\wExp{L}{i}{o}$ with \m{x\in\Mlocinf(L)} and \m{\cE_\nu\nmid r} for all~$\nu$, then
\m{x\in\wExp{L}{i}{o}}; $gx-x\in\wExp{L}{i}{o}$ for \m{x\in\wExp{L}{i}{o}}, \m{g\in\Gamma'};
$x\in\wExp{L}{i}{\std_0}$ if and only if $x\in\Minf(L)\cap(\WC\otimes L)_i$; $\pr_i(x)\in\wExp{L}{i}{\std_0}$ for \m{x\in\Minf(L)}; the map \m{1\otimes\inclY{S}} sends
\m{\wExp{Y_S}{i}{o}} into $\wExp{\extL{S}}{i}{o}$ and \m{\Mloc(Y_S)} into \m{\Mloc(\extL{S})}; and
$\cost{\nu}{n}\le3\nu$ if the chart is unramified and $p$ does not divide~$n$ or \m{n=0} (Part~1,
Lemma~6.1).
\item\label{r:deep} \emph{Deep elements} (Part~1, Lemma~6.2). If \m{x\in\wExp{L}{i}{o}} and
$x\in\invMod{\wExp{L}{i}{o'}}$, then $x\in\wExp{L}{i}{\impE{o}{o'}{\cdot}{\cdot}}$.
\item\label{r:ideal} \emph{The lifting ideal} (Part~1, Lemmas 9.1, 9.3 to 9.7, Theorem~D). \emph{(a)}
$s\in\liftId{S}{r}{l}{i}$ if and only if there is \m{t\in\Ainf} such that every $(s-rt)\constAt{S}{i}{c}$
is the image of an element of $\wExp{\extL{S}}{i}{\std_l}$. \emph{(b)} If
$p^{c}\in\liftId{S}{p^{c+1}}{l}{i}$, then $\liftId{S}{p^{c+1}}{l}{i}\subseteq\liftId{S}{0}{l}{i}$.
\emph{(c)} $p^{c}z^{pa}\in\liftId{S}{r}{1}{i^{p}}$ if and only if $p^{c}z^{a}\in\liftId{S}{r}{0}{i}$;
if $s\in\liftId{S}{r}{0}{i}$ and \m{u\le n} for all \m{u\in U_{Y_S}(i)}, then
$s\,\cE^{n}\in\liftId{S}{r}{1}{i}$. \emph{(d)} If \m{f\in\fS} lies in $\liftId{S}{r}{l}{i}$, there is
\m{t\in\fS} such that every $y\in\locLP{l}{i}(\trivL{S})$ has an $x\in\locLP{l}{i}(\extL{S})$ over
\m{(f-rt)y}; and $\constAt{S}{i}{c}\in\locLP{l}{i}(\trivL{S})$. \emph{(e)} If \m{pa\le n},
$p^{c}z^{n}\in\liftId{S}{r}{1}{i^{p}}$ and $c+v_p(n!)\le\kappa+v_p((pa)!)$, then
$p^{\kappa}z^{a}\in\liftId{S}{r}{0}{i}$. \emph{(f)} If $p^{c}z^{a}\cE^{n}\in\liftId{S}{r}{1}{i}$,
$c+\fd n+v_p(n!)\le b$ and \m{a+n\le\kappa}, then $p^{b}z^{\kappa}\in\liftId{S}{r}{1}{i}$. \emph{(g)} If
\m{u\le n} for all letters~$j$ and \m{u\in U_{Y_S}(j)}, then for every \m{\kappa\ge0} there is \m{a\ge0} with
$z^{a}\in\liftId{S}{p^{\kappa}}{0}{j}$ for all~$j$ and $a\,p\le e_{\fc}(n+\kappa)+a$. \emph{(h)} If the
chart is unramified, \m{p\ge5} and $u\le(p-3)(\kappa+\rho)$ for all~$j$ and \m{u\in U_{Y_S}(j)}, then
$\Kills{S}{i}{\kappa}{\rho}$. \emph{(i)} $\sum_{a\in\twoWin{n}{c}}(v_p(a)+d)\le(4m+2)\lfloor\log_pH\rfloor+(4n+2)d$
if \m{n\le pm} and \m{|c|+n\le H}.
\end{enumerate}}

\inputbox{\textbf{Inputs from Parts 1 and~2, continued.}
\begin{enumerate}[label=\textup{Q\arabic*.},ref=Q\arabic*,leftmargin=2.8em,itemsep=2.5pt,start=4]
\item\label{r:points} \emph{Points} (Part~2, Lemmas 4.1, 4.4, 4.5, 4.7, Proposition~4.6). \emph{(a)} If
$\Kills{S}{i}{\kappa}{\rho}$, then $\ptsOver{S}{c}{i}{\kappa}{\raisedE{-\rho}{\cdot}{\cdot}}{\raisedE{-\rho}{\cdot}{\cdot}}$
is not empty. \emph{(b)} For a finite set~$I$ of steps \m{(\nu,n)} there is an additive map~$\mathsf{A}$ of
\m{\Mloc(\extL{S})} with $\mathsf{A}(x)-p^{\epsilon}x\in\locKP{\kerX{S}{i}{o}}$ for
$x\in\deepE{S}{i}{o}$, where $\epsilon=\sum_{(\nu,n)\in I}\cost{\nu}{n}$, and
$\mathsf{A}(x)\in\deepE{S}{i}{o'}$ if $I$ contains all steps from $o$ to~$o'$; such a map preserves every
$\Zp$-submodule of \m{\Mloc(\extL{S})} that contains $\locKP{\kerX{S}{i}{o}}$, on its elements in
\m{\deepE{S}{i}{o}}; and $\locKP{\kerX{S}{i}{o}}\subseteq\deepE{S}{i}{o}$. \emph{(c)} A point of
$\ptsOver{S}{c}{i}{e}{o_1}{o}$ gives a point of $\ptsOver{S}{c}{i}{e'}{\impE{o_1}{o'}{\cdot}{\cdot}}{o'}$
if $I$ contains the steps from $o$ to~$o'$ and \m{e+\epsilon\le e'}; in particular, if $\NoMult{p}{n}{\beta}$
and all these steps are at \m{\nu=0} and in $\twoWin{n}{\beta}\cup\twoWin{n}{-\beta}$, one may take
\m{e'=e+(8n+4)\fd}. \emph{(d)} $\varphi^{l}$ maps $\ptsOver{S}{c}{i^{p^{-l}}}{e}{o}{o'}$ into
$\ptsOver{S}{c}{i}{e}{\frobX{l}{0}{o}{\cdot}{\cdot}}{\frobX{l}{1}{o'}{\cdot}{\cdot}}$. \emph{(e)}
(Strictness.) Every \m{w\in\kerX{S}{i}{o}} is the image of a $y\in\wExp{Y_S}{i}{o}$, and \m{y\in\Mloc(Y_S)}
if $w\in\Mloc(\extL{S})$.
\item\label{r:walk} \emph{The walk} (Part~2, Theorems B, C and~D, Definition~2.6). Theorem~B: deep points of
$S$ and~$S'$ against each other at every cut, with a point of the third side, give
$\CupB{\Pi}{E}{c}{c'}$. Theorem~C: the rank-free exponent
$H\ge3(\kappa+(48(\rho+\eta+W)+4)\fd)$ for sides with jumps in opposite windows and three cokernels killed
by $p^{\kappa}\cE^{\rho}$. Theorem~D: the long gap. They are cited with their hypotheses where used.
\item\label{r:auer} \emph{A relation among elements of one weight} (Part~2, Proposition~7.1). Unramified
chart; \m{m_L(i,u)\le n}; \m{j^{p^{-\nu_0}}=i}; $o(\nu_0,u)<o(\nu_0,v)$ for \m{v\in U_L(i)}, \m{v<u};
$x_0,\dots,x_{hn}\in\Mloc(L)\cap\wExp{L}{j}{o}$. Then some combination \m{\sum_kc_kx_k} with \m{c_k\in\Zp},
one of them~$1$, lies in $\locKP{\wExp{L}{j}{o_+}}+\locKP{\wExp{L}{j}{\raiseX{o}{\nu_0}{u}{\cdot}{\cdot}}}$,
where $o_+$ is $o$ raised by~$1$ at \m{\nu=0}.
\end{enumerate}}

\inputbox{\textbf{Inputs from the background library.} Each is a declaration of the background library, read in its
Lean text; the number in brackets is the unit that owns it.
\begin{enumerate}[label=\textup{K\arabic*.},ref=K\arabic*,leftmargin=2.8em,itemsep=2.5pt]
\item\label{k:len} \emph{Relative lengths} [57]. For $\Zp$-submodules: \emph{(a)}
$\relLen{\Zp}{N}{N''}\le\relLen{\Zp}{N}{N'}+\relLen{\Zp}{N'}{N''}$; \emph{(b)}
$\relLen{\Zp}{N}{f^{-1}(N')}=\relLen{\Zp}{f(N)}{N'}$ for a linear map~$f$; \emph{(c)}
$\relLen{\Zp}{N}{\bigcap_kN_k}\le\sum_k\relLen{\Zp}{N}{N_k}$ for finitely many~$N_k$; \emph{(d)} if
\m{p^{a}N\subseteq N'} and any \m{m+1} elements~$x_k$ of~$N$ with \m{px_k\in N'} have a combination
\m{\sum_kc_kx_k\in N'} with \m{c_k\in\Zp}, one of them~$1$, then $\relLen{\Zp}{N}{N'}\le a\,m$. Elementary.
\item\label{k:frob} \emph{The linear Frobenius} [86]. For potentially crystalline lattices: \emph{(a)}
$\linF{T}{st}{m+m'}=\linF{T}{s}{m}\,\linF{T}{t}{m'}$, and \m{\linF{T}{g}{0}=1} for inertial~$g$ near~$1$;
\emph{(b)} a $\Zp$-linear map \m{f\colon T\to T'} that commutes with~$s$ and with all~$g$ near~$1$ satisfies
$(1\otimes f)\,\linF{T}{s}{m}=\linF{T'}{s}{m}\,(1\otimes f)$; \m{\linF{T}{s}{0}=1} for inertial~$s$ if $T$ is
crystalline on~$\Gamma$; \emph{(c)} an element of~$B_{\max}$ fixed near~$1$ is fixed by every inertial
element. In the second reader's comparison, Fontaine gives the Weil action on the periods of a
finite-dimensional representation (J.-M.~Fontaine, \emph{Représentations $\ell$-adiques potentiellement
semi-stables}, Astérisque~223, Exposé~VIII, §§2.3.2--2.3.3, §2.3.7), and the $B_{\max}$-linear extension on
arbitrary modules is not located. \emph{Relation}: not settled; the existence of the linear extension is
part of what the formal library proves, and no published statement of it in this generality is known to us.
\item\label{k:crit} \emph{A crystalline criterion} [86]. If $T$ is potentially crystalline and
\m{\linF{T}{s}{0}=1} for every inertial \m{s\in\Gamma}, then $T$ is crystalline on~$\Gamma$. \emph{What is
quoted} is this statement of the formal library, for any group~$\Gamma$ over~$\GQp$ of open image, any
module~$T$ over a commutative $\Zp$-algebra, and with $B_{\max}$ as period ring. \emph{Relation to the
literature}: a weaker form of Fontaine's criterion, namely only the implication from ``inertia acts trivially
on the potentially crystalline periods'' to ``crystalline''; the place is not confirmed. In the second
reader's comparison, unramified descent and the Weil inertia account for the classical criterion for
finite-dimensional representations (Fontaine, \emph{Représentations $p$-adiques semi-stables},
Astérisque~223, Exposé~III, §§1.8.5--1.8.6, §5.6.8(iii); Exposé~VIII, §2.3), and a published descent for
arbitrary modules and for the spanning condition over~$B_{\max}$ is not located. So the Input is not a
citation of a published theorem: it is a statement that the formal library proves, close to a classical one.
\end{enumerate}}

\inputbox{\textbf{Inputs from the background library, continued.}
\begin{enumerate}[label=\textup{K\arabic*.},ref=K\arabic*,leftmargin=2.8em,itemsep=2.5pt,start=4]
\item\label{k:dual} \emph{Duals of tensor products} [86]. Let $s$ be inertial, $Y$ and~$L$ lattices in classes
with relations $\mathsf{K}$, $\mathsf{K}'$, and $q$, $q'$ monic polynomials over~$\Zp$ with
\m{\mathsf{K}(s,q)}, \m{\mathsf{K}'(s,q')}. Then a monic polynomial over~$\Zp$ of degree at most
\m{\deg q\cdot\deg q'} kills \m{(s^{-1},0)} on $(Y\otimes_{\cO}L)^{\vee}(1)$. Reference to be confirmed.
\item\label{k:bk} \emph{Crystalline cocycles and jumps} [86, 87]. For a potentially crystalline lattice~$T$:
$Z^{1}_{\mathrm{f}}(T)$ is a finite $\cO$-module, and for each~$\gamma$ the multiplicities of the jumps
of~$T$ at~$\gamma$ have sum $\operatorname{rk}_{\cO}(T)$. For \m{u\ne0} the multiplicity of the jump~$u$ of
\m{\extL{S}} at~$\gamma$ is at most that of~$Y_S$. \emph{What is quoted} for the first statement is the
finite generation, over the coefficient ring, of the module of the \emph{integral} crystalline cocycles
defined with~$B_{\max}$; the formal library deduces it from a rank formula for that module, which it proves
itself. \emph{Relation to the literature}: the Bloch--Kato local dimension formula (S.~Bloch, K.~Kato,
\emph{$L$-functions and Tamagawa numbers of motives}, §3.7.2--3, Corollary~3.8.4) concerns $H^{1}_{\mathrm{f}}$
of the rational representation, defined with~$B_{\mathrm{cris}}$. The passage from it to the module of
integral cocycles with~$B_{\max}$ needs three things: the count of the coboundaries, the comparison of the
invariants over $B_{\max}$ and over~$B_{\mathrm{cris}}$ (Colmez, \emph{Théorie d'Iwasawa des représentations
de de~Rham d'un corps local}, §III.4, for finite-dimensional representations), and the continuity of
crystalline cocycles. The second reader sketches this bridge and marks it as not located in the literature;
it is not made here. So the first statement is a special case of the Bloch--Kato formula only after that
bridge. Of it this note uses the finite generation alone (Definition~\ref{d:univ}); the rank formula is not
used. The second statement is classical for field coefficients (Fontaine, Exposé~III, Theorem~3.8(i),
Remark~3.9(ii), §5.1.7), its formulation for coefficient fibres not located; the third is elementary.
\item\label{k:tame} \emph{Tame inertia} [89, 92]. Let $e$ be prime to~$p$, $y\in\bar{\mathbb{Q}}_p$ with
$y^{e}$ a uniformizer of~$\Qp$, and $f$ a homomorphism from the inertia group of~$\GQp$ to a monoid, with open
kernel and \m{f(g)^{e}=1} for all~$g$. Then $f$ is trivial on the stabilizer of~$y$. The file names it Hilbert's
ramification theory for infinite Galois extensions; reference to be confirmed.
\end{enumerate}}

\inputbox{\textbf{Inputs from the background library, concluded.}
\begin{enumerate}[label=\textup{K\arabic*.},ref=K\arabic*,leftmargin=2.8em,itemsep=2.5pt,start=7]
\item\label{k:frame} \emph{The frame theorem with a divisor; Galois and the Kisin module} [83, 82]. Let $L$ be
a lattice datum. A \kw{crystalline period} is a period \m{\mu\in\Per(L)} whose values on \m{1\otimes L} are
given by a $\Zp$-linear map \m{L\to B_{\max}}. \emph{(a)} Let \m{d\in\Ainf} be such that $d$ divides~$r$
in~$\Ainf$ whenever $d$ divides \m{\cE^{m}r} in~$B^{+}$. If \m{x\in\Minf(L)} and
$\cE^{2h_L}\mu(x)\in B^{+}d$ for every crystalline period~$\mu$, then \m{x\in d\,\Minf(L)}. The
coordinate~$z$ is such a~$d$. \emph{(b)} \m{\cE^{h_L}\mu(x)} lies in the image of~$\fS_B$ for every
crystalline period~$\mu$ and $x\in\mathfrak{M}(L)$. \emph{(c)} If \m{g\in\GQp} fixes~$\WW$, then $z$ divides
\m{gs-s} in~$B^{+}$ for every \m{s\in\fS_B}. \emph{(d)} Let \m{g\in\GQp}, \m{d\in\Ainf} and \m{m\ge0}; if $d$
divides \m{gs-s} in~$B^{+}$ for all \m{s\in\fS_B} and \m{b\in\BdR} has \m{\cE^{m}b} in the image of~$\fS_B$,
then $\cE^{2m}(gb-b)\in B^{+}d$. \emph{(e)} For \m{g\in\Gamma} with $g\Gamma'g^{-1}\subseteq\Gamma'$:
$g\,\Minf(L)\subseteq\Minf(L)$, and $g^{-1}\circ\mu\circ g\in\Per(L)$ for every \m{\mu\in\Per(L)}.
\emph{(f)} A $\Zp$-linear combination of periods that vanishes on \m{1\otimes L} vanishes on \m{\Minf(L)}.
Here $\fS_B$ is the ring of Part~1, Setting~2.1. Reference to be confirmed (the Kisin lattice of a crystalline
lattice).
\item\label{k:chart} \emph{Existence of charts} [83]. \emph{(a)} For \m{h\ge1}, \m{x\in\mathcal{O}_C} and a
unit~$v$ with \m{x^{e}v=-p} there is a chart with residue degree~$h$, index~$e$, coordinate~$z$ with
\m{\theta(z)=x}, and constant unit series. \emph{(b)} If a closed subgroup of~$\GQp$ acts transitively on the
$n$-th roots of~$p$ for all \m{n>0}, and $u,v\in\mathcal{O}_C^\flat$ have untilts whose $e$-th power is~$p$,
then some element of the subgroup sends $[u]$ to~$[v]$. The inertia group has this property. \emph{(c)} Every
\m{s\in\fS} is \m{w+zt} with \m{w\in\WW}, \m{t\in\fS}; and if $z$ divides \m{\cE^{a}r} in~$B^{+}$, then $z$
divides~$r$ in~$\Ainf$. Elementary.
\end{enumerate}}

\sect{Constants that lift to deep points}{s:few}
Throughout this section the setting is Setting~\ref{set:few}. The chart is unramified, so \m{\cE=z+p},
\m{\fd=0} and $\cost{\nu}{n}=v_p(n)+3\nu$. For a letter~$i$ and an exponent function~$o$ write
\[\mathsf{P}_i(o)=\locKP{\wExp{\extL{S}}{i}{\raisedE{0}{\cdot}{\cdot}}}\cap\deepE{S}{i}{o},\qquad
\mathsf{P}^{0}_i(o)=\locKP{\kerX{S}{i}{o}},\]
so that $\deepC{S}{i}{0}{o}=\projCloc{S}(\mathsf{P}_i(o))$ by Definition~\ref{d:dc}. Write $i^-$ for the
letter \m{i^{p^{-1}}}.

\begin{lemma}[deepening at one Frobenius index: the loss of constants]\label{l:rel}
In Setting~\ref{set:few}, assume of its conditions only~\textup{(F6)}. Let $i$, $i'$ be letters and
\m{\nu_0\ge0} an integer with \m{i^{p^{-\nu_0}}=i'}, let $o$, $o'$ be exponent functions and \m{A,B\ge0}
integers. Suppose that:
\begin{enumerate}[label=\textup{(\roman*)},leftmargin=2.3em,itemsep=2pt]
\item for every \m{v\in\mathbb{Z}}, $p$ divides no integer~$n$ with $o(0,v)<n<o(0,v)+p^{\nu_0}$;
\item for every \m{\nu\ge0} and $u\in U_{Y_S}(i^{p^{-\nu}})$: \m{-u\le o(\nu,u)}, and
$o'(\nu,u)\le o(\nu,u)$ if \m{\nu\ne\nu_0};
\item the function $u\mapsto o'(\nu_0,u)$ is non-increasing on \m{U_{Y_S}(i')};
\item for every \m{u\in U_{Y_S}(i')}: $\sum_{o(\nu_0,u)\le n<o'(\nu_0,u)}\cost{\nu_0}{n}\le A+B_u$, where
\m{B_u=B} if \m{|u|\le\eta} and \m{B_u=0} otherwise.
\end{enumerate}
Then
\[\relLen{\Zp}{\deepC{S}{i}{0}{o}}{\deepC{S}{i}{0}{o'}}\ \le\ A\cdot n_{\mathrm{all}}(i')+B\cdot n_{\mathrm{lo}}(i').\]
\end{lemma}
\leanline{\lean{relLengthdeepConstantsle}}
\begin{proof}
Multiplying (iv) by \m{h\,m_{Y_S}(i',u)} and summing over \m{u\in U_{Y_S}(i')} gives, by~(F6),
\[\begin{gathered}\sum_{(u,n)\in I}h\,m_{Y_S}(i',u)\,\cost{\nu_0}{n}\ \le\ A\,n_{\mathrm{all}}(i')+B\,n_{\mathrm{lo}}(i'),\\
I=\{(u,n): u\in U_{Y_S}(i'),\ o(\nu_0,u)\le n<o'(\nu_0,u)\}.\end{gathered}\]
We prove that the relative length is at most the left side, by induction on the number of elements of~$I$,
for all~$o$ at once.

If no \m{v\in U_{Y_S}(i')} has $o(\nu_0,v)<o'(\nu_0,v)$, then \m{o'\le o} on all jump sets by~(ii), so
$\deepE{S}{i}{o}\subseteq\deepE{S}{i}{o'}$ (Input~\ref{r:exp}) and the relative length is~$0$. Otherwise let
$u$ be the least such~$v$ and let \m{o^{+}} be the function $\raiseX{o}{\nu_0}{u}{\cdot}{\cdot}$. The
hypotheses hold for \m{(o^{+},o')}: (i) because \m{o^{+}(0,\cdot)} differs from \m{o(0,\cdot)} only if
\m{\nu_0=0}, when the intervals of~(i) are empty; (ii), (iii) and~(iv) at once. Its set of steps is
$I\smallsetminus\{(u,o(\nu_0,u))\}$. By Input~\ref{k:len}(a) and the induction hypothesis it remains to show
\begin{equation}\label{e:onestep}
\relLen{\Zp}{\deepC{S}{i}{0}{o}}{\deepC{S}{i}{0}{o^{+}}}\ \le\ \cost{\nu_0}{o(\nu_0,u)}\cdot h\,m_{Y_S}(i',u).
\end{equation}
By (ii) and Input~\ref{r:exp}, $\mathsf{P}^{0}_i(o)\subseteq\locKP{\wExp{\extL{S}}{i}{\raisedE{0}{\cdot}{\cdot}}}$,
and $\mathsf{P}^{0}_i(o)\subseteq\deepE{S}{i}{o}$ (Input~\ref{r:points}(b)); so
$\mathsf{P}^{0}_i(o)\subseteq\mathsf{P}_i(o)$, and \m{\projCloc{S}} vanishes on~$\mathsf{P}^{0}_i(o)$. By
Input~\ref{k:len}(b) the left side of~\eqref{e:onestep} is the relative length of \m{\mathsf{P}_i(o)} and the
preimage of $\projCloc{S}(\mathsf{P}_i(o^{+}))$, which contains $N'=\mathsf{P}_i(o^{+})+\mathsf{P}^{0}_i(o)$. We
apply Input~\ref{k:len}(d) to \m{N=\mathsf{P}_i(o)} and~$N'$, with the exponent
$\cost{\nu_0}{o(\nu_0,u)}$ and with \m{m=h\,m_{Y_S}(i',u)}.

\emph{The power of~$p$.} Let $x\in\mathsf{P}_i(o)$. The only step from $o$ to~$o^{+}$ is
\m{(\nu_0,o(\nu_0,u))}. Let $\mathsf{A}$ be the additive map of Input~\ref{r:points}(b) for this one step, of
cost $\epsilon=\cost{\nu_0}{o(\nu_0,u)}$. Then $\mathsf{A}(x)\in\deepE{S}{i}{o^{+}}$,
$\mathsf{A}(x)-p^{\epsilon}x\in\mathsf{P}^{0}_i(o)$, and \m{\mathsf{A}(x)} lies in
$\locKP{\wExp{\extL{S}}{i}{\raisedE{0}{\cdot}{\cdot}}}$, a $\Zp$-module that contains $\mathsf{P}^{0}_i(o)$.
So $p^{\epsilon}x\in N'$.

\emph{The relations.} Let $x_0,\dots,x_m\in\mathsf{P}_i(o)$ with \m{p\,x_k\in N'}, and write
\m{p\,x_k=x'_k+w_k} with $x'_k\in\mathsf{P}_i(o^{+})$ and $w_k\in\mathsf{P}^{0}_i(o)$. By strictness
(Input~\ref{r:points}(e)), $w_k=(1\otimes\inclY{S})(y_k)$ with $y_k\in\Mloc(Y_S)\cap\wExp{Y_S}{i}{o}$. By
the minimality of~$u$ and by~(iii), every \m{v\in U_{Y_S}(i')} with \m{v<u} has
$o(\nu_0,v)\ge o'(\nu_0,v)\ge o'(\nu_0,u)>o(\nu_0,u)$. So Input~\ref{r:auer} applies to the lattice
datum~$Y_S$, the letters $i'$ and~$i$, the jump~$u$ with \m{n=m_{Y_S}(i',u)}, and the elements~$y_k$: there are
\m{c_k\in\Zp}, one of them equal to~$1$, and
\[\sum_kc_ky_k=y'+y'',\qquad y'\in\locKP{\wExp{Y_S}{i}{o_+}},\quad y''\in\locKP{\wExp{Y_S}{i}{o^{+}}},\]
where $o_+$ is $o$ raised by~$1$ at \m{\nu=0}. Put $x^{\ast}=\sum_kc_kx_k\in\mathsf{P}_i(o)$ and
$w=(1\otimes\inclY{S})(y')\in\mathsf{P}^{0}_i(o_+)$. We show \m{x^{\ast}\in N'}.

First, $p\,x^{\ast}-w\in\deepE{S}{i}{o^{+}}$: it lies in \m{\Mloc(\extL{S})} over a constant, and
\[p\,x^{\ast}-w=\sum_kc_kx'_k+(1\otimes\inclY{S})(y''),\]
where each $x'_k$ is invariant modulo $\wExp{\extL{S}}{i}{o^{+}}$ and the last term lies in
$\wExp{\extL{S}}{i}{o^{+}}$, a module stable under \m{x\mapsto gx-x} for \m{g\in\Gamma'}
(Input~\ref{r:exp}).

Secondly, let $I_2$ be the set of the pairs \m{(0,n)} with $o(0,v)<n<o(0,v)+p^{\nu_0}$ for some
\m{v\in U_{Y_S}(i)}. By~(i) each of them has $\cost{0}{n}=v_p(n)=0$; so the additive map~$\mathsf{B}$ of
Input~\ref{r:points}(b) for~$I_2$ has cost~$0$. Since $w\in\mathsf{P}^{0}_i(o_+)\subseteq\deepE{S}{i}{o_+}$, the
element \m{\mathsf{B}(w)} lies in $\mathsf{P}^{0}_i(o_+)$ and in \m{\deepE{S}{i}{o_2}}, where
$o_2(0,v)=o(0,v)+p^{\nu_0}$ and $o_2(\nu,v)=o(\nu,v)$ for \m{\nu>0}: the steps from $o_+$ to~$o_2$ are
in~$I_2$. By Input~\ref{r:deep}, $\mathsf{B}(w)\in\wExp{\extL{S}}{i}{\impE{o_+}{o_2}{\cdot}{\cdot}}$. Now
$\impE{o_+}{o_2}{0}{v}-p^{\nu_0}\ge o(0,v)$: this is clear if \m{o(0,v)\ge0} or
$o(0,v)+p^{\nu_0}\le0$, and the remaining case $o(0,v)<0<o(0,v)+p^{\nu_0}$ is excluded by~(i), as $p$
divides~$0$. Hence, by Input~\ref{r:exp},
\[w_3=\cE^{-p^{\nu_0}}\,\mathsf{B}(w)\ \in\ \mathsf{P}^{0}_i(o).\]
Since \m{\cE=z+p} and $\cE_{\nu_0}=z^{p^{\nu_0}}+p$, we have $\cE^{p^{\nu_0}}=\cE_{\nu_0}+p\,r$ with
\m{r\in\fS}. Put $w_4=r\,w_3-(\mathsf{B}(x^{\ast})-x^{\ast})\in\mathsf{P}^{0}_i(o)$; the second term lies in
$\mathsf{P}^{0}_i(o)$ because $\mathsf{B}$ has cost~$0$. As $\mathsf{B}$ is additive,
\[p\,(x^{\ast}-w_4)=p\,\mathsf{B}(x^{\ast})-p\,r\,w_3=\mathsf{B}(p\,x^{\ast}-w)+\cE^{p^{\nu_0}}w_3-p\,r\,w_3
=\mathsf{B}(p\,x^{\ast}-w)+\cE_{\nu_0}\,w_3 .\]
Here $\mathsf{B}(p\,x^{\ast}-w)\in\deepE{S}{i}{o^{+}}$ (Input~\ref{r:points}(b), for the module
$\deepE{S}{i}{o^{+}}$ itself), and $\cE_{\nu_0}w_3\in\wExp{\extL{S}}{i}{o^{+}}$ because multiplication
by~$\cE_{\nu_0}$ raises the exponents at~$\nu_0$ by~$1$. So for every \m{g\in\Gamma'} the element
$p\,\bigl(g(x^{\ast}-w_4)-(x^{\ast}-w_4)\bigr)$ lies in $\wExp{\extL{S}}{i}{o^{+}}$; as no~$\cE_\nu$
divides~$p$, Input~\ref{r:exp} gives $g(x^{\ast}-w_4)-(x^{\ast}-w_4)\in\wExp{\extL{S}}{i}{o^{+}}$. Therefore
$x^{\ast}-w_4\in\mathsf{P}_i(o^{+})$ and \m{x^{\ast}\in N'}.
\end{proof}

\begin{lemma}[the colength of the constants that lift]\label{l:col}
In Setting~\ref{set:few}, assume \textup{(F1)} to~\textup{(F6)}. For a letter~$i$ let $o_i$ be the following
exponent function: $o_i=\raisedE{2\rho+\eta}{\cdot}{\cdot}$ if \m{\mathsf{s}(i)=1}; and if
\m{\mathsf{s}(i)=-1},
\[o_i(0,u)=\beta(i)+W+2\rho,\qquad o_i(1,u)=W+1-u+2\rho_u,\qquad o_i(\nu,u)=-u\ \ (\nu\ge2),\]
where \m{\rho_u=\rho} if \m{|u|\le\eta} and \m{\rho_u=0} otherwise. Then for every letter~$i$, with $N_i$ the
image of the map $c\mapsto\constAt{S}{i}{c}$ on~$C_S$,
\[\relLen{\Zp}{N_i}{\deepC{S}{i}{0}{o_i}}\ \le\ (2\rho+1)\,n_{\mathrm{lo}}(i)+3\,(W+1)\,n_{\mathrm{all}}(i^-)
+6\rho\,n_{\mathrm{lo}}(i^-).\]
\end{lemma}
\leanline{\lean{existsrelLengthdeepConstantsle}}
\begin{proof}
By (F2), (F3) and~(F5), \m{\rho\ge1}, every jump of~$Y_S$ is at most \m{(p-3)\rho}, and a jump~$u$ at~$i$ with
\m{|u|>\eta} has $|u|\ge\beta(i)>6(\rho+\eta+W)$.

\emph{A supply of points.} Let \m{s\in\{0,\rho\}}, \m{d\in C_S} and $n$ an integer with
\begin{equation}\label{e:supply}
n-u\le\impE{\raisedE{s-\rho}{\cdot}{\cdot}}{\raisedE{2\rho+\eta}{\cdot}{\cdot}}{0}{u}\qquad\text{for all }u\in U_{Y_S}(i).
\end{equation}
Then $p^{s}\constAt{S}{i}{d}\in\deepC{S}{i}{n}{\raisedE{2\rho+\eta}{\cdot}{\cdot}}$. Indeed
$\Kills{S}{i}{s}{\rho-s}$ by Input~\ref{r:ideal}(h), which gives a point over \m{p^{s}d} with exponents and
depth $\raisedE{s-\rho}{\cdot}{\cdot}$ (Input~\ref{r:points}(a)); the steps to the depth
$\raisedE{2\rho+\eta}{\cdot}{\cdot}$ are at \m{\nu=0} and lie in
$\twoWin{6(\rho+\eta+W)}{\beta(i)}\cup\twoWin{6(\rho+\eta+W)}{-\beta(i)}$, so by (F4) and
Input~\ref{r:points}(c), with \m{\fd=0}, there is a point over \m{p^{s}d} with exponents
$\impE{\raisedE{s-\rho}{\cdot}{\cdot}}{\raisedE{2\rho+\eta}{\cdot}{\cdot}}{\cdot}{\cdot}$ and depth
$\raisedE{2\rho+\eta}{\cdot}{\cdot}$; at \m{\nu\ge1} these exponents are~$-u$, and \eqref{e:supply} is the
comparison at \m{\nu=0}.

\emph{The low weights.} We show, for every integer~$t$ with \m{-\eta\le t\le0},
\begin{equation}\label{e:low}
\relLen{\Zp}{N_i}{\deepC{S}{i}{t}{\raisedE{2\rho+\eta}{\cdot}{\cdot}}}\ \le\ \rho\sum_{-\eta\le u<t}h\,m_{Y_S}(i,u).
\end{equation}
For \m{t=-\eta} the supply with \m{s=0} and \m{n=-\eta} gives
$N_i\subseteq\deepC{S}{i}{-\eta}{\raisedE{2\rho+\eta}{\cdot}{\cdot}}$: in~\eqref{e:supply} the right side is
\m{2\rho+\eta-u} for \m{u<-\rho}, and $\min(2\rho+\eta-u,0)$ for \m{u\ge-\rho}, when \m{u\ge-\eta} because
no jump lies in \m{[-\rho,-\eta)}. Let \m{-\eta\le t<0} and suppose~\eqref{e:low} for~$t$. By
Input~\ref{k:len}(a) and~(b) it suffices to bound by $\rho\cdot h\,m_{Y_S}(i,t)$ the relative length of
$N=\locKP{\wExp{\extL{S}}{i}{\raisedE{t}{\cdot}{\cdot}}}\cap\deepE{S}{i}{\raisedE{2\rho+\eta}{\cdot}{\cdot}}$
and the preimage~$N'$ of $\deepC{S}{i}{t+1}{\raisedE{2\rho+\eta}{\cdot}{\cdot}}$ under \m{\projCloc{S}}; we
use Input~\ref{k:len}(d). For \m{x\in N} over \m{\constAt{S}{i}{d}}, the supply with \m{s=\rho} and
\m{n=t+1\le0} gives \m{p^{\rho}x\in N'}: the right side of~\eqref{e:supply} is now \m{2\rho+\eta-u} for
\m{u<0} and $\min(2\rho+\eta-u,0)\ge t+1-u$ for \m{u\ge0}. Let \m{x_0,\dots,x_m\in N} with \m{m=h\,m_{Y_S}(i,t)}. The
multiplicity of the jump \m{t\ne0} of \m{\extL{S}} is at most that of~$Y_S$ (Input~\ref{k:bk}), and
$\raisedE{t}{0}{t}=0<\raisedE{t}{0}{v}$ for \m{v<t}; so by Input~\ref{r:auer}, for the lattice datum
\m{\extL{S}} and \m{\nu_0=0}, a combination $x^{\ast}=\sum_kc_kx_k$ with \m{c_k\in\Zp}, one of them~$1$,
lies in $\wExp{\extL{S}}{i}{o^{+}}$, where $o^{+}$ is $\raisedE{t}{\cdot}{\cdot}$ raised by~$1$ at
\m{(0,t)} (the other module of Input~\ref{r:auer} is contained in this one). By Input~\ref{r:deep},
$x^{\ast}\in\wExp{\extL{S}}{i}{\impE{o^{+}}{\raisedE{2\rho+\eta}{\cdot}{\cdot}}{\cdot}{\cdot}}$, and
$\impE{o^{+}}{\raisedE{2\rho+\eta}{\cdot}{\cdot}}{0}{u}\ge t+1-u$: for \m{u\le t} the left side is
\m{2\rho+\eta-u}, and for \m{u>t} it is at least $\min(2\rho+\eta-u,0)\ge t+1-u$. So
$x^{\ast}\in\locKP{\wExp{\extL{S}}{i}{\raisedE{t+1}{\cdot}{\cdot}}}$ and \m{x^{\ast}\in N'}.
This proves~\eqref{e:low}; for \m{t=0} its right side is at most $\rho\,n_{\mathrm{lo}}(i)$ by~(F6).

If \m{\mathsf{s}(i)=1} the lemma is proved. Let \m{\mathsf{s}(i)=-1}, and let $\mathsf{N}$ be the least multiple
of~$p$ with $\mathsf{N}\ge\beta(i)+W+2\rho$.

\emph{The index~$0$.} Apply Lemma~\ref{l:rel} with \m{\nu_0=0}, \m{i'=i}, \m{A=0}, \m{B=\rho+1},
$o=\raisedE{2\rho+\eta}{\cdot}{\cdot}$ and \m{o'(0,u)=\mathsf{N}}, \m{o'(\nu,u)=-u} for \m{\nu\ge1}.
Hypotheses (i) to~(iii) are clear. For (iv) let \m{u\in U_{Y_S}(i)}. If \m{|u|\le\eta}, the sum is at most
$\sum_{0\le n\le\mathsf{N}}v_p(n)=v_p(\mathsf{N}!)$, and by Legendre's formula and~(F2)
\[(p-1)\,v_p(\mathsf{N}!)\le\mathsf{N}-1\le\beta(i)+W+2\rho+p-2\le(p-1)(\rho+1).\]
If \m{|u|>\eta}, then \m{-u\ge\beta(i)}, and the sum runs over
$\beta(i)+2\rho+\eta\le n<\mathsf{N}$. A multiple~$n$ of~$p$ in this range has
$n\le\mathsf{N}-p<\beta(i)+W+2\rho$, so it lies in $\twoWin{6(\rho+\eta+W)}{\beta(i)}$ and is not zero,
against~(F4); so the sum is~$0$. Hence
$\relLen{\Zp}{\deepC{S}{i}{0}{o}}{\deepC{S}{i}{0}{o'}}\le(\rho+1)\,n_{\mathrm{lo}}(i)$.

\emph{The index~$1$.} Apply Lemma~\ref{l:rel} with \m{\nu_0=1}, \m{i'=i^-}, \m{A=3(W+1)}, \m{B=6\rho}, the
function~$o'$ just defined in the place of~$o$, and $o_i$ in the place of~$o'$. (i): there is no multiple
of~$p$ strictly between $\mathsf{N}$ and \m{\mathsf{N}+p}. (ii): \m{-u\le\mathsf{N}} for \m{u\ge-\beta(i)-W},
and \m{o_i\le o'} at \m{\nu\ne1}. (iii): for \m{u<u'} in \m{U_{Y_S}(i^-)} one has
$o_i(1,u)\ge o_i(1,u')$ unless \m{|u'|\le\eta<|u|}, when $u'-u\ge\beta(i^-)-\eta>2\rho$. (iv): the sum has
\m{W+1+2\rho_u} terms \m{\cost{1}{n}} with $-u\le n<W+1-u+2\rho_u$; these~$n$ lie in
$\twoWin{6(\rho+\eta+W)}{\beta(i^-)}\cup\twoWin{6(\rho+\eta+W)}{-\beta(i^-)}$, so by (F4) and
Input~\ref{r:exp} each term is at most~$3$. Hence
$\relLen{\Zp}{\deepC{S}{i}{0}{o'}}{\deepC{S}{i}{0}{o_i}}\le3(W+1)\,n_{\mathrm{all}}(i^-)+6\rho\,n_{\mathrm{lo}}(i^-)$.

The three bounds add up, by Input~\ref{k:len}(a), to the bound of the lemma.
\end{proof}

\begin{proposition}[few low weights: a submodule of small colength on which cup products bound]\label{p:few}
In Setting~\ref{set:few}, assume \textup{(F1)} to~\textup{(F6)}, and suppose \m{\eta\ge1}. Then there is a
$\Zp$-submodule~$\mathsf{C}$ of~$C_S$ with
\[\relLen{\Zp}{C_S}{\mathsf{C}}\ \le\ 11\sum_{i}\bigl(n_{\mathrm{lo}}(i)\,(\rho+1)+n_{\mathrm{all}}(i)\,(W+1)\bigr),\]
the sum over the letters, such that for every side~$T$ over~$\cO$, every pairing~$\Pi$ of $T$ and~$S$ and
every \m{c\in C_T}: if $\JumpWin{\eta}{-\mathsf{s}(i)}{\beta(i)}{W}{U_{Y_T}(i)}$ for every letter~$i$, then
$\CupB{\Pi}{0}{c}{d}$ for every \m{d\in \mathsf{C}}.
\end{proposition}
\leanline{\lean{existsscaledCupBoundsfewlowjumps}}
\begin{proof}
Let $\mathsf{C}$ be the set of the \m{d\in C_S} with $\constAt{S}{i}{d}\in\deepC{S}{i}{0}{o_i}$ for every letter~$i$,
with $o_i$ as in Lemma~\ref{l:col}. By Input~\ref{k:len}(b) and~(c) and Lemma~\ref{l:col},
\begin{align*}\relLen{\Zp}{C_S}{\mathsf{C}}&\le\sum_i\bigl((2\rho+1)\,n_{\mathrm{lo}}(i)+3(W+1)\,n_{\mathrm{all}}(i^-)+6\rho\,n_{\mathrm{lo}}(i^-)\bigr)\\
&=\sum_i\bigl((8\rho+1)\,n_{\mathrm{lo}}(i)+3(W+1)\,n_{\mathrm{all}}(i)\bigr),\end{align*}
since \m{i\mapsto i^-} permutes the letters; this is at most the bound stated.

Let $T$, $\Pi$, $c$ be as in the statement and \m{d\in \mathsf{C}}. We apply Theorem~B of Part~2
(Input~\ref{r:walk}) to the pairing~$\Pi$ of $T$ and~$S$ with \m{\kappa=e=e'=E=0}, with every letter a cut, so
that \m{\mathsf{l}(i)=1}, with \m{\mathsf{Q}(i)} the statement \m{\mathsf{s}(i)=1}, and with the families
\[\begin{gathered}o^{T}=\impE{\raisedE{-\rho}{\cdot}{\cdot}}{\raisedE{2\rho+\eta}{\cdot}{\cdot}}{\cdot}{\cdot},\qquad
o^{S}_i=\impE{\raisedE{0}{\cdot}{\cdot}}{o_i}{\cdot}{\cdot},\\
\mathsf{z}=\impE{\raisedE{-\rho}{\cdot}{\cdot}}{\raisedE{\rho}{\cdot}{\cdot}}{\cdot}{\cdot},\qquad
\mathsf{z}'=\raisedE{\rho}{\cdot}{\cdot},\end{gathered}\]
the first, third and fourth being the same at all letters. Explicitly, at \m{\nu\ge1} the functions $o^{T}$
and~$\mathsf{z}$ are \m{u\mapsto-u}, and
\[o^{T}(0,u)=\begin{cases}0,&-\rho\le u\le2\rho+\eta,\\2\rho+\eta-u,&\text{otherwise,}\end{cases}\qquad
\mathsf{z}(0,s)=\begin{cases}0,&|s|\le\rho,\\ \rho-s,&|s|>\rho;\end{cases}\]
in particular $\mathsf{z}(\nu,\cdot)$ is non-increasing and
\begin{equation}\label{e:z}
\mathsf{z}(0,s)\ge-s-\rho\quad\text{for all }s .
\end{equation}
For the jumps~$v$ of~$Y_S$ at~$j$, \m{o^{S}_j(0,v)} is \m{2\rho+\eta-v} or~$0$ if \m{\mathsf{s}(j)=1}, the value~$0$
only for \m{v\ge0}; it is \m{\beta(j)+W+2\rho} for \m{v<0} and $0$ for \m{v\ge0} if \m{\mathsf{s}(j)=-1}; and
\begin{equation}\label{e:oS}
o^{S}_j(\nu,v)\ge-v\ \text{ and }\ o_j(\nu,v)\ge-v\qquad\text{for all }\nu\ge0\text{ and }v\in U_{Y_S}(j^{p^{-\nu}}),
\end{equation}
because \m{v\ge-\beta(j)-W} at \m{\nu=0}.

\emph{The points.} By (F2) every jump of $T$, of~$S$ and of the third side is at most \m{(p-3)\rho}: those of
the third side are $0$ or \m{-u-v-1} with $u$, $v$ jumps of $T$ and~$S$ at the same letter, in opposite
windows. So $\Kills{T}{i}{0}{\rho}$ and $\Kills{\third{\Pi}}{i}{0}{\rho}$ (Input~\ref{r:ideal}(h)), and as
in the proof of Lemma~\ref{l:col}, by (F4) and Input~\ref{r:points}(a) and~(c) with \m{\fd=0}, the sets
$\ptsOver{T}{c}{i}{0}{o^{T}}{\raisedE{2\rho+\eta}{\cdot}{\cdot}}$ and
$\ptsOver{\third{\Pi}}{1}{i}{0}{\mathsf{z}}{\mathsf{z}'}$ are not empty: the steps of the third side are within \m{\rho+2\eta+W+1} of $0$ or~\m{\pm\beta(i)}, hence in the
windows of~(F4), because \m{\rho\ge1} by (F2) and~(F3). (The formal statement assumes \m{\eta\ge1} and
passes it to its automated arithmetic; no step of the argument given here uses it.) Since \m{d\in \mathsf{C}}, there is $x\in\mathsf{P}_i(o_i)$ over \m{\constAt{S}{i}{d}}; by
Input~\ref{r:deep} it is a point of $\ptsOver{S}{d}{i}{0}{o^{S}_i}{o_i}$.

\emph{The deep point of~$T$.} We check Definition~2.6 of Part~2 for $T$ against~$S$ at~$i$, with the depth
$\raisedE{2\rho+\eta}{\cdot}{\cdot}$ and \m{r=0}. Let \m{\nu\ge0}, $u$ a jump of~$T$ and $v$ a jump of~$S$
at \m{i^{p^{-\nu}}}, and \m{s=-u-v-1}. The left side of clause~(a) is \m{s+u=-v-1} for \m{\nu\ge1}, and for
\m{\nu=0} it is, by~\eqref{e:z},
\[-\mathsf{z}(0,s)-(2\rho+\eta-u)\ \le\ -v-1-\rho-\eta .\]
Both are smaller than~$-v$ and, by \eqref{e:oS} and the list of the values of \m{o^{S}_i(0,v)}, smaller than
\m{o^{S}_i(\nu,v)}. For clause~(b) let \m{\mathsf{s}(i)=1}. Then the jumps of~$T$ at~$i$ are at most~$\eta$, so
$o^{T}(0,u)\ge0=\std_1(0,u)$, and $o^{T}(\nu,u)=-u=\std_1(\nu,u)$ for \m{\nu\ge1}. With \m{r=0} the left
side is that of~(a); the additional bound $\frobX{1}{0}{o^{S}_{i^-}}{\nu}{v}$ is $0$ for \m{\nu=0}, where
the left side is at most \m{-v-1-\rho-\eta<0} because \m{v\ge-\eta}; and it is
$o^{S}_{i^-}(\nu-1,v)\ge-v$ for \m{\nu\ge1}, by~\eqref{e:oS}.

\emph{The deep point of~$S$.} We check the definition for $S$ against~$T$ at~$i$, for the proposition
\m{\mathsf{s}(i)=-1}, with the depth~$o_i$ and \m{r=1}. Let $v$ be a jump of~$S$ and $u$ a jump of~$T$ at
\m{i^{p^{-\nu}}}, \m{s=-u-v-1}. Clause~(a) asks that
$-\mathsf{z}(\nu,s)-o_i(\nu,v)<\min(-u,o^{T}(\nu,u))$. For \m{\nu\ge1} the left side is at most
\m{s+v=-u-1}, by~\eqref{e:oS}, and \m{o^{T}(\nu,u)=-u}. For \m{\nu=0} and \m{\mathsf{s}(i)=1} it is at most
\m{-u-1-\rho-\eta} by~\eqref{e:z}. For \m{\nu=0} and \m{\mathsf{s}(i)=-1} it is
$-\mathsf{z}(0,s)-\beta(i)-W-2\rho\le-u-1-\rho$, by~\eqref{e:z} and \m{v\ge-\beta(i)-W}. In both cases it
is smaller than~$-u$, than \m{2\rho+\eta-u}, and than~$0$ when \m{u\ge-\rho}, hence smaller than
\m{o^{T}(0,u)}.

Clause~(b), for \m{\mathsf{s}(i)=-1}: \m{o^{S}_i(0,v)\ge0} and $o^{S}_i(\nu,v)\ge-v$ for \m{\nu\ge1}, so
\m{\std_1\le o^{S}_i} on the jumps. It remains to show
\[-\frobX{1}{0}{\mathsf{z}}{\nu}{s}-o_i(\nu,v)\ <\ \min\bigl(o^{T}(\nu,u),\ \frobX{1}{0}{o^{T}}{\nu}{u}\bigr).\]
For \m{\nu=0} the left side is $-(\beta(i)+W+2\rho)$, the right side is \m{\min(o^{T}(0,u),0)}, and
$o^{T}(0,u)\ge2\rho+\eta-\beta(i)-W$ or \m{o^{T}(0,u)=0}. For \m{\nu\ge2} the left side is \m{s+v=-u-1} and
the right side is~$-u$. Let \m{\nu=1}; the letter is~$i^-$, the left side is
$-\mathsf{z}(0,s)-(W+1-v+2\rho_v)$, the right side is $\min(-u,o^{T}(0,u))$. So we need
\[\mathsf{z}(0,s)>-s-W-2-2\rho_v\qquad\text{and}\qquad\mathsf{z}(0,s)>v-W-1-2\rho_v-o^{T}(0,u).\]
The first: if \m{|v|\le\eta} it follows from~\eqref{e:z}; if \m{|v|>\eta} and \m{|s|>\rho} from
$\mathsf{z}(0,s)=\rho-s$; and if \m{|v|>\eta} and \m{|s|\le\rho}, then $u$ is a far jump as well, in the
opposite window, so \m{|u+v|\le W}, \m{s\ge-W-1} and $\mathsf{z}(0,s)=0>-s-W-2$. The second: if
$o^{T}(0,u)=2\rho+\eta-u$ it follows from~\eqref{e:z}. If \m{o^{T}(0,u)=0}, then $-\rho\le u\le2\rho+\eta$,
so \m{|u|\le\eta}; for \m{|v|\le\eta}, \eqref{e:z} gives $\mathsf{z}(0,s)\ge u+v+1-\rho\ge v+1-2\rho$; for
\m{|v|>\eta} one has \m{|s|>\rho} and $\mathsf{z}(0,s)=\rho+u+v+1\ge v+1$.

All hypotheses of Theorem~B of Part~2 hold, and it gives $\CupB{\Pi}{0}{c}{d}$.
\end{proof}

\sect{Tame charts}{s:tame}
In this section $\fc$ is a tame chart for~$\Gamma$ (Definition~\ref{d:tame}), with its root of
unity~$\varrho$, its tame generator~$\sigma$ and its element~$\varpi$. For a polynomial~$q$ over~$\Zp$ write
\m{q=\sum_kq_kX^{k}}; it kills \m{(\sigma,0)} on a lattice if its image in~\m{\cO[X]} does.

\begin{lemma}[a polynomial in the tame generator makes elements divisible by the coordinate]\label{l:zsum}
In Setting~\ref{set:base}, let $\fc$ be a tame chart for~$\Gamma$ with tame generator~$\sigma$, $L$ a lattice
datum over~$\cO$, and $q$ a polynomial over~$\Zp$ with $q(\linF{L}{\sigma}{0})=0$. Then for every
$x\in\mathfrak{M}(L)$ there is \m{y\in\Minf(L)} with
\[z\cdot y=\sum_{k=0}^{\deg q}q_k\,\sigma^{k}x .\]
\end{lemma}
\leanline{\lean{existszsmuleqsumtameGen}}
\begin{proof}
Since $\Gamma'$ is normal, each \m{\sigma^{k}} preserves \m{\Minf(L)}, and $\mu_k=\sigma^{-k}\circ\mu\circ\sigma^{k}$
is a period for every period~$\mu$ (Input~\ref{k:frame}(e)); so the right side~$w$ lies in \m{\Minf(L)}. By
Input~\ref{k:frame}(a) it suffices to show $\cE^{2h_L}\mu(w)\in B^{+}z$ for every crystalline
period~$\mu$. Let $\lambda\colon L\to B_{\max}$ give the values of~$\mu$ on \m{1\otimes L}, and for
\m{g\in\Gamma} let $\Lambda_g\colon B_{\max}\otimes L\to B_{\max}$ be the $B_{\max}$-linear map with
$\Lambda_g(1\otimes t)=j_\Gamma(g)^{-1}\lambda(gt)$. Then
\begin{equation}\label{e:Lam}
\Lambda_{g'}\bigl((j_\Gamma(g)\otimes g)\,y\bigr)=j_\Gamma(g)\,\Lambda_{g'g}(y)\qquad(g,g'\in\Gamma,\ y\in B_{\max}\otimes L),
\end{equation}
and $\Lambda_{\sigma^{k}g}=\Lambda_{\sigma^{k}}$ for \m{g\in\Gamma'}, because $\mu$ commutes with the element
$\sigma^{k}g\sigma^{-k}$ of~$\Gamma'$.

Let $y\in B_{\max}\otimes L$ be fixed near~$1$. By~\eqref{e:Lam} the element $\Lambda_{\sigma^{k}}(y)$
of~$B_{\max}$ is fixed near~$1$, hence fixed by the inertial element~$\sigma^{k}$ (Input~\ref{k:frob}(c)). On
such~$y$, $(\linF{L}{\sigma}{0})^{k}=\linF{L}{\sigma^{k}}{0}$ acts as $j_\Gamma(\sigma^{k})\otimes\sigma^{k}$
(Input~\ref{k:frob}(a)), so by~\eqref{e:Lam} again
\[\begin{gathered}\Lambda_1\bigl((\linF{L}{\sigma}{0})^{k}y\bigr)=j_\Gamma(\sigma^{k})\,\Lambda_{\sigma^{k}}(y)=\Lambda_{\sigma^{k}}(y),\\
\text{and}\qquad\sum_kq_k\Lambda_{\sigma^{k}}(y)=\Lambda_1\bigl(q(\linF{L}{\sigma}{0})\,y\bigr)=0 .\end{gathered}\]
As $L$ is potentially crystalline, such~$y$ span \m{B_{\max}\otimes L}; so $\sum_kq_k\Lambda_{\sigma^{k}}=0$.
On \m{1\otimes t} this says $\sum_kq_k\mu_k(1\otimes t)=0$, hence $\sum_kq_k\mu_k(x)=0$
(Input~\ref{k:frame}(f)). Therefore
\[\mu(w)=\sum_kq_k\,\sigma^{k}\bigl(\mu_k(x)\bigr)=\sum_kq_k\bigl(\sigma^{k}\mu_k(x)-\mu_k(x)\bigr).\]
Each $\mu_k$ is a crystalline period, so \m{\cE^{h_L}\mu_k(x)} lies in the image of~$\fS_B$
(Input~\ref{k:frame}(b)); \m{\sigma^{k}} fixes~$\WW$, so $z$ divides \m{\sigma^{k}s-s} for \m{s\in\fS_B}
(Input~\ref{k:frame}(c)); and Input~\ref{k:frame}(d) gives
$\cE^{2h_L}(\sigma^{k}\mu_k(x)-\mu_k(x))\in B^{+}z$.
\end{proof}

\begin{lemma}[the extension is killed by the product with \m{X-1}]\label{l:extkill}
In Setting~\ref{set:base}, let $\fc$ be a tame chart for~$\Gamma$ with tame generator~$\sigma$, $S$ a side
over~$\cO$, and $q$ a polynomial over~$\Zp$ that kills \m{(\sigma,0)} on~$Y_S$. Then
\[\bigl(q\cdot(X-1)\bigr)\bigl(\linF{\extL{S}}{\sigma}{0}\bigr)=0 .\]
\end{lemma}
\leanline{\lean{aevallinearFrobextLatticeeqzero}}
\begin{proof}
The maps \m{\inclY{S}} and \m{\projC{S}} commute with~$\Gamma$, so the linear Frobenius commutes with
\m{1\otimes\inclY{S}} and \m{1\otimes\projC{S}}, and it is the identity on $B_{\max}\otimes\trivL{S}$
(Input~\ref{k:frob}(b)). For $y\in B_{\max}\otimes\extL{S}$ the element $\linF{\extL{S}}{\sigma}{0}y-y$
therefore maps to~$0$ in $B_{\max}\otimes C_S$, so it is $(1\otimes\inclY{S})(x)$ with
$x\in B_{\max}\otimes Y_S$, and
$q(\linF{\extL{S}}{\sigma}{0})(1\otimes\inclY{S})(x)=(1\otimes\inclY{S})\bigl(q(\linF{Y_S}{\sigma}{0})x\bigr)=0$.
\end{proof}

\begin{proposition}[tame charts: a cokernel killer independent of the index]\label{p:tamekill}
In Setting~\ref{set:base}, let $\fc$ be a tame chart for~$\Gamma$ with tame generator~$\sigma$, $S$ a side
over~$\cO$, \m{n,m,\kappa\ge0} integers, $i$ a letter, and $q$ a monic polynomial over~$\Zp$ that kills
\m{(\sigma,0)} on~$Y_S$. Suppose that \m{u\le n} for every letter~$j$ and every \m{u\in U_{Y_S}(j)}, that
\[e_{\fc}\le p^{m},\qquad n+\kappa+2<p,\qquad m\,(\deg q+1)\le\kappa .\]
Then $\Kills{S}{i}{\kappa}{0}$.
\end{proposition}
\leanline{\lean{killsCokernelofisKilledBy}}
\begin{proof}
Put $Q=q\cdot(X-1)=\sum_kQ_kX^{k}$, monic of degree \m{\deg q+1}, with \m{Q(1)=0}, and \m{r=p^{\kappa+1}}.
By Input~\ref{r:ideal}(g) there is \m{a_0\ge0} with $z^{a_0}\in\liftId{S}{r}{0}{i}$ and
$a_0(p-1)\le e_{\fc}(n+\kappa+1)$; as \m{n+\kappa+1<p-1}, \m{a_0<e_{\fc}}. For \m{0\le a<e_{\fc}} let
\m{N_a} be the number of the integers~$k$ with \m{a<k<e_{\fc}} and $Q(\varrho^{k})\in p\WW$. We show, by
descending induction on~$a$ from~$a_0$,
\begin{equation}\label{e:down}
p^{mN_a}z^{a}\in\liftId{S}{r}{0}{i}\qquad(0\le a\le a_0).
\end{equation}
For \m{a=a_0} it holds. Let \m{0<a\le a_0} and suppose~\eqref{e:down} for~$a$; write \m{c=mN_a}.

\emph{If $Q(\varrho^{a})\in p\WW$}, then \m{N_{a-1}=N_a+1}. By Input~\ref{r:ideal}(c),
$p^{c}z^{pa}\in\liftId{S}{r}{1}{i^{p}}$, and by Input~\ref{r:ideal}(e), with \m{p(a-1)\le pa}, the element
\m{p^{c+m}z^{a-1}} lies in $\liftId{S}{r}{0}{i}$ as soon as $c+v_p((pa)!)\le c+m+v_p((p(a-1))!)$, that is
\m{1+v_p(a)\le m}; and $p^{v_p(a)}\le a<e_{\fc}\le p^{m}$.

\emph{If \m{Q(\varrho^{a})} is a unit of~$\WW$}, then \m{N_{a-1}=N_a}. Let \m{f=p^{c}z^{a}}. By
Input~\ref{r:ideal}(d) there is \m{t_0\in\fS} such that for every \m{c'\in C_S} some
$x\in\locLP{0}{i}(\extL{S})$ lies over $(f-r\,t_0)\constAt{S}{i}{c'}$; write \m{t_0=w+z\,t} with
\m{w\in\WW}, \m{t\in\fS} (Input~\ref{k:chart}(c)). By Input~\ref{r:exp}, $x\in\Minf(\extL{S})$, and it is
a point over~$F_0$; so $x\in\mathfrak{M}(\extL{S})$. By Lemmas \ref{l:extkill} and~\ref{l:zsum} there is
$y\in\Minf(\extL{S})$ with $z\,y=\sum_kQ_k\sigma^{k}x$. Now \m{\sigma^{k}} fixes~$\WW$, hence commutes
with~$\pr_i$ and fixes \m{\constAt{S}{i}{c'}}, and $\sigma^{k}z=\varrho^{k}z$; so, with
$T=\sum_kQ_k\varrho^{k}\sigma^{k}(t)\in\Ainf$,
\begin{align*}z\cdot(1\otimes\projC{S})(\pr_iy)&=\sum_kQ_k\bigl(p^{c}\varrho^{ka}z^{a}-r\,w-r\,\varrho^{k}z\,\sigma^{k}(t)\bigr)\constAt{S}{i}{c'}\\
&=z\,\bigl(Q(\varrho^{a})\,p^{c}z^{a-1}-r\,T\bigr)\constAt{S}{i}{c'},\end{align*}
the term with~$w$ vanishing because \m{Q(1)=0}. The coordinate~$z$ is a unit of~$\WC$, and
$\pr_iy\in\wExp{\extL{S}}{i}{\std_0}$ (Input~\ref{r:exp}); so by Input~\ref{r:ideal}(a),
$Q(\varrho^{a})\,p^{c}z^{a-1}\in\liftId{S}{r}{0}{i}$, and the unit can be removed.

This proves~\eqref{e:down}. For \m{a=0} it says $p^{mN_0}\in\liftId{S}{r}{0}{i}$. The elements
\m{\varrho^{k}}, \m{0<k<e_{\fc}}, are distinct modulo~$p$, and the reduction of~$Q$ is monic, so
\m{N_0\le\deg q+1} and \m{mN_0\le\kappa}. Hence $p^{\kappa}\in\liftId{S}{p^{\kappa+1}}{0}{i}$, and
Input~\ref{r:ideal}(b) gives $p^{\kappa}\in\liftId{S}{0}{0}{i}$.
\end{proof}

\begin{proposition}[existence of tame charts]\label{p:tameexist}
In Setting~\ref{set:base}, suppose that \m{j_\Gamma} is surjective. Let \m{e\ge0} and \m{h_0\ge1} be integers
such that $p$ does not divide~$h_0$ and $e$ divides \m{p^{h_0}-1}. Then there is a tame chart for~$\Gamma$
whose index is~$e$.
\end{proposition}
\leanline{\lean{existstameChart}}
\begin{proof}
Put \m{N=p^{h_0}-1}; it is prime to~$p$, and \m{e\ge1}. Choose $\varpi\in\bar{\mathbb{Q}}_p$ with
\m{\varpi^{e}=p}, an element \m{\varpi^{\flat}} of the tilt~$\mathcal{O}_C^{\flat}$ given by a compatible system
of $p$-power roots of~$\varpi$ (its \kw{untilt} is~$\varpi$), and a primitive $e$-th root of unity~$a$ in the
subfield $\mathbb{F}_{p^{h_0}}$ of the tilt. By Input~\ref{k:chart}(a), for \m{x=\varpi} and \m{v=-1}, there is a chart with residue degree~$h_0$,
index~$e$, constant unit series, and coordinate \m{z=[\varpi^{\flat}]} with \m{\theta(z)=\varpi}.

\emph{The subgroup.} Let \m{\chi_N} be the character of~$\GQp$ on the $N$-th roots of unity, and
$\Delta=\ker\chi_N\cap\operatorname{Stab}(\varpi)$, an open subgroup. An element of~$\GQp$ fixes~$\WW$ if
and only if it lies in \m{\ker\chi_N}; so $\Delta$ is the group of the elements that fix $\WW$
and~\m{\theta(z)}. For \m{g\in\GQp} the quotient \m{g\varpi/\varpi} is an $e$-th, hence an $N$-th root of
unity. It follows that $\Delta$ is normal: for \m{\delta\in\Delta} and \m{g\in\GQp}, $\delta$ fixes
$g^{-1}\varpi=(g^{-1}\varpi/\varpi)\,\varpi$, so \m{g\delta g^{-1}} fixes~$\varpi$. It also follows that
\m{g^{N}\in\Delta} for \m{g\in\ker\chi_N}, because $g$ fixes the root of unity \m{g\varpi/\varpi}. Finally
$g^{h_0}\in\ker\chi_N$ for every~$g$: it acts on the $N$-th roots of unity as a power of the $p^{h_0}$-th
power map, which is the identity. So every element of \m{\GQp/\Delta} has order dividing \m{h_0N}, the order
of this finite group divides a power of \m{h_0N}, and it is prime to~$p$. Put $\Gamma'=j_\Gamma^{-1}(\Delta)$;
it is normal, of index prime to~$p$ since $j_\Gamma$ is surjective, it contains the elements near~$1$,
$j_\Gamma(\Gamma')=\Delta$, and its elements fix~$\varpi$.

\emph{The elements of inertia.} The inertia group is closed and acts transitively on the $n$-th roots
of~$p$ for every \m{n>0}; so by Input~\ref{k:chart}(b) and the surjectivity of~$j_\Gamma$, for every
$v\in\mathcal{O}_C^{\flat}$ whose untilt has $e$-th power~$p$ there is an inertial \m{s\in\Gamma} with
\m{s\cdot z=[v]}. Inertial elements fix~$\WW$. For $v=\varepsilon^{l}\varpi^{\flat}$, \m{l\ge0}, the
element~$s$ fixes \m{\theta(z)}, so it lies in~$\Gamma'$ and $s\cdot z=[\varepsilon]^{l}z$: with~$\Gamma'$
the chart is a Kummer chart. For $v=a\,\varpi^{\flat}$ we get the tame generator~$\sigma$, with
\m{\sigma z=[a]z} and, since $\sigma$ fixes \m{\varrho=[a]\in\WW}, $\sigma^{k}z=\varrho^{k}z$ for all~$k$.
The index $e$ divides~$N$, so $p$ does not divide it.
\end{proof}

\sect{Cuts and windows}{s:cuts}
In this section the setting is Setting~\ref{set:cut}.

\begin{lemma}[the cokernel killer at a cut]\label{l:atcut}
In Setting~\ref{set:cut}, assume \textup{(W3)}, \textup{(W4)}, \textup{(W5)}, \textup{(W6)}, \textup{(W8)},
\textup{(W9)}, \textup{(W10)} and~\textup{(W11)}. Let $T$ be a side over~$\cO$ such that for every letter~$j$
and every integer~$n'$ with \m{\beta(j)+W\le n'}, every \m{u\in U_{Y_T}(j)} satisfies \m{u\le n'}. Then
$\Kills{T}{i}{\kappa}{\rho(i)}$ for every letter~$i$ with the property~$\mathsf{cut}$.
\end{lemma}
\leanline{\lean{killsCokernelatcut}}
\begin{proof}
Let $i$ be a cut. If the second alternative of (W10) holds, then by (W11) every jump of~$T$ is at most
$(p-3)\rho_0\le(p-3)(\kappa+\rho(i))$, and Input~\ref{r:ideal}(h) gives the claim. Suppose the first:
\m{\rho(i)=0} and $\beta(i^{p^{a}})+W\le G$ for \m{1\le a\le J}. Put \m{r=p^{\kappa+1}}. By (W8) every jump
of~$T$ is at most~$2B$, so by Input~\ref{r:ideal}(g) there is \m{a_0\ge0} with
$z^{a_0}\in\liftId{T}{r}{0}{j}$ for every letter~$j$ and $a_0(p-1)\le e_{\fc}(2B+\kappa+1)$.

\emph{Claim.} Let \m{J'\le J}, let $j$ be a letter with $\beta(j^{p^{a}})+W\le G$ for \m{1\le a\le J'}, and
\m{x\ge0} an integer with $p\,a_0\le p^{J'+1}x$. Then $p^{J'X}z^{x+n}\in\liftId{T}{r}{0}{j}$. For
\m{J'=0}, \m{a_0\le x}. Let \m{J'\ge1}. The claim for \m{J'-1}, the letter~$j^{p}$ and the integer~$px$ gives
$p^{(J'-1)X}z^{px+n}\in\liftId{T}{r}{0}{j^{p}}$; the jumps of~$T$ at~$j^{p}$ are at most~$G$, so
$p^{(J'-1)X}z^{px+n}\cE^{G}\in\liftId{T}{r}{1}{j^{p}}$ (Input~\ref{r:ideal}(c)). By Legendre's formula and
(W9), \m{v_p(G!)\le n}, so $\fd\,G+v_p(G!)\le X$ by~(W3), and \m{px+n+G\le p(x+n)} by~(W9);
Input~\ref{r:ideal}(f) gives $p^{J'X}z^{p(x+n)}\in\liftId{T}{r}{1}{j^{p}}$, and Input~\ref{r:ideal}(c) the
claim.

Take \m{J'=J}, \m{j=i} and $x_0=2X+\lfloor\kappa/4\rfloor+1$. By (W5) and~(W6),
\[p\,a_0\le2\,e_{\fc}\,(2B+\kappa+1)\le p^{J+1}\bigl(2X+\tfrac{\kappa+1}{4}\bigr)\le p^{J+1}x_0 ,\]
so $p^{JX}z^{m}\in\liftId{T}{r}{0}{i}$ with \m{m=x_0+n}, and $p^{JX}z^{pm}\in\liftId{T}{r}{1}{i^{p}}$.
Now $v_p((pm)!)=m+v_p(m!)\le2m-1$, and
\[JX+2m-1=(J+4)X+2n+1+2\lfloor\kappa/4\rfloor\le\kappa ,\]
because \m{n\le X} by (W3) and $\kappa-2\lfloor\kappa/4\rfloor\ge6(J+1)X+1$ by~(W4). So
Input~\ref{r:ideal}(e), with the exponent~$0$ of~$z$, gives $p^{\kappa}\in\liftId{T}{p^{\kappa+1}}{0}{i}$,
and Input~\ref{r:ideal}(b) gives $p^{\kappa}\in\liftId{T}{0}{0}{i}$.
\end{proof}

\begin{proposition}[cuts with windows]\label{p:cuts}
In Setting~\ref{set:cut}, assume \textup{(W1)} to~\textup{(W11)}, and let \m{E\ge0} be an integer. Suppose
that:
\begin{enumerate}[label=\textup{(\roman*)},leftmargin=2.3em,itemsep=2pt]
\item \m{\rho(i)+W\le p\,n} for every letter~$i$;
\item \m{51\,\kappa\le E};
\item for every letter~$i$ there is an integer~$l$ with \m{0<l\le2J+1} such that \m{i^{p^{-l}}} has the
property~$\mathsf{cut}$;
\item for every letter~$i$ and every integer \m{l\ge0} such that \m{i^{p^{-l}}} has the property~$\mathsf{cut}$
and no \m{i^{p^{-a}}} with \m{0<a<l} has it, there is an integer \m{t\ge0} for which one of the following
holds:
\begin{itemize}[leftmargin=2.4em,itemsep=1pt]
\item[\textup{($\alpha$)}] for every \m{0\le a<l}: $\mathsf{s}(i^{p^{-a}})=\mathsf{s}(i^{p^{-t}})$ or
$\beta(i^{p^{-a}})+V\le\eta$;
\item[\textup{($\beta$)}] the chart is unramified, \m{p\ge5}, \m{\rho(i^{p^{-l}})=0}, \m{t+J+1=l}, \m{\rho(i)=0} or \m{t=0},
\[\rho_0+G\sum_{0\le a\le J}p^{a}<\beta(i^{p^{-t}}),\]
and for every \m{0\le a<l} with \m{a\ne t}: $\mathsf{s}(i^{p^{-a}})\ne\mathsf{s}(i^{p^{-t}})$ or
$\beta(i^{p^{-a}})+V\le\eta$; and $\beta(i^{p^{-a}})+W\le G$ if \m{t<a}.
\end{itemize}
\end{enumerate}
Then $\CupB{\Pi}{E}{c}{c'}$ for all \m{c\in C_S} and \m{c'\in C_{S'}}.
\end{proposition}
\leanline{\lean{scaledCupBoundsofcutwindows}}
\begin{proof}
\emph{The data of the walk.} For a letter~$i$ let \m{\mathsf{l}(i)} be the least \m{l>0} such that
\m{i^{p^{-l}}} is a cut; $\mathsf{l}(i)\le2J+1$ by~(iii). Let \m{t(i)} be an integer~$t$ of~(iv) for $i$ and
\m{l=\mathsf{l}(i)}, and $\mathsf{s}^{\circ}(i)=\mathsf{s}(i^{p^{-t(i)}})$. Let \m{\eta_\nu=\eta} for
\m{\nu\le2J} and \m{\eta_\nu=0} for \m{\nu>2J}. For a sign~$\epsilon$ let
\[\mathsf{R}^{\epsilon}_i(\nu,u)=\max\bigl(\std_{\mathsf{l}(i)}(\nu,u),\,-u-\eta_\nu\bigr)\ \text{ if }\mathsf{s}^{\circ}(i)=\epsilon,\qquad
\mathsf{R}^{\epsilon}_i(\nu,u)=-u-\eta_\nu\ \text{ otherwise.}\]
We apply Theorem~B of Part~2 (Input~\ref{r:walk}) with the property~$\mathsf{cut}$, with \m{\mathsf{Q}(i)} the
statement $\mathsf{s}^{\circ}(i)=1$, with the families \m{o=\mathsf{R}^{+1}}, \m{o'=\mathsf{R}^{-1}} and
$z(i)=z'(i)=\raisedE{-\rho(i)}{\cdot}{\cdot}$, and with \m{e=e'=25\kappa}; by~(ii), \m{e+e'+\kappa\le E}. The
functions \m{z(i)(\nu,\cdot)} are non-increasing.

\emph{The third side.} A jump of \m{\third{\Pi}} at~$j$ is $0$ or \m{-u-v-1} with $u$, $v$ jumps of $S$
and~$S'$ at~$j$; by (W1) and~(W2) it is at most \m{\beta(j)+W}. So Lemma~\ref{l:atcut} applies to the third
side, and to $S$ and~$S'$. For a cut~$i$, Input~\ref{r:points}(a) gives a point of
$\ptsOver{\third{\Pi}}{1}{i}{\kappa}{z(i)}{z'(i)}$.

\emph{The deep points.} Let $\epsilon$, $\epsilon'$ be signs and $T$, $T'$ sides over~$\cO$ such that
$\JumpWin{\eta}{-\epsilon\mathsf{s}(j)}{\beta(j)}{V}{U_{Y_T}(j)}$ and
$\JumpWin{\eta}{\epsilon\mathsf{s}(j)}{\beta(j)}{V}{U_{Y_{T'}}(j)}$ for every letter~$j$; let \m{c\in C_T} and
$i$ a cut. We claim that $T$ has a deep point over \m{p^{25\kappa}c} at~$i$ against~$T'$, for the data
$(\mathsf{cut},\mathsf{l}(i),\mathsf{R}^{\epsilon},\mathsf{R}^{\epsilon'},z)$ and the proposition
$\mathsf{s}^{\circ}(i)=\epsilon$. For $(T,T',\epsilon,\epsilon')=(S,S',1,-1)$ and \m{(S',S,-1,1)} these are
hypotheses (i) and~(ii) of Theorem~B of Part~2, by~(W2).

Let $I$ be the set of the pairs \m{(\nu,n)} with \m{\nu<4J+2} and
$n\in\twoWin{W_\nu}{\pm\beta(i^{p^{-\nu}})}$, where \m{W_0=\rho(i)+W}, \m{W_\nu=W} for
\m{\nu>0}, and the sign is~$+$ if and only if $\mathsf{s}(i^{p^{-\nu}})=\epsilon$. By Input~\ref{r:ideal}(i),
with (i) and~(W8), the costs of the pairs of~$I$ with a given~$\nu$ have sum at most
\[(4n+2)\lfloor\log_p(2B)\rfloor+(4W_\nu+2)(\fd+3\nu)\ \le\ 72\,X ,\]
by (W3) and \m{W\le n}; for \m{\nu=0} one uses that \m{\fd=0} if \m{\rho(i)\ne0}, by~(W10). So the costs
of~$I$ have sum at most \m{(4J+2)\cdot72X}, and by~(W4)
\begin{equation}\label{e:25}
\kappa+4X+(4J+2)\cdot72\,X\le25\,\kappa .
\end{equation}
Let \m{t=t(i)}. Two cases.

\emph{Case~1: $\mathsf{s}^{\circ}(i)\ne\epsilon$, or alternative~\textup{($\alpha$)} of~\textup{(iv)} holds.}
By Lemma~\ref{l:atcut} and Input~\ref{r:points}(a) there is a point over \m{p^{\kappa}c} at~$i$ with exponents
and depth $\raisedE{-\rho(i)}{\cdot}{\cdot}$. Let
$d(\nu,u)=\raisedE{\rho(i)}{\nu}{u}+\eta_{\lfloor\nu/2\rfloor}$. A step \m{(\nu,n)} from
$\raisedE{-\rho(i)}{\cdot}{\cdot}$ to~$d$ has \m{\nu\le4J+1} and
$-u-(W_\nu-W)\le n<-u+(W_\nu-W)+\eta$ for a jump~$u$ of~$T$ at \m{i^{p^{-\nu}}}; as $-u$ lies in
\m{[-\eta,\eta]} or in $\epsilon\mathsf{s}(i^{p^{-\nu}})\cdot[\beta,\beta+V]$ and \m{V+\eta\le W}, the step
lies in~$I$. By Input~\ref{r:points}(c) and~\eqref{e:25} there is a point over \m{p^{25\kappa}c} with
depth~$d$ and exponents $\impE{\raisedE{-\rho(i)}{\cdot}{\cdot}}{d}{\cdot}{\cdot}$. We take this depth~$d$
and \m{r=0}; the letter~$i$ is a cut.

\emph{Case~2: $\mathsf{s}^{\circ}(i)=\epsilon$ and \textup{($\alpha$)} fails.} Then ($\beta$) holds: the chart
is unramified, \m{p\ge5}, $\mathsf{l}(i)=t+J+1$, and $i^{p^{-\mathsf{l}(i)}}$ is a cut at which
\m{\rho=0}. Let \m{j=i^{p^{-t}}}. We apply Theorem~D of Part~2 (Input~\ref{r:walk}) to~$T$ at~$j$ with
\m{l=J+1}, \m{\gamma=G}, \m{\rho=\rho_0}, $N=\rho_0+G\sum_{0\le a\le J}p^{a}$, \m{A=2B} and \m{\delta=X}. Its
hypotheses: $\Kills{T}{j^{p^{-(J+1)}}}{\kappa}{0}$ by Lemma~\ref{l:atcut}; $\Kills{T}{j}{0}{\rho_0}$ by
(W11) and Input~\ref{r:ideal}(h); the jumps of~$T$ at~$j$ are at most \m{\eta\le G}, because
$\mathsf{s}(j)=\epsilon$ puts the far window of~$T$ at~$j$ on the negative side, and those at
\m{j^{p^{-a}}}, \m{0<a\le J}, are at most~$G$ by~($\beta$); \m{N<\beta(j)\le2B} by ($\beta$) and~(W8);
$3e_{\fc}\le p^{J+1}$ by~(W6); $e_{\fc}\cdot2B\le p^{J+1}X$ by~(W5). It gives a point of
$\ptsOver{T}{c}{j}{\kappa+4X}{\shStd{J+1}{o_1}{\cdot}{\cdot}}{\shStd{J+1}{o_2}{\cdot}{\cdot}}$ with
$o_1(u)=\min(2B,-u-N)$ and \m{o_2(u)=2B}, and Input~\ref{r:points}(d) moves it by~\m{\varphi^{t}} to a
point over \m{p^{\kappa+4X}c} at~$i$ with depth $\frobX{t}{1}{\shStd{J+1}{o_2}{\cdot}{\cdot}}{\cdot}{\cdot}$.
Let \m{d(t,u)=2B} and $d(\nu,u)=\std_{t+J+1}(\nu,u)+3\eta_{\lfloor\nu/2\rfloor}$ for \m{\nu\ne t}. The steps
from that depth to~$d$ have \m{\nu\le4J+1}, \m{\nu\ne t}, and \m{0\le n<3\eta} or
\m{-u\le n<-u+3\eta}; as \m{V+4\eta\le W} they lie in~$I$. By Input~\ref{r:points}(c) and~\eqref{e:25} there
is a point over \m{p^{25\kappa}c} at~$i$ with depth~$d$. We take this~$d$ and \m{r=\mathsf{l}(i)}.

\emph{The comparisons.} In both cases it remains to compare exponents. Write~$o^{\ast}$ for the exponents of the
point found. It must be shown that \m{o^{\ast}\ge\mathsf{R}^{\epsilon}_i} on the jumps of~$T$, so that the point
lies in $\ptsOver{T}{c}{i}{25\kappa}{\mathsf{R}^{\epsilon}_i}{d}$ (Input~\ref{r:exp}), and that the depth~$d$
and the integer~$r$ satisfy clauses (a) and~(b) of Definition~2.6 of Part~2 against $\mathsf{R}^{\epsilon'}$
and the family~$z$, for the proposition $\mathsf{s}^{\circ}(i)=\epsilon$. The formal text settles these
inequalities by a decision procedure for linear arithmetic after a case distinction; here they are done by
hand. Put \m{\mathsf{l}=\mathsf{l}(i)}; let \m{\nu\ge0}, let $u$ be a jump of~$T$ and $v$ a jump of~$T'$ at
\m{i^{p^{-\nu}}}, and \m{s=-u-v-1}. By the choice of~$z$, \m{z(i)(0,s)=-\rho(i)-s} and \m{z(i)(\nu,s)=-s} for
\m{\nu\ge1}. Three remarks are used.

\emph{(1)} For every letter~$j$, \m{\mathsf{R}^{\epsilon'}_j(\nu,v)\ge-v-\eta_\nu}, by the definition. Hence the
right side of clause~(a), which is $\min\bigl(-v,\mathsf{R}^{\epsilon'}_i(\nu,v)\bigr)$, is at least
\m{-v-\eta_\nu}; and the right side of the second inequality of clause~(b), which is
$\min\bigl(\mathsf{R}^{\epsilon'}_i(\nu,v),\frobX{\mathsf{l}}{0}{\mathsf{R}^{\epsilon'}_{i'}}{\nu}{v}\bigr)$
with \m{i'=i^{p^{-\mathsf{l}}}}, is at least \m{\min(-v-\eta_\nu,0)} for \m{\nu<\mathsf{l}} and at least
$\min(-v-\eta_\nu,-v-\eta_{\nu-\mathsf{l}})$ for \m{\nu\ge\mathsf{l}}.

\emph{(2)} \m{\eta_\nu\le\eta_{\lfloor\nu/2\rfloor}} for all~$\nu$, and
\m{\eta_{\nu-\mathsf{l}}\le\eta_{\lfloor\nu/2\rfloor}} for \m{\nu\ge\mathsf{l}}: if \m{\lfloor\nu/2\rfloor>2J},
then \m{\nu\ge4J+2} and \m{\nu-\mathsf{l}\ge2J+1}, because \m{\mathsf{l}\le2J+1}. Moreover
\m{\eta_\nu=\eta_{\lfloor\nu/2\rfloor}=\eta} for \m{\nu<\mathsf{l}}.

\emph{(3)} By the windows of $T$ and~$T'$: at a letter~$j'$ with \m{\mathsf{s}(j')=\epsilon} the jumps of~$T$
are at most~$\eta$ and those of~$T'$ at least~$-\eta$; at a letter with \m{\mathsf{s}(j')\ne\epsilon} the jumps
of~$T$ are at least~$-\eta$ and those of~$T'$ at most~$\eta$; at a letter with \m{\beta(j')+V\le\eta} all jumps
lie in \m{[-\eta,\eta]}; and at every letter~$j'$ a jump outside \m{[-\eta,\eta]}, a \emph{far} jump, has
absolute value in \m{[\beta(j'),\beta(j')+V]}.

\emph{The comparisons in Case~1.} Here $o^{\ast}=\impE{\raisedE{-\rho(i)}{\cdot}{\cdot}}{d}{\cdot}{\cdot}$, that
is,
\[\begin{gathered}
o^{\ast}(0,u)=\begin{cases}\rho(i)-u+\eta,&u<-\rho(i),\\ \min(\rho(i)-u+\eta,\,0),&u\ge-\rho(i),\end{cases}\\
o^{\ast}(\nu,u)=\begin{cases}-u+\eta_{\lfloor\nu/2\rfloor},&u<0,\\ \min(-u+\eta_{\lfloor\nu/2\rfloor},\,0),&u\ge0\end{cases}
\qquad(\nu\ge1).
\end{gathered}\]
So \m{o^{\ast}(\nu,u)\ge-u} for \m{\nu\ge1}, and \m{o^{\ast}(\nu,u)\ge-u-\eta_\nu} for all~$\nu$: at \m{\nu=0}
this is clear for \m{u<-\rho(i)} and for \m{u\ge-\eta}, and no jump of~$T$ at the cut~$i$ lies in
\m{[-\rho(i),-\eta)}, because by~(W10) either \m{\rho(i)=0}, or \m{\rho(i)<\beta(i)} while a jump \m{u<-\eta}
is far and has \m{u\le-\beta(i)}. If \m{\mathsf{s}^{\circ}(i)\ne\epsilon}, then
\m{\mathsf{R}^{\epsilon}_i(\nu,u)=-u-\eta_\nu} and \m{o^{\ast}\ge\mathsf{R}^{\epsilon}_i} is proved. If
\m{\mathsf{s}^{\circ}(i)=\epsilon}, then ($\alpha$) holds, and with~(3) it gives
\begin{equation}\label{e:alpha}
u\le\eta\quad\text{and}\quad v\ge-\eta\qquad\text{at every }\nu<\mathsf{l}.
\end{equation}
At these~$\nu$ the formulas give \m{o^{\ast}(\nu,u)\ge0=\std_{\mathsf{l}}(\nu,u)}, by~(2); and at
\m{\nu\ge\mathsf{l}}, \m{o^{\ast}(\nu,u)\ge-u=\std_{\mathsf{l}}(\nu,u)}. So
$o^{\ast}\ge\max(\std_{\mathsf{l}},-u-\eta_\nu)=\mathsf{R}^{\epsilon}_i$ on the jumps.

Clause~(a), with \m{d(0,u)=\rho(i)-u+\eta} and \m{d(\nu,u)=-u+\eta_{\lfloor\nu/2\rfloor}} for \m{\nu\ge1}: the
left side \m{-z(i)(\nu,s)-d(\nu,u)} is \m{-v-1-\eta} at \m{\nu=0} and \m{-v-1-\eta_{\lfloor\nu/2\rfloor}} at
\m{\nu\ge1}. By~(2) it is smaller than \m{-v-\eta_\nu}, and (1) gives~(a). Clause~(b): let
\m{\mathsf{s}^{\circ}(i)=\epsilon}. The first inequality, \m{\std_{\mathsf{l}}\le\mathsf{R}^{\epsilon}_i}, holds
by the definition of~$\mathsf{R}^{\epsilon}_i$. In the second, as \m{r=0}, the left side is that of~(a). At
\m{\nu<\mathsf{l}} it is \m{-v-1-\eta}, by~(2), which is smaller than \m{-v-\eta} and, by~\eqref{e:alpha},
negative. At \m{\nu\ge\mathsf{l}} it is \m{-v-1-\eta_{\lfloor\nu/2\rfloor}}, smaller than \m{-v-\eta_\nu} and
than \m{-v-\eta_{\nu-\mathsf{l}}} by~(2). By~(1), clause~(b) holds.

\emph{The comparisons in Case~2.} Here \m{\mathsf{s}(j)=\epsilon} for \m{j=i^{p^{-t}}},
\m{r=\mathsf{l}=t+J+1} and \m{\rho(i^{p^{-\mathsf{l}}})=0}; moreover \m{\rho(i)=0} unless \m{t=0}, and
in that case \m{j=i}. By (W1), (W8) and~($\beta$), $\eta\le W\le G\le N<\beta(j)$; so by~(3) every jump of $T$ or~$T'$
at~$j$ has absolute value at most \m{\beta(j)+V}, and by (W1), (W7) and~(W8), applied to the letter~$j$,
\begin{equation}\label{e:2B}
\beta(j)+V+\eta<\beta(j)+W\le2B-\rho(j).
\end{equation}
By ($\beta$) and~(3),
\begin{equation}\label{e:beta}
u\ge-\eta\ \text{ and }\ v\le\eta\quad\text{at }\nu<\mathsf{l},\ \nu\ne t;\qquad
u\le\eta\ \text{ and }\ v\ge-\eta\quad\text{at }\nu=t,
\end{equation}
and a jump~$u$ of~$T$ at~$j$ with \m{u\ge-N} is not far, since a far one has \m{u\le-\beta(j)<-N}; so it has
\m{|u|\le\eta}. The point moved by~\m{\varphi^{t}} has the exponents
$\frobX{t}{0}{\shStd{J+1}{o_1}{\cdot}{\cdot}}{\cdot}{\cdot}$, which are~$0$ at \m{\nu<\mathsf{l}}, \m{\nu\ne t},
are \m{\min(2B,-u-N)} at \m{\nu=t}, and are~$-u$ at \m{\nu\ge\mathsf{l}}. As \m{d(\nu,u)=3\eta\ge0} at
\m{\nu<\mathsf{l}}, \m{\nu\ne t}, and \m{d(\nu,u)=-u+3\eta_{\lfloor\nu/2\rfloor}} at \m{\nu\ge\mathsf{l}},
Input~\ref{r:points}(c) gives
\[\begin{gathered}o^{\ast}(\nu,u)=0\ \ (\nu<\mathsf{l},\ \nu\ne t),\qquad
o^{\ast}(\nu,u)\ge-u\ \ (\nu\ge\mathsf{l}),\\
o^{\ast}(t,u)=\begin{cases}2B,&u<-N,\\0,&u\ge-N.\end{cases}\end{gathered}\]
This is at least $\mathsf{R}^{\epsilon}_i(\nu,u)=\max(\std_{\mathsf{l}}(\nu,u),-u-\eta_\nu)$. At
\m{\nu<\mathsf{l}}, \m{\nu\ne t}: \m{-u-\eta\le0} by~\eqref{e:beta}. At \m{\nu=t}: if \m{u<-N}, then
\m{-u-\eta\le-u\le\beta(j)+V\le2B} by~\eqref{e:2B}; if \m{u\ge-N}, then \m{|u|\le\eta} and \m{-u-\eta\le0}. At
\m{\nu\ge\mathsf{l}}: \m{\std_{\mathsf{l}}(\nu,u)=-u}.

Clause~(a). At \m{\nu<\mathsf{l}}, \m{\nu\ne t}, the left side is \m{s-3\eta}; for \m{\nu=0} this uses
\m{\rho(i)=0}, which holds because then \m{t\ne0}. It is smaller than \m{-v-\eta}, since \m{u\ge-\eta}
by~\eqref{e:beta}. At \m{\nu=t} the left side is \m{s-2B} if \m{t\ge1} and \m{\rho(i)+s-2B} if \m{t=0}; it is
smaller than \m{-v-\eta} because \m{\rho(i)-u+\eta\le2B}: indeed \m{-u\le\beta(j)+V}, and \m{\rho(i)} is~$0$ or
is~\m{\rho(j)}, so \eqref{e:2B} applies. At \m{\nu\ge\mathsf{l}} the left side is
\m{-v-1-3\eta_{\lfloor\nu/2\rfloor}}, smaller than \m{-v-\eta_\nu} by~(2). By~(1), clause~(a) holds.

Clause~(b). The first inequality holds by the definition of~$\mathsf{R}^{\epsilon}_i$, as
\m{\mathsf{s}^{\circ}(i)=\epsilon}. For the second, \m{z(i^{p^{-\mathsf{l}}})=\raisedE{0}{\cdot}{\cdot}}, so
$\frobX{\mathsf{l}}{0}{z(i^{p^{-\mathsf{l}}})}{\nu}{s}$ is~$0$ at \m{\nu<\mathsf{l}} and~$-s$ at
\m{\nu\ge\mathsf{l}}. At \m{\nu<\mathsf{l}} the left side is therefore \m{-d(\nu,u)}, and by~(1) it suffices
that \m{-d(\nu,u)<\min(-v-\eta,0)}. At \m{\nu=t} this reads \m{-2B<\min(-v-\eta,0)}, and
\m{v+\eta\le\beta(j)+V+\eta<2B} by~\eqref{e:2B}. At \m{\nu\ne t} it reads \m{-3\eta<\min(-v-\eta,0)}, which
holds because \m{v\le\eta} by~\eqref{e:beta} and \m{\eta\ge1} by~(W7). At \m{\nu\ge\mathsf{l}} the left side is
\m{s-d(\nu,u)=-v-1-3\eta_{\lfloor\nu/2\rfloor}}, smaller than \m{-v-\eta_\nu} and than
\m{-v-\eta_{\nu-\mathsf{l}}} by~(2). By~(1), clause~(b) holds, and the claim is proved in both cases.

All hypotheses of Theorem~B of Part~2 hold: every letter has a cut among the \m{i^{p^{-a}}}, \m{a>0}, by~(iii);
its hypotheses (i) and~(ii) are the claim; (iii) is the point of the third side; and (iv) was noted. It gives
$\CupB{\Pi}{E}{c}{c'}$.
\end{proof}

\begin{proof}[Proof of Theorem C]
By (W1) and~(W7), \m{W\ge4}; by the hypotheses \m{G\le\eta\le W}, $B\ge11(p^{J}+1)W\ge22W$ and
\m{X\ge16}. The objects $\mathsf{cut}$, $\rho(\cdot)$, $n$, $\kappa$, $\rho_0$ of Setting~\ref{set:cut} do not occur in
the theorem; the proof chooses them, and it replaces~$G$. So Proposition~\ref{p:cuts} is applied to another
instance of Setting~\ref{set:cut}, with the same $S$, $S'$, $\Pi$, $\mathsf{s}$, $\beta$, $\eta$, $V$, $W$,
$X$, $B$ and~$J$; condition~(W12) need not hold in that instance, and it is not among the hypotheses of the
Proposition. Call a letter~$i$ a \kw{cut} if no \m{i^{p^{a}}} with \m{1\le a\le J} is thick. We apply
Proposition~\ref{p:cuts} with this property, with \m{\rho(i)=0} for all~$i$, with \m{2W} in the place of~$n$
and of~$G$, with $\rho_0=\lfloor3B/5\rfloor$, \m{\kappa=12(J+1)X+1} and \m{E=10^{4}X(J+1)}, and with
$B$, $J$, $X$, $V$, $W$, $\eta$, $\beta$, $\mathsf{s}$ as they are. Conditions (W1), (W2), (W5), (W6), (W7) are
assumed, (W4) holds by the choice of~$\kappa$, and (W3) reads
$(J+1+\fd+\lfloor\log_p(2B)\rfloor)\cdot2W+\fd\cdot2W\le X$, which follows from the hypothesis on~$X$. (W8):
$\beta(i)+W\le B+W\le2B$. (W9): \m{2W\le(p-1)\cdot2W}. (W10): at a cut~$i$ the letters \m{i^{p^{a}}},
\m{1\le a\le J}, are not thick, so $\beta(i^{p^{a}})+W=G+W\le2W$. (W11):
$B+W\le2\lfloor3B/5\rfloor$, because \m{B\ge5W+8}. Hypothesis~(i): \m{W\le2pW}. (ii):
$51\kappa=612(J+1)X+51\le10^{4}X(J+1)$.

\emph{A thick letter that is not a cut.} Let $j$ be thick and not a cut. Then \m{j^{p^{-(J+1)}}} is a cut.
Indeed some \m{j^{p^{a'}}}, \m{1\le a'\le J}, is thick; if \m{j^{p^{-(J+1)+b}}} were thick for some
\m{1\le b\le J}, the hypothesis of the theorem for this letter and the exponents \m{a=J+1-b} and~$a'$ would
say that \m{j^{p^{a'}}} is not thick.

\emph{Hypothesis~(iii).} Suppose that no \m{i^{p^{-l}}}, \m{1\le l\le2J+1}, is a cut. As
\m{i^{p^{-(J+1)}}} is not a cut, some \m{i^{p^{-a}}} with \m{1\le a\le J} is thick; it is not a cut, so
\m{i^{p^{-(a+J+1)}}} is a cut, and \m{a+J+1\le2J+1}.

\emph{Hypothesis~(iv).} Let $i$ and~$l$ be as there. If \m{i^{p^{-a}}} is thick with \m{0<a<l}, then it is
not a cut, so \m{i^{p^{-(a+J+1)}}} is a cut and \m{l\le a+J+1}; and \m{l<a+J+1} is impossible, because the cut
\m{i^{p^{-l}}} would have the thick letter \m{i^{p^{-a}}} among its translates by \m{p^{b}},
\m{1\le b\le J}. So among the \m{i^{p^{-a}}}, \m{0\le a<l}, only those with \m{a=0} and \m{a=l-J-1} can be
thick. If ($\alpha$) holds with \m{t=0} we are done. Otherwise some \m{a<l} has
$\mathsf{s}(i^{p^{-a}})\ne\mathsf{s}(i)$ and $\beta(i^{p^{-a}})+V>\eta$; then \m{i^{p^{-a}}} is thick,
because \m{G+V\le\eta}, and \m{a\ne0}, so \m{a=l-J-1}. Take \m{t=a}. The signs are not all equal, so the
chart is unramified and \m{p\ge7}; \m{\rho=0} everywhere; \m{t+J+1=l}; and
\[\rho_0+2W\sum_{0\le b\le J}p^{b}<\tfrac35B+4Wp^{J}\le\tfrac35B+\tfrac4{11}B<B=\beta(i^{p^{-t}}).\]
For \m{b<l}, \m{b\ne t}: if \m{i^{p^{-b}}} is thick then \m{b=0} and
$\mathsf{s}(i)\ne\mathsf{s}(i^{p^{-t}})$; if it is not thick then $\beta(i^{p^{-b}})+V=G+V\le\eta$; and
for \m{t<b} it is not thick, so $\beta(i^{p^{-b}})+W=G+W\le2W$. This is~($\beta$).

Proposition~\ref{p:cuts} gives $\CupB{\Pi}{10^{4}X(J+1)}{c}{c'}$.
\end{proof}

\begin{proof}[Proof of Theorem D]
By (W1) and~(W7), \m{W\ge4}, so \m{G\ge20}, \m{B\ge G>0} and \m{K_0\ge1}. Put
$\lambda=\lfloor\log_p(2B)\rfloor$. For a letter~$i$ let $\rho(i)=\max(\lfloor B/p\rfloor-7K_0,0)$ if
\m{i^{p}} is thick and \m{\rho(i)=0} otherwise, and call~$i$ a \kw{cut} if \m{\rho(i)<\beta(i)}. Since
$\lfloor B/p\rfloor<B$ and \m{\beta(i)\in\{B,G\}}, a letter~$i$ is not a cut if and only if
\begin{equation}\label{e:noncut}
i\text{ is not thick},\qquad i^{p}\text{ is thick},\qquad G+7K_0\le\lfloor B/p\rfloor .
\end{equation}
As in the proof of Theorem~C, Proposition~\ref{p:cuts} is applied to another instance of
Setting~\ref{set:cut}, with the same $S$, $S'$, $\Pi$, $\mathsf{s}$, $\beta$, $\eta$, $V$, $W$ and~$B$, in
which (W12) is not needed: with this property and this~$\rho$, with $K_0$ in the place of~$n$,
\m{J=1}, $\rho_0=\lfloor B/p\rfloor+12K_0$, \m{G+W} in the place of~$G$, \m{X=(2+\lambda)K_0},
\m{\kappa=24X+1} and $E=10^{4}K_0(\lambda+1)$. The chart is unramified, so \m{e_{\fc}=1}, \m{\fd=0}. (W1),
(W2), (W7) are assumed; (W3) and (W4) hold by the choices. (W5): $2B<2K_0p^{2}\le p^{2}X$. (W6):
\m{8\le p^{2}}. (W8): \m{W\le K_0}, and $\rho(i)+W\le\max(\lfloor B/p\rfloor-6K_0,\,W)\le B$, so
$\beta(i)+\rho(i)+W\le2B$. (W9):
$G+W\le G+K_0<K_0(p/2+1)\le(p-1)K_0$. (W10): let $i$ be a cut; if \m{i^{p}} is not thick, then
\m{\rho(i)=0} and \m{\beta(i^{p})+W=G+W}; if \m{i^{p}} is thick, then \m{\rho(i)<\beta(i)}, the chart is
unramified, \m{p\ge5}, and $\rho_0\le\kappa+\rho(i)$ because $\rho(i)\ge\lfloor B/p\rfloor-7K_0$ and
\m{19K_0\le\kappa}. (W11): with $B\le p\lfloor B/p\rfloor+p-1$ and $\lfloor B/p\rfloor\le K_0p-1$,
\[(p-3)\rho_0-(B+K_0)\ \ge\ (12p-37)K_0-3\lfloor B/p\rfloor-p+1\ \ge\ (9p-37)K_0-p+4\ >0 .\]
Hypothesis~(i): $\rho(i)+W=W\le pK_0$ if \m{\rho(i)=0}, and
$\rho(i)+W\le\lfloor B/p\rfloor-6K_0<pK_0$ otherwise. (ii):
$51\kappa=1224(2+\lambda)K_0+51\le10^{4}K_0(\lambda+1)$. (iii): if \m{i^{p^{-1}}} is not a cut, it is not
thick by~\eqref{e:noncut}, so \m{i^{p^{-2}}} is a cut, again by~\eqref{e:noncut}; thus \m{l\le2}.

(iv): let $i$ and~$l$ be as there; \m{l\le2}. Take \m{t=0}. For \m{l\le1}, ($\alpha$) holds. Let \m{l=2}:
\m{i^{p^{-1}}} is not a cut, so by~\eqref{e:noncut} it is not thick, $i$ is thick and
$G+7K_0\le\lfloor B/p\rfloor$. If ($\alpha$) fails, then $\mathsf{s}(i^{p^{-1}})\ne\mathsf{s}(i)$, and
($\beta$) holds: the chart is unramified and \m{p\ge7}; \m{\rho(i^{p^{-2}})=0} because \m{i^{p^{-1}}} is not thick; \m{t+J+1=2};
$\beta(i^{p^{-1}})+W=G+W$; and, as $B-\lfloor B/p\rfloor\ge(p-1)(G+7K_0)$, \m{W\le K_0} and
\m{2G<K_0p},
\begin{align*}B-\rho_0-(G+W)(1+p)\ &\ge\ (p-1)(G+7K_0)-12K_0-(p+1)(G+K_0)\\
&=(6p-20)K_0-2G\ >0 .\end{align*}
Proposition~\ref{p:cuts} gives $\CupB{\Pi}{10^{4}K_0(\lambda+1)}{c}{c'}$.
\end{proof}

\sect{The two theorems on lattices}{s:lattices}
\begin{lemma}[orthogonality from the universal sides]\label{l:orth}
In Setting~\ref{set:base}, let $\Gamma$ carry a topology for which $j_\Gamma$ is an open embedding. Let $\fc$
be a Kummer chart for~$\Gamma$, \m{\omega\in\cO} with \m{\omega^{p^h-1}=1}, $\mathcal{F}$ a filter on~$\Gamma$
every member of which contains~$\Gamma'$, $U$ and~$U'$ families of finite sets of integers indexed by the maps
$\gamma\colon\cO\to C$, $\mathcal{Q}$ a class of lattices over~$\cO$ contained in
$\InCl{\mathcal{F}}{\bot}{U}$, $L$ a lattice with $L\in\InCl{\mathcal{F}}{\bot}{U'}$, \m{e\ge0} an integer
and $x\in Z^{1}_{\mathrm{f}}(L)$. Suppose that $\CupB{\Pi}{e}{c}{x}$ for every \m{Y\in\mathcal{Q}}, every
pairing~$\Pi$ of the sides $\univS{\fc}{\omega}{\mathcal{F}}{Y;U}$ and $\univS{\fc}{\omega}{\mathcal{F}}{L;U'}$, and every
$c\in Z^{1}_{\mathrm{f}}(Y)$. Then $p^{e}x\in\cupOrth{\mathcal{Q}}{L}$.
\end{lemma}
\leanline{\lean{memcupOrthogonalofscaledCupBounds}}
\begin{proof}
A pairing of the two universal sides is a pairing of the lattices $Y$ and~$L$, their cocycle maps are the
identity, and the topology of~$\Gamma$ is one of those of Definition~2.4 of Part~2. So the hypothesis says
that the cup product of $c$ and~\m{p^{e}x} bounds, for every \m{Y\in\mathcal{Q}}, every pairing and every
$c\in Z^{1}_{\mathrm{f}}(Y)$; this is the definition of $\cupOrth{\mathcal{Q}}{L}$.
\end{proof}

\begin{proof}[Proof of Theorem A]
The filter~$\top$ has the single member~$\Gamma$, which contains~$\Gamma'$. For a lattice~$Y$ in
$\InCl{\top}{\bot}{U}$ write $S_Y=\univS{\fc}{\omega}{\top}{Y;U}$, and $S_L=\univS{\fc}{\omega}{\top}{L;U'}$. Their
jump sets at~$i$ are $\letterJ{\fc}{\omega}{U}{i}$ and $\letterJ{\fc}{\omega}{U'}{i}$, in the windows of
the statement.

(1) Let $x\in Z^{1}_{\mathrm{f}}(L)$, $Y$ in the class, $\Pi$ a pairing of $S_Y$ and~$S_L$, and
$c\in Z^{1}_{\mathrm{f}}(Y)$. We apply Theorem~C of Part~2 (Input~\ref{r:walk}) with \m{\kappa=0}, the
given~$\rho$, and \m{H=0}. Its hypothesis~(i) holds by Input~\ref{r:ideal}(h): the jumps of $S_Y$, $S_L$ and
of the third side are at most $\beta(i)+W+\eta\le(p-3)\rho$. Hypotheses (ii) and~(iii) are assumed. (iv):
in the first case of~(1) no element of the two jump sets lies in \m{[-\rho,0)}, as $0$ does not; in the
second, the chart is unramified, \m{p\ge5}, \m{\eta\ge1}, and
$\letterJ{\fc}{\omega}{U'}{i}\subseteq[-\eta,0]$. (v): \m{\fd=0}. So $\CupB{\Pi}{0}{c}{x}$, and
Lemma~\ref{l:orth} with \m{e=0} gives $x\in\cupOrth{\InCl{\top}{\bot}{U}}{L}$.

(2) For a letter~$i$ let $g_i$ be the number of the maps~$\gamma$ with $\gamma(\omega)=\theta([i])$. The side
$S_L$ satisfies Setting~\ref{set:few} with the given $\rho,\eta,W,\beta,\mathsf{s}$ and
\[n_{\mathrm{lo}}(i)=h\,m_0\,g_i,\qquad n_{\mathrm{all}}(i)=h\,\operatorname{rk}_{\cO}(L)\,g_i :\]
(F1) to (F5) are hypotheses of the theorem, and (F6) holds because \m{m_L(i,u)} is the sum over these~$\gamma$
of the multiplicities of~$u$, which for each~$\gamma$ have sum at most~$m_0$ over \m{|u|\le\eta} and sum
$\operatorname{rk}_{\cO}(L)$ over all~$u$ (Input~\ref{k:bk}). Let $\mathsf{C}\subseteq Z^{1}_{\mathrm{f}}(L)$ be the
$\Zp$-submodule of Proposition~\ref{p:few}. For \m{d\in \mathsf{C}}, $Y$ in the class, a pairing~$\Pi$ of $S_Y$
and~$S_L$ and $c\in Z^{1}_{\mathrm{f}}(Y)$, the proposition gives $\CupB{\Pi}{0}{c}{d}$; by
Lemma~\ref{l:orth}, $\mathsf{C}\subseteq\cupOrth{\InCl{\top}{\bot}{U}}{L}$. Distinct letters have distinct
\m{\theta([i])}, and there are $\operatorname{rk}_{\Zp}(\cO)$ maps~$\gamma$ in all, so
$\sum_ig_i\le\operatorname{rk}_{\Zp}(\cO)$ and
\begin{align*}\relLen{\Zp}{Z^{1}_{\mathrm{f}}(L)}{\cupOrth{\InCl{\top}{\bot}{U}}{L}}&\le\relLen{\Zp}{Z^{1}_{\mathrm{f}}(L)}{\mathsf{C}}\\
&\le11\,h\,\bigl(m_0(\rho+1)+\operatorname{rk}_{\cO}(L)(W+1)\bigr)\cdot\operatorname{rk}_{\Zp}(\cO).\end{align*}
Let \m{f_{\cO}} be the degree of the residue field of~$\cO$ over~$\mathbb{F}_p$. The length over~$\Zp$ of an
$\cO$-module of finite length is $f_{\cO}$ times its length over~$\cO$, and
$\operatorname{rk}_{\Zp}(\cO)=\operatorname{len}_{\Zp}(\cO/p\cO)=f_{\cO}\operatorname{len}_{\cO}(\cO/p\cO)$.
Dividing by~$f_{\cO}$ gives~(2).
\end{proof}

\begin{proof}[Proof of Theorem B]
\emph{The tame chart.} Let \m{M_0} be the least common multiple of the \m{p^{a}-1}, $1\le a\le\mathsf{r}$.
It divides \m{p^{\mathsf{r}!}-1}, and $p$ does not divide~$\mathsf{r}!$ because \m{\mathsf{r}<p}. By
Proposition~\ref{p:tameexist} there is a tame chart~$\fc$ for~$\Gamma$ with index \m{e_{\fc}=M_0}; let
$\sigma$ be its tame generator and $\varpi$ its element. Then $e_{\fc}\le\prod_{a\le\mathsf{r}}p^{a}\le
p^{\mathsf{r}^{2}}$, and \m{\fd\le1}: the index is prime to~$p$, so $\fd$ is $0$ or~$1$. As a coefficient
root we take \m{\omega=1}.

\emph{The lattices are crystalline on~$\Gamma'$.} Let $\mathcal{F}'$ be the filter of the subsets of~$\Gamma$
that contain~$\Gamma'$. We claim: if $\mathsf{K}_1$ is a relation with
$\TameType{p}{\mathsf{r}}{\mathsf{n}}{\mathsf{K}_1(s,\cdot)}$ for all inertial~$s$, then every lattice~$Y$ in
$\InCl{\mathcal{F}}{\mathsf{K}_1}{U}$ lies in $\InCl{\mathcal{F}'}{\bot}{U}$. Indeed $Y$ is potentially
crystalline. As $j_\Gamma$ is bijective, $g\mapsto\linF{Y}{j_\Gamma^{-1}(g)}{0}$ is a map from the inertia
group of~$\GQp$ to the endomorphisms of \m{B_{\max}\otimes Y}; it is a homomorphism of monoids and is trivial
near~$1$ (Input~\ref{k:frob}(a)), and its $M_0$-th power is trivial, because \m{X^{M_0}-1} kills \m{(s,0)}
on~$Y$ for inertial~$s$ by the definitions of the tame type and of the class. By Input~\ref{k:tame}, for
\m{y=\varpi} and the exponent \m{M_0=e_{\fc}}, it is trivial on the stabilizer of~$\varpi$, which contains
\m{j_\Gamma(s)} for the inertial \m{s\in\Gamma'}. By Input~\ref{k:crit}, applied to the group~$\Gamma'$ with the restriction of~$j_\Gamma$, for which
the linear Frobenius of an inertial element is the same, $Y$ is crystalline on~$\Gamma'$. The condition on the jumps is unchanged, and the relation~$\bot$ imposes nothing.
The claim applies to $\mathsf{K}$ and to~$\mathsf{K}^{-}$, since the inverse of an inertial element is
inertial.

\emph{The killers.} By the hypothesis on~$\mathsf{K}$ at $\sigma$ and at~$\sigma^{-1}$ there are monic
polynomials $q$, $q'$ over~$\Zp$ of degree at most~$\mathsf{n}$ with $\mathsf{K}(\sigma,q)$ and
$\mathsf{K}(\sigma^{-1},q')$. Let $Y$ be a lattice in $\InCl{\mathcal{F}}{\mathsf{K}}{U}$. Then $q$ kills
\m{(\sigma,0)} on~$Y$ and $q'$ kills \m{(\sigma^{-1},0)} on~$Y$; and since $L\in\InCl{\mathcal{F}}{\mathsf{K}^{-}}{U'}$,
$q'$ kills \m{(\sigma,0)} on~$L$ and $q$ kills \m{(\sigma^{-1},0)} on~$L$. By Input~\ref{k:dual} for
\m{s=\sigma^{-1}}, a monic polynomial~$q_3$ over~$\Zp$ of degree at most~$\mathsf{n}^{2}$ kills
\m{(\sigma,0)} on $(Y\otimes_{\cO}L)^{\vee}(1)$.

\emph{The walk.} Let $S_Y=\univS{\fc}{1}{\mathcal{F}'}{Y;U}$ and $S_L=\univS{\fc}{1}{\mathcal{F}'}{L;U'}$, let $\Pi$ be a pairing of
$S_Y$ and~$S_L$, $c\in Z^{1}_{\mathrm{f}}(Y)$ and $a\in Z^{1}_{\mathrm{f}}(L)$. The jump sets of the two
sides are in the windows of the statement, with \m{\beta(i)=\beta} and $\mathsf{s}(i)=\mathsf{s}$ for all~$i$,
and the jumps of the two sides and of the third side are at most \m{n=\beta+6(W+\eta)}. Put
$\kappa=\mathsf{r}^{2}(\mathsf{n}+1)^{2}$. By Proposition~\ref{p:tamekill} with \m{m=\mathsf{r}^{2}}, for the
polynomials $q$, $q'$ and~$q_3$,
\[\Kills{S_Y}{i}{\kappa}{0},\qquad\Kills{S_L}{i}{\kappa}{0},\qquad\Kills{\third{\Pi}}{i}{\kappa}{0}\]
for every letter~$i$: \m{n+\kappa+2<p} is the hypothesis, and
$\mathsf{r}^{2}(\mathsf{n}^{2}+1)\le\kappa$. We apply Theorem~C of Part~2 (Input~\ref{r:walk}) with this
$\kappa$, with \m{\rho=0} and $H=3\kappa+150(\eta+W+1)$. (i) has just been shown. (ii): a non-zero element
of $\twoWin{6(\eta+W)}{\beta}\cup\twoWin{6(\eta+W)}{-\beta}$ has absolute value at most
\m{\beta+6(\eta+W)<p}. (iii) is assumed. (iv): the interval \m{[-0,0)} is empty. (v):
$3(\kappa+(48(\eta+W)+4)\fd)\le3\kappa+144(\eta+W)+12\le H$, as \m{\fd\le1}. So
$\CupB{\Pi}{H}{c}{a}$, and Lemma~\ref{l:orth}, for the filter~$\mathcal{F}'$, the class
$\mathcal{Q}=\InCl{\mathcal{F}}{\mathsf{K}}{U}$ and \m{e=H}, gives
$p^{H}a\in\cupOrth{\InCl{\mathcal{F}}{\mathsf{K}}{U}}{L}$.
\end{proof}

\sect{Where the results are used}{s:used}
The unit on Selmer groups of the handout bounds the loss at one place above~$\ell$ by the four theorems, in
its note \emph{Columns and rows as Galois representations} (the lemma called \emph{the crystalline case} there,
and the proof of Theorem~E of that note). Theorem~B serves the tame case of that proof, in which $\ell$ is one of
the primes of the second level of the data. The other three serve the crystalline case, through the lemma, whose
integers $Q$, $h$ and~$\lambda$ are those of that note (its~$h$ is not the residue degree of a chart):
Theorem~C (sparse thick letters) when \m{\ell\le Q}, through either of its alternatives: the Kummer chart
unramified and \m{\ell\ge7}, where the signs may vary, or constant signs; Theorem~D (two sizes) when
\m{Q<\ell} and \m{h^{3}\le\lambda}; and Theorem~A (few low weights), with both of its assertions, when
\m{Q<\ell} and \m{\lambda<h^{3}}. Nothing else of the three parts is used there. Inside the unit,
Proposition~\ref{p:few} is used only for Theorem~A, the results on tame charts only for Theorem~B, and
Proposition~\ref{p:cuts} only for Theorems C and~D.

\emph{A dictionary.} The notes on the unit on Selmer groups write $\mathcal{Q}(F,K,U)$ for the class
$\InCl{\mathcal{F}}{\mathsf{K}}{U}$, $Z_{\mathrm{f}}(L)$ for $Z^{1}_{\mathrm{f}}(L)$, ``of tame type
\m{(\ell,\varrho,M)}'' for $\TameType{\ell}{\varrho}{M}{\cdot}$, $f_k$ and~$e_k$ for the residue degree~$h$ and
the index~$e_{\fc}$ of a chart, $\epsilon$ for the signs~$\mathsf{s}$, and \m{B_{\mathrm{cris}}\otimes L} for
\m{B_{\max}\otimes L}, with the remark that the ring is Colmez's~$B_{\max}$. ``The note on crystalline lifts''
there means the three parts on this unit; the theorems it quotes are Theorems A to~D of this part.

\sect{Notes on the formal text}{s:notes}
\emph{Arguments as the text gives them.} The arguments were read from the Lean text and the census of the
workspace, and were not printed by Lean. The hypotheses of Settings \ref{set:few} and~\ref{set:cut} are
section hypotheses of the formal text; each declaration includes those that its statement here names, by the
commands of the text that include and omit them, which were read. Several declarations carry implicit
arguments that their statements do not use: a letter and an exponent function in
Proposition~\ref{p:few}, a function of the letters in Theorem~C, an integer in Lemma~\ref{l:atcut} and
Theorem~C, and the function~$\beta$ in Lemma~\ref{l:rel}. The universal side has four explicit arguments: the chart, the
root, the hypothesis on the filter and the membership of the lattice in a class. Its notation shows the filter
in the slot of the hypothesis on it, and the lattice with its jump sets in the slot of the membership. The statement of
Lemma~\ref{l:atcut} is the formal one with its auxiliary integer~$n'$.

\emph{The number~$\fd$ for a tame chart.} For a tame chart with index different from~$1$ the number~$\fd$ is
\m{v_p(e_{\fc})+1=1}, not~$0$. An earlier state of the reader's dossier put it equal to~$0$ in the proof of
Theorem~B; the second reader of the workspace corrected this, the dossier now says so, and the formal proof
agrees: it unfolds the definition and uses only that $p$ does not divide the index. The constant \m{150(\eta+W+1)} absorbs the term.

\emph{The product in Lemma~\ref{l:extkill}.} The polynomial is the product \m{q\cdot(X-1)}, evaluated at the
linear Frobenius; the factor \m{X-1} is what removes the constant~$w$ in the proof of
Proposition~\ref{p:tamekill}.

\emph{The filter.} The hypothesis of the formal text on the filter~$\mathcal{F}$ says that every member
of~$\mathcal{F}$ contains~$\Gamma'$, not that $\Gamma'$ is a member.

\emph{Lengths.} In Theorem~A the length is over~$\cO$, and the factor $\operatorname{len}_{\cO}(\cO/p\cO)$
is the ramification index of~$\cO$; the formal proof passes through lengths over~$\Zp$ as in the proof above.

\emph{Arithmetic.} The formal text settles the comparisons of sizes and exponents by decision procedures.
The note gives them by hand in Lemma~\ref{l:col}, Proposition~\ref{p:few}, Lemma~\ref{l:atcut},
Proposition~\ref{p:cuts} and the proofs of Theorems A, B, C and~D. In Proposition~\ref{p:cuts} the final
comparisons use (W1), (W2), (W7), (W8), (W10), the choice of~\m{t(i)} in~(iv) and, in Case~2, the inequality
\m{N<\beta(j)}; of the family \m{\mathsf{R}^{\epsilon'}} of the other side they use only the lower bound
\m{-v-\eta_\nu}, as the formal proof does. The hypothesis \m{\eta\ge1} of Proposition~\ref{p:few} is not
used by the argument given for it. Theorems C and~D do not involve the objects
$\mathsf{cut}$, $\rho(\cdot)$, $n$, $\kappa$, $\rho_0$ of Setting~\ref{set:cut} (Theorem~D not $J$ and~$X$
either): the formal statements omit the corresponding hypotheses, and their proofs choose these objects. In the proof of Lemma~\ref{l:col} the integer~$\mathsf{N}$ is the least
multiple of~$p$ that is at least \m{\beta(i)+W+2\rho}.

\emph{The background.} The period ring is Colmez's~$B_{\max}$ (Setting~\ref{set:base}). The Inputs K1
to~K8 are statements of the background library; the authors report that Lean's kernel accepted their proofs,
and they are not proved in this note. Two of them must not be read as citations of published theorems.
The crystalline criterion (Input~\ref{k:crit}) is one implication of Fontaine's criterion, in a generality
(arbitrary modules, any group of open image, the ring~$B_{\max}$) for which the second reader finds no
published statement; its place in the literature is not confirmed. The finiteness of the module of crystalline
cocycles (Input~\ref{k:bk}) is a statement about integral cocycles defined with~$B_{\max}$; the Bloch--Kato
formula is about the rational cohomology defined with~$B_{\mathrm{cris}}$, and the bridge between the two
(coboundaries, the comparison of the two period rings, continuity) is sketched by the second reader and not
located in the literature. For the linear Frobenius the second reader finds the classical statement for
representations over a field in Fontaine's exposés and no source for the form on arbitrary modules. No place
in the literature was confirmed for the frame theorem, for the tame inertia statement in the form used, and
for the killing of duals of tensor products.

\subhead{How this note was written.} From the reader's dossier of the unit: the statements of all members and
the structured proofs of the declarations treated here; and from the second reader's two reports on the unit,
whose corrections were applied. The three Lean files were read in full. The proofs of Lemma~\ref{l:rel} and
of Proposition~\ref{p:tameexist}, which the reader had marked as least certain, were read line by line in the
Lean text: in the first, the identity $p(x^{\ast}-w_4)=\mathsf{B}(px^{\ast}-w)+\cE_{\nu_0}w_3$ and the three
places where hypothesis~(i) enters were recomputed by hand; in the second, the normality of~$\Delta$ and the
bound for the exponent of the quotient. The verifications of the deep points in Proposition~\ref{p:few}, the
induction in Proposition~\ref{p:tamekill} and the reductions of Theorems C and~D to
Proposition~\ref{p:cuts} were worked out by hand; the formal proofs leave them to automation. The comparisons
at the end of the proof of Proposition~\ref{p:cuts} were first left to the formal text; the referee of this
note carried them out and found that they also use~(W10), and they were then checked once more and written
into the proof. The program check of the unit's dossier
covered \UnitSteps{} leaf steps of \UnitProofs{} proofs.

\head{Appendix A. Concordance with the Lean artifact}
{\fontsize{8.6}{11.6}\selectfont\renewcommand{\leanfont}{\ttfamily\fontsize{7.7}{10}\selectfont\color{leancol}}
\setlength{\LTleft}{0pt}\setlength{\LTright}{\fill}\setlength{\tabcolsep}{0pt}
\begin{longtable}{@{}p{30mm}p{56mm}p{43mm}r@{}}
\textbf{In this note} & \textbf{Lean declaration} & \textbf{File} & \textbf{Lines}\\\midrule\endhead
\ConcordanceRows
\end{longtable}}

Two statements as the artifact has them:
\quoted{fewlowjumps}
\quoted{tametypesexponent}

\emph{Legend.} As in Parts 1 and~2; in addition \texttt{InClass F K U} is the class
$\InCl{\mathcal{F}}{\mathsf{K}}{U}$, \texttt{cupOrthogonal} is $\cupOrth{\mathcal{Q}}{L}$,
\texttt{crysCocycles} is \m{Z^{1}_{\mathrm{f}}}, \texttt{relLength} is the relative length,
\texttt{linearFrob} is the linear Frobenius, \texttt{IsWeilPair (toGp s) 0} says that $s$ is inertial,
\texttt{jump} is the multiplicity of a jump, \texttt{k.f} is~$h$ and \texttt{k.e} is~$e_{\fc}$,
\texttt{differentExp} is~$\fd$, \texttt{letterFrob} is \m{i\mapsto i^{p}} and \texttt{stepBack i a} is
\m{i^{p^{-a}}}; \texttt{zetaE}, \texttt{tameGen} and \texttt{ϖ} of a tame chart are $\varrho$, $\sigma$
and~$\varpi$. In the first file \texttt{a} is~$\rho$, \texttt{b} is~$\beta$, \texttt{ε} is~$\mathsf{s}$,
\texttt{n} and \texttt{r} are $n_{\mathrm{lo}}$ and~$n_{\mathrm{all}}$, and in its first lemma \texttt{a}
is~$\nu_0$ and \texttt{c}, \texttt{c'} are $o$, $o'$. In the second file \texttt{s}, \texttt{s₀} are $\rho$,
$\rho_0$, \texttt{b} is~$B$, \texttt{T} is the property of being thick, \texttt{K} is~$K_0$, \texttt{e} is~$E$,
and \texttt{ν} in the proof of Proposition~\ref{p:tamekill} is~$a$. In the third file \texttt{ζ}
is~$\omega$, \texttt{r} and \texttt{N} are $\mathsf{r}$ and~$\mathsf{n}$, \texttt{m} is~$m_0$ and \texttt{g}
is~$\gamma$.

\head{Notation}
\begingroup\fontsize{8.6}{10.9}\selectfont\raggedright
\setlength{\columnsep}{9mm}\setlength{\parskip}{0pt}
\begin{multicols}{2}
\nt{$\ell$}{the prime of Chapter~VIII of the handout; it is~$p$ here}
\nt{$p$, $\Qp$, $\Zp$, $\bar{\mathbb{Q}}_p$, $C$, $\mathcal{O}_C$, $\GQp$, $\mathbb{F}_p$}{the prime, the $p$-adic numbers and integers, an algebraic closure, its completion, the integers, the Galois group, the prime field}
\nt{$\Ainf$, $\WC$, $\mathcal{O}_C^\flat$}{Witt vectors of the tilt and of its fraction field, the tilt}
\nt{$A_{\max}$, $B^{+}$, $B_{\max}$, $\xi$}{Colmez's period rings (Setting~\ref{set:base}; Part~1, Setting~2.1): the $p$-adic completion of \m{\Ainf[X]/(pX-\xi)}, \m{A_{\max}[1/p]}, \m{B^{+}[1/([\varepsilon]-1)]}; a generator of \m{\ker\theta}}
\nt{$\varphi$, $\theta$, $\varepsilon$, $[\,\cdot\,]$}{Frobenius, Fontaine's map, a system of roots of unity in the tilt, the Teichmüller lift}
\nt{$\BdR$, $\Fil^{n}\BdR$}{the field of de Rham periods and its filtration}
\nt{$\Gamma$, $\Gamma'$, $j_\Gamma$, $\chi$, $\chi_N$, $\Delta$}{the group, the subgroup of the chart, the map to~$\GQp$, the cyclotomic character, the character on the $N$-th roots of unity, an open subgroup}
\nt{$\fc$, $\fS$, $\fS_B$, $\WW$, $h$, $X$, $z$, $\cE$, $\cE_\nu$, $e_{\fc}$, $\fd$, $F_0$}{the chart, its ring, the ring of Part~1, $W(\mathbb{F}_{p^h})$, the residue degree, the variable, the coordinate, the Eisenstein element and its translates, the index, the different exponent, a value field}
\nt{$\varrho$, $\sigma$, $\varpi$, $\varpi^{\flat}$}{the root of unity, the tame generator and the root of~$p$ of a tame chart; an element of the tilt above~$\varpi$}
\nt{$\cO$, $\omega$, $\gamma$, $f_{\cO}$}{coefficients, the root of unity, a map \m{\cO\to C}, the residue degree of~$\cO$}
\nt{$i$, $i'$, $j$, $i^-$}{letters; \m{i^-=i^{p^{-1}}}}
\nt{$L$, $Y$, $T$, $T'$}{lattices and lattice data; $T$, $T'$ are also sides}
\nt{$Z^{1}_{\mathrm{f}}(L)$, $\operatorname{rk}_{\cO}(L)$, $U_L(i)$, $m_L(i,u)$, $h_L$}{crystalline cocycles, the rank, the jump set and the multiplicity of a jump at a letter, the jump bound}
\nt{$U$, $U'$, $U(\gamma)$}{families of finite sets of integers}
\nt{$\pr_i$, $\Per(L)$, $\Minf(L)$, $\Mlocinf(L)$, $\Mloc(L)$, $\mathfrak{M}(L)$}{projector, periods, Kisin modules}
\nt{\m{\wExp{L}{i}{o}}, \m{\locLP{l}{i}(L)}, \m{\locKP{Q}}, \m{\invMod{Q}}}{exponent modules; elements of \m{\Mloc}; invariants modulo~$Q$}
\nt{$\std_l$, \m{\cost{\nu}{n}}, $\impE{o}{o'}{\nu}{u}$, $\frobX{l}{n}{o}{\nu}{u}$, $\raiseX{o}{\nu}{u}{\nu'}{v}$, $\shStd{l}{o_1}{\nu}{u}$}{standard exponents, cost, improved, shifted and raised exponents}
\nt{$S$, $S'$, $Y_S$, $C_S$, $\alpha_S$, $S_Y$, $S_L$}{sides and their parts; universal sides}
\nt{\m{\extL{S}}, \m{\trivL{S}}, \m{\projC{S}}, \m{\inclY{S}}, \m{\projCloc{S}}}{extension, trivial datum, projection, inclusion}
\nt{\m{\constAt{S}{i}{c}}, $\liftId{S}{r}{l}{i}$, $\Kills{S}{i}{\kappa}{\rho}$}{constants, lifting ideal, cokernel killer}
\nt{$\raisedE{n}{\nu}{u}$, \m{\twoWin{n}{c}}, $\NoMult{p}{n}{\beta}$, $\JumpWin{\eta}{\mathsf{s}}{\beta}{W}{U}$}{raised exponents and windows}
\nt{$\Pi$, \m{\third{\Pi}}, $\CupB{\Pi}{e}{c}{c'}$}{a pairing, its third side, the scaled cup product bounds}
\nt{\m{\kerX{S}{i}{o}}, \m{\deepE{S}{i}{o}}, $\Reaches{S}{e}{\mathsf{A}}$, $\ptsOver{S}{c}{i}{e}{o_1}{o'}$}{kernel with exponents, deep elements, cost of a map, points}
\nt{\m{\linF{T}{s}{m}}}{the linear Frobenius}
\nt{$\InCl{\mathcal{F}}{\mathsf{K}}{U}$, $\mathcal{F}$, $\mathcal{F}'$, $\mathsf{K}$, $\mathsf{K}^{-}$, $\mathsf{K}_0$, $\mathsf{K}_1$, $\top$, $\bot$, $\mathcal{Q}$}{a class of lattices; filters; relations and a property of polynomials; the trivial filter, the empty relation; a class}
\nt{$\cupOrth{\mathcal{Q}}{L}$, \m{\relLen{R}{N}{N'}}, $\operatorname{len}$, $R$}{the orthogonal module, the relative length, length, a ring}
\nt{\m{\deepC{S}{i}{n}{o}}, $\TameType{p}{\mathsf{r}}{\mathsf{n}}{\mathsf{K}_0}$, $\letterJ{\fc}{\omega}{U}{i}$, $\univS{\fc}{\omega}{\mathcal{F}}{L;U}$}{constants that lift, tame types, jump sets at a letter, the universal side}
\nt{$\mathsf{cut}$, $\mathsf{thick}$, $\mathsf{Q}$, $\mathsf{l}(i)$, $t(i)$, $\mathsf{s}^{\circ}(i)$}{properties of letters, the distance to the nearest cut, the integer of~(iv), a sign}
\nt{$o$, $o'$, $o^{+}$, $o_+$, $o_1$, $o_2$, $o_i$, $o^{T}$, $o^{S}_i$, $\mathsf{z}$, $\mathsf{z}'$, $z(i)$, $z'(i)$, $d$, $\mathsf{R}^{\epsilon}_i$}{exponent functions and families of them; $o_1$, $o_2$ are also functions of~$u$ alone; $z$ is also the coordinate, $d$ also an element}
\nt{$\nu$, $\nu'$, $\nu_0$, $l$, $u$, $v$, $s$}{indices, levels, jumps; $s$ is also an element of~$\Gamma$}
\nt{$\rho$, $\eta$, $W$, $V$, $\beta$, $B$, $G$, $J$, $J'$, $X$, $\kappa$, $\rho_0$, $\rho_u$, $W_\nu$, $\eta_\nu$, $n$, $n'$, $K_0$, $E$, $H$, $\lambda$}{the integers of the Settings and of the proofs; $X$ is also the variable, $\lambda$ also a linear map}
\nt{$n_{\mathrm{lo}}$, $n_{\mathrm{all}}$, $m_0$, $g_i$, $\mathsf{r}$, $\mathsf{n}$, $M_0$, $h_0$, $N$, $\mathsf{N}$, $N_a$}{bounds for multiplicities, the number of maps above a letter, the sizes of a tame type, a least common multiple, a residue degree, \m{p^{h_0}-1}, a multiple of~$p$, a number of roots}
\nt{$\mathsf{s}$, $\epsilon$, $\epsilon'$}{signs; $\epsilon$ is also a cost}
\nt{$o^{\ast}$}{the exponents of the point found, in the proof of Proposition~\ref{p:cuts}}
\nt{$\mathcal{Q}(F,K,U)$, $Z_{\mathrm{f}}(L)$, $f_k$, $e_k$, \m{(\ell,\varrho,M)}, $B_{\mathrm{cris}}$}{the notation of the notes on Selmer groups (Section~\ref{s:used}); $B_{\mathrm{cris}}$ is also the ring defined with divided powers, which is not used}
\nt{$A$, $B_u$, $a$, $a'$, $a_0$, $b$, $c$, $c'$, $e$, $e'$, $k$, $m$, $r$, $t$, $x_0$, $\gamma$, $\delta$}{integers; $B$ of Lemma~\ref{l:rel} is an integer; $a$ is also a root of unity and a cocycle, $c$ an element, $k$, $r$, $t$ also elements, $\delta$ an element of~$\Delta$}
\nt{$g$, $g'$, $\sigma^{k}$}{elements of~$\Gamma$ or of~$\GQp$}
\nt{$x$, $x'_k$, $x_k$, $x^{\ast}$, $y$, $y'$, $y''$, $y_k$, $w$, $w_k$, $w_3$, $w_4$, $d$, $t_0$, $f$}{elements of modules and rings}
\nt{$\mathsf{P}_i(o)$, $\mathsf{P}^{0}_i(o)$, $N$, $N'$, $N_i$, $I$, $I_2$}{modules and sets of steps of Section~\ref{s:few}; $N$ is also an integer}
\nt{$\mathsf{C}$}{the submodule of~$C_S$ of Proposition~\ref{p:few}}
\nt{$\mathsf{A}$, $\mathsf{B}$}{additive maps}
\nt{$q$, $q'$, $q_3$, $q_k$, $Q$, $Q_k$, $T$}{polynomials and their coefficients; an element of~$\Ainf$}
\nt{$\mu$, $\mu_k$, $\Lambda_g$}{periods and linear maps}
\nt{$v_p(n)$, $\lfloor\log_pn\rfloor$, $\deg q$, $\operatorname{rk}_{\Zp}(\cO)$, $\operatorname{Stab}(\varpi)$}{the exponent of~$p$ in~$n$, the integral logarithm, the degree, the rank over~$\Zp$, the stabilizer of~$\varpi$ in~$\GQp$}
\end{multicols}
\endgroup

\end{document}
